% This LaTeX document needs to be compiled with XeLaTeX.
\documentclass[10pt]{article}
\usepackage[utf8]{inputenc}
\usepackage{amsmath}
\usepackage{amsfonts}
\usepackage{amssymb}
\usepackage[version=4]{mhchem}
\usepackage{stmaryrd}
\usepackage{hyperref}
\hypersetup{colorlinks=true, linkcolor=blue, filecolor=magenta, urlcolor=cyan,}
\urlstyle{same}
\usepackage{graphicx}
\usepackage[export]{adjustbox}
\graphicspath{ {./images/} }
\usepackage{caption}
\usepackage{bbold}
\usepackage{multirow}
\usepackage{fvextra, csquotes}
% \usepackage[fallback]{xeCJK}
\usepackage{polyglossia}
\usepackage{fontspec}
\usepackage{newunicodechar}
\IfFontExistsTF{Noto Serif CJK TC}
{\setCJKmainfont{Noto Serif CJK TC}}
{\IfFontExistsTF{STSong}
  {\setCJKmainfont{STSong}}
  {\IfFontExistsTF{Droid Sans Fallback}
    {\setCJKmainfont{Droid Sans Fallback}}
    {\setCJKmainfont{SimSun}}
}}

\setmainlanguage{english}
\IfFontExistsTF{CMU Serif}
{\setmainfont{CMU Serif}}
{\IfFontExistsTF{DejaVu Sans}
  {\setmainfont{DejaVu Sans}}
  {\setmainfont{Georgia}}
}

 Finite-Sample Properties of OLS }
%New command to display footnote whose markers will always be hidden
\let\svthefootnote\thefootnote
\newcommand\blfootnotetext[1]{%
  \let\thefootnote\relax\footnote{#1}%
  \addtocounter{footnote}{-1}%
  \let\thefootnote\svthefootnote%
}
%Overriding the \footnotetext command to hide the marker if its value is `0`
\let\svfootnotetext\footnotetext
\renewcommand\footnotetext[2][?]{%
  \if\relax#1\relax%
    \ifnum\value{footnote}=0\blfootnotetext{#2}\else\svfootnotetext{#2}\fi%
  \else%
    \if?#1\ifnum\value{footnote}=0\blfootnotetext{#2}\else\svfootnotetext{#2}\fi%
    \else\svfootnotetext[#1]{#2}\fi%
  \fi
}
\newunicodechar{×}{\ifmmode\times\else{$\times$}\fi}
\newunicodechar{⋮}{\ifmmode\vdots\else{$\vdots$}\fi}
\newunicodechar{⇔}{\ifmmode\Leftrightarrow\else{$\Leftrightarrow$}\fi}
\newunicodechar{¥}{\ifmmode\text{\yen}\else\yen\fi}
\title{Chapter 7: Extremum Estimators}

\author{}
\date{}

\begin{document}
\maketitle

\section*{CHAPTER 7}
\section*{Extremum Estimators}
\begin{abstract}
In the previous chapters, the generalized method of moments (GMM) has been our choice for estimating the various models we have considered. The method of maximum likelihood (ML) has been discussed, but only in passing. In this chapter, we expand our repertoire of estimation techniques by examining a class of estimators called "extremum estimators," which includes (linear and nonlinear) least squares, (linear and nonlinear) GMM, and ML as special cases. As has been emphasized by Amemiya (1985) and others, this unified approach is useful for bringing out the common structure that underlies those apparently diverse estimation principles.\\
Section 7.1 introduces extremum estimators. The two sections that follow develop the asymptotic properties of extremum estimators. Section 7.4 defines the trio of test statistics - the Wald, Langrange multipler, and likelihood-ratio statis-tics-for extremum estimators in general, and shows how they can be specialized to individual cases. Section 7.5 is a brief discussion of the computational aspects of extremum estimators.\\
The discussion of this chapter may appear daunting, as it combines a heavy use of asymptotic theory with calculus. The essence, however, is really straightforward. If you can read graduate textbooks on micro and macroeconomics without too much difficulty, you should be able to understand most of the discussion on the second (if not the first) reading.\\
If you are not interested in the asymptotics of extremum estimators per se but plan to study Chapter 8, there is no need to read this chapter in its entirety. Just read Section 7.1 and then try to understand the statements in Propositions 7.5, 7.6, 7.8, 7.9, and 7.11.
\end{abstract}

A Note on Notation: Unlike in the previous chapters, we use the parameter vector with subscript 0 for its true value. So $\boldsymbol{\theta}_{0}$ is the true parameter value, while $\boldsymbol{\theta}$ represents a hypothetical parameter value. This notation is standard in the "highbrow" literature on nonlinear estimation. Also, we will use " $t$ " rather than " $i$ " for the observation index.

\subsection*{7.1 Extremum Estimators}
An estimator $\hat{\boldsymbol{\theta}}$ is called an extremum estimator if there is a scalar objective function $Q_{n}(\boldsymbol{\theta})$ such that


\begin{equation*}
\hat{\boldsymbol{\theta}} \text { maximizes } Q_{n}(\boldsymbol{\theta}) \text { subject to } \boldsymbol{\theta} \in \Theta \subset \mathbb{R}^{p}, \tag{7.1.1}
\end{equation*}


where $\Theta$, called the parameter space, is the set of possible parameter values. In this book, we restrict our attention to the case where $\Theta$ is a subset of the finite dimensional Euclidean space $\mathbb{R}^{p}$. The objective function $Q_{n}(\boldsymbol{\theta})$ depends not only on $\boldsymbol{\theta}$ but also on the sample or the data, $\left(\mathbf{w}_{1}, \mathbf{w}_{2}, \ldots, \mathbf{w}_{n}\right)$, where $\mathbf{w}_{t}$ is the $t$-th observation and $n$ is the sample size. In our notation the dependence of the objective function on the sample of size $n$ is signalled by the subscript $n$. The maximum likelihood (ML) and the linear and nonlinear generalized method of moments (GMM) estimators are particular extremum estimators. This section presents these and other extremum estimators and indicates how they are related to each other.

\section*{"Measurability" of $\hat{\boldsymbol{\theta}}$}
To be sure, the maximization problem (7.1.1) may not necessarily have a solution. But recall from calculus the following fact:

Let $h: \mathbb{R}^{p} \rightarrow \mathbb{R}$ be continuous and $A \subset \mathbb{R}^{p}$ be a compact (closed and bounded) set. Then $h$ has a maximum on the set $A$. That is, there exists an $\mathbf{x}^{*} \in A$ such that $h\left(\mathbf{x}^{*}\right) \geq h(\mathbf{x})$ for all $\mathbf{x}$ in $A$.

Therefore, if $Q_{n}(\boldsymbol{\theta})$ is continuous in $\boldsymbol{\theta}$ for any data ( $\mathbf{w}_{1}, \mathbf{w}_{2}, \ldots, \mathbf{w}_{n}$ ) and $\Theta$ is compact, then there exists a $\boldsymbol{\theta}$ that solves the maximization problem in (7.1.1) for any given data $\left(\mathbf{w}_{1}, \ldots, \mathbf{w}_{n}\right)$. In the event of multiple solutions, we would choose one from them. So $\hat{\boldsymbol{\theta}}$, thus uniquely determined for any data ( $\mathbf{w}_{1}, \ldots, \mathbf{w}_{n}$ ), is a function of the data. Strictly speaking, however, being a function of the vector of random variables ( $\mathbf{w}_{1}, \ldots, \mathbf{w}_{n}$ ) is not enough to make $\hat{\boldsymbol{\theta}}$ a well-defined random variable; $\hat{\boldsymbol{\theta}}$ needs to be a "measurable" function of $\left(\mathbf{w}_{1}, \ldots, \mathbf{w}_{n}\right)$. (If a function is continuous, it is measurable.) The following lemma shows that $\hat{\boldsymbol{\theta}}$ is measurable if $Q_{n}(\boldsymbol{\theta})$ is

Lemma 7.1 (Existence of extremum estimators): Suppose that (i) the parameter space $\Theta$ is a compact subset of $\mathbb{R}^{p}$, (ii) $Q_{n}(\boldsymbol{\theta})$ is continuous in $\boldsymbol{\theta}$ for any data $\left(\mathbf{w}_{1}, \ldots, \mathbf{w}_{n}\right)$, and (iii) $Q_{n}(\boldsymbol{\theta})$ is a measurable function of the data for all $\boldsymbol{\theta}$ in $\Theta$. Then there exists a measurable function $\hat{\boldsymbol{\theta}}$ of the data that solves (7.1.1).

See, e.g., Gourieroux and Monfort (1995, Property 24.1) for a proof. In all the examples of this chapter, $Q_{n}(\boldsymbol{\theta})$ is measurable.

In most applications, we do not know the upper or lower bound for the true parameter vector. Even if we do, those bounds are not included in the parameter space, so the parameter space is not closed (see the CES example below for an example). So the compactness assumption for $\Theta$ is something we wish to avoid. In some of the asymptotic results of this chapter, we will replace the compactness assumption by some other conditions that are satisfied in many applications.

\section*{Two Classes of Extremum Estimators}
In the next two sections, we will develop asymptotic results for extremum estimators and then specialize those results to the following two classes of extremum estimators.

\begin{enumerate}
  \item M-Estimators. An extremum estimator is an M-estimator if the objective function is a sample average:
\end{enumerate}


\begin{equation*}
Q_{n}(\boldsymbol{\theta})=\frac{1}{n} \sum_{t=1}^{n} m\left(\mathbf{w}_{t} ; \boldsymbol{\theta}\right), \tag{7.1.2}
\end{equation*}


where $m$ is a real-valued function of ( $\mathbf{w}_{t}, \boldsymbol{\theta}$ ). Two examples of an M-estimator we study are the maximum likelihood (ML) and the nonlinear least squares (NLS).\\
2. GMM. An extremum estimator is a GMM estimator if the objective function can be written as


\begin{equation*}
Q_{n}(\boldsymbol{\theta})=-\frac{1}{2} \mathbf{g}_{n}(\boldsymbol{\theta})^{\prime} \widehat{\mathbf{W}} \mathbf{g}_{n}(\boldsymbol{\theta}) \text { with } \underset{(K \times 1)}{\mathbf{g}_{n}(\boldsymbol{\theta})} \equiv \frac{1}{n} \sum_{t=1}^{n} \mathbf{g}\left(\mathbf{w}_{t} ; \boldsymbol{\theta}\right), \tag{7.1.3}
\end{equation*}


where $\widehat{\mathbf{W}}$ is a $K \times K$ symmetric and positive definite matrix that defines the distance of $\mathbf{g}_{n}(\boldsymbol{\theta})$ from zero. It can depend on the data. Maximizing this objective function is equivalent to minimizing the distance $\mathbf{g}_{n}(\boldsymbol{\theta})^{\prime} \widehat{\mathbf{W}} \mathbf{g}_{n}(\boldsymbol{\theta})$, which is how we defined GMM in Chapter 3. As will become clear, the deflation of the distance by 2 simplifies the expression for the derivative of the objective function. We have introduced GMM estimation for linear models where $\mathbf{g}\left(\mathbf{w}_{t} ; \boldsymbol{\theta}\right)$ is linear in $\boldsymbol{\theta}$. Below we will show that GMM estimation can be extended easily to nonlinear models.

There are extremum estimators that do not fall into either class. A prominent example is classical minimum distance estimators. Their objective function can be\\
written as $-\mathbf{g}_{n}(\boldsymbol{\theta})^{\prime} \widehat{\mathbf{W}} \mathbf{g}_{n}(\boldsymbol{\theta})$ but $\mathbf{g}_{n}(\boldsymbol{\theta})$ is not necessarily a sample mean. The consistency theorem of the next section is general enough to cover this class as well.

In the rest of this section, we present examples of M-estimators and GMM estimators. Each example is a combination of the model (a set of data generating processes [DGPs]) and the estimation method.

\section*{Maximum Likelihood (ML)}
A prime example of an M-estimator is the maximum likelihood for the case where $\left\{\mathbf{w}_{t}\right\}$ is i.i.d. The model here is a set of i.i.d. sequences $\left\{\mathbf{w}_{t}\right\}$ where the density of $\mathbf{w}_{t}$ is a member of the family of densities indexed by a finite-dimensional vector $\boldsymbol{\theta}: f\left(\mathbf{w}_{t} ; \boldsymbol{\theta}\right), \boldsymbol{\theta} \in \Theta$. (Since $\left\{\mathbf{w}_{t}\right\}$ is identically distributed, the functional form of $f(\cdot ; \cdot)$ does not depend on $t$.) The functional form of $f(\cdot ; \cdot)$ is known. The model is parametric because the parameter vector $\boldsymbol{\theta}$ is finite dimensional. At the true parameter value $\theta_{0}$, the density of the true DGP (the DGP that actually generated the data) is $f\left(\mathbf{w}_{t} ; \boldsymbol{\theta}_{0}\right) .{ }^{1}$ We say that the model is correctly specified if $\boldsymbol{\theta}_{0} \in \Theta$.

Since $\left\{\mathbf{w}_{t}\right\}$ is independently distributed, the joint density of the data ( $\mathbf{w}_{1}$, $\mathbf{w}_{2}, \ldots, \mathbf{w}_{n}$ ) is


\begin{equation*}
f\left(\mathbf{w}_{1}, \mathbf{w}_{2}, \ldots, \mathbf{w}_{n} ; \boldsymbol{\theta}_{0}\right)=\prod_{t=1}^{n} f\left(\mathbf{w}_{t} ; \boldsymbol{\theta}_{0}\right) . \tag{7.1.4}
\end{equation*}


With the distribution of the data thus completely specified, the natural estimation method is maximum likelihood, which proceeds as follows. With the true parameter vector $\boldsymbol{\theta}_{0}$ replaced by its hypothetical value $\boldsymbol{\theta}$, this density, viewed as a function of $\theta$, is called the likelihood function. The maximum likelihood (ML) estimator of $\boldsymbol{\theta}_{0}$ is the $\boldsymbol{\theta}$ that maximizes the likelihood function. Since the log transformation is a monotone transformation, maximizing the likelihood function is equivalent to maximizing the log likelihood function: ${ }^{2}$


\begin{equation*}
\log f\left(\mathbf{w}_{1}, \mathbf{w}_{2}, \ldots, \mathbf{w}_{n} ; \boldsymbol{\theta}\right)=\log \left[\prod_{t=1}^{n} f\left(\mathbf{w}_{t} ; \boldsymbol{\theta}\right)\right]=\sum_{t=1}^{n} \log f\left(\mathbf{w}_{t} ; \boldsymbol{\theta}\right) . \tag{7.1.5}
\end{equation*}


The ML estimator of $\boldsymbol{\theta}_{0}$, therefore, is an M-estimator with

\footnotetext{${ }^{1}$ Had we adhered to the notational convention of the previous sections, the true parameter value would be denoted by $\theta$ and a hypothetical value by $\bar{\theta}$. In this chapter, we use $\theta_{0}$ for the true value and $\theta$ for a hypothetical value of the parameter vector.\\
${ }^{2}$ For $\theta$ such that $f\left(\mathbf{w}_{t} ; \theta\right)=0$, the $\log$ cannot be taken. This does not present a problem because such $\theta$ never maximizes the (level) likelihood function $f\left(\mathbf{w}_{1}, \ldots, \mathbf{w}_{n} ; \boldsymbol{\theta}\right)$. For such $\theta$ we can assign an arbitrarily small number to $\log f\left(\mathbf{w}_{t} ; \theta\right)$ so that the maximizer of the level likelihood is also the maximization of the log likelihood. See White (1994, p. 15) for a more formal treatment of this issue.
}

\begin{equation*}
m\left(\mathbf{w}_{i} ; \boldsymbol{\theta}\right)=\log f\left(\mathbf{w}_{i} ; \boldsymbol{\theta}\right), \text { that is, } Q_{n}(\boldsymbol{\theta})=\frac{1}{n} \sum_{t=1}^{n} \log f\left(\mathbf{w}_{t} ; \boldsymbol{\theta}\right) . \tag{7.1.6}
\end{equation*}


Before illustrating ML by an example, it is useful to make three remarks.

\begin{itemize}
  \item (Serially correlated observations) The average log likelihood would not take a form as simple as this if $\left\{\mathbf{w}_{t}\right\}$ were serially correlated. Examples of ML for serially correlated observations will be studied in the next chapter.
  \item (Alternative estimators) Maximum likelihood is not the only way to estimate the parameter vector. For example, let $\boldsymbol{\mu}(\boldsymbol{\theta})$ be the expectation of $\mathbf{w}_{t}$ implied by the density function $f\left(\mathbf{w}_{t} ; \boldsymbol{\theta}\right)$. It is a known function of $\boldsymbol{\theta}$. By construction, the following zero-mean condition holds:
\end{itemize}


\begin{equation*}
\mathrm{E}\left[\mathbf{w}_{t}-\boldsymbol{\mu}\left(\boldsymbol{\theta}_{0}\right)\right]=\mathbf{0} . \tag{7.1.7}
\end{equation*}


The parameter vector $\boldsymbol{\theta}_{0}$ could be estimated by GMM with $\mathbf{g}\left(\mathbf{w}_{t} ; \boldsymbol{\theta}\right)=\mathbf{w}_{t}-\boldsymbol{\mu}(\boldsymbol{\theta})$ in the GMM objective function (7.1.3).

\begin{itemize}
  \item (Efficiency of ML) As will be shown in the next two sections, ML is consistent and asymptotically normal under suitable conditions. It was widely believed that ML is efficient (i.e., achieves minimum asymptotic variance) in the class of all consistent and asymptotically normal estimators. This belief was shown to be wrong by a counterexample (described, for example, in Amemiya, 1985, Example 4.2.4). Nevertheless, ML is efficient in quite general classes of asymptotically normal estimators. One such general class is GMM estimators. In Section 7.3, we will note that the asymptotic variance of any GMM estimator of $\boldsymbol{\theta}_{0}$ is no smaller than that of the ML estimator.
\end{itemize}

Example 7.1 (Estimating the mean of a normal distribution): Let the data $\left(w_{1}, \ldots, w_{n}\right)$ be a scalar i.i.d. sequence with the distribution of $w_{t}$ given by $N\left(\mu, \sigma^{2}\right)$. So $\boldsymbol{\theta}=\left(\mu, \sigma^{2}\right)^{\prime}$ and

$$
f\left(w_{t} ; \mu, \sigma^{2}\right)=\frac{1}{\sqrt{2 \pi \sigma^{2}}} \exp \left[-\frac{\left(w_{t}-\mu\right)^{2}}{2 \sigma^{2}}\right] .
$$

The average log likelihood of the data ( $w_{1}, \ldots, w_{n}$ ) is


\begin{equation*}
\frac{1}{n} \sum_{t=1}^{n} \log f\left(w_{t} ; \mu, \sigma^{2}\right)=-\frac{1}{2} \log (2 \pi)-\frac{1}{2} \log \left(\sigma^{2}\right)-\frac{1}{n} \sum_{t=1}^{n}\left[\frac{\left(w_{t}-\mu\right)^{2}}{2 \sigma^{2}}\right] . \tag{7.1.8}
\end{equation*}


The ML estimator of $\left(\mu_{0}, \sigma_{0}^{2}\right)$ is an extremum estimator with $Q_{n}(\boldsymbol{\theta})$ taken to be this average log likelihood. Suppose that the variance is assumed to be positive but that there is no a priori restriction on $\mu_{0}$, so the parameter space $\Theta$ is $\mathbb{R} \times \mathbb{R}_{++}$, where $\mathbb{R}_{++}$is the set of positive real numbers. It is easy to check that the ML estimator of $\mu_{0}$ is the sample mean of $w_{t}$. The GMM estimator of $\mu_{0}$ based on the zero-mean condition $\mathrm{E}\left(w_{t}-\mu_{0}\right)=0$, too, is the sample mean of $w_{t}$.

\section*{Conditional Maximum Likelihood}
In most applications, the vector $\mathbf{w}_{t}$ is partitioned into two groups, $y_{t}$ and $\mathbf{x}_{t}$, and the researcher's interest is to examine how $\mathbf{x}_{t}$ influences the conditional distribution of $y_{t}$ given $\mathbf{x}_{t}$. We have been calling the variable $y_{t}$ the dependent variable and $\mathbf{x}_{t}$ regressors. None of the results of this chapter depends on $y_{t}$ being a scalar; $y_{t}$ can be a vector and still all the results will remain valid with the scalar " $y_{t}$ " replaced by a vector " $\mathbf{y}_{t}$." We nevertheless use the scalar notation, simply because in all the examples we consider in this chapter, there is only one dependent variable.

Let $f\left(y_{t} \mid \mathbf{x}_{t} ; \boldsymbol{\theta}_{0}\right)$ be the conditional density of $y_{t}$ given $\mathbf{x}_{t}$, and let $f\left(\mathbf{x}_{t} ; \boldsymbol{\psi}_{0}\right)$ be the marginal density of $\mathbf{x}_{t}$. Then


\begin{equation*}
f\left(y_{t}, \mathbf{x}_{t} ; \boldsymbol{\theta}_{0}, \boldsymbol{\psi}_{0}\right)=f\left(y_{t} \mid \mathbf{x}_{t} ; \boldsymbol{\theta}_{0}\right) f\left(\mathbf{x}_{t} ; \boldsymbol{\psi}_{0}\right) \tag{7.1.9}
\end{equation*}


is the joint density of $\mathbf{w}_{t}=\left(y_{t}, \mathbf{x}_{t}^{\prime}\right)^{\prime}$. Suppose, for now, that $\boldsymbol{\theta}_{0}$ and $\boldsymbol{\psi}_{0}$ are not functionally related (see the next paragraph for an example of a functional relationship). The average log likelihood of the data $\left(\mathbf{w}_{1}, \ldots, \mathbf{w}_{n}\right)$ is


\begin{equation*}
\frac{1}{n} \sum_{t=1}^{n} \log f\left(\mathbf{w}_{t} ; \boldsymbol{\theta}, \boldsymbol{\psi}\right)=\frac{1}{n} \sum_{t=1}^{n} \log f\left(y_{t} \mid \mathbf{x}_{t} ; \boldsymbol{\theta}\right)+\frac{1}{n} \sum_{t=1}^{n} \log f\left(\mathbf{x}_{t} ; \boldsymbol{\psi}\right) . \tag{7.1.10}
\end{equation*}


The first term on the right-hand side is the average log conditional likelihood. The conditional ML estimator of $\boldsymbol{\theta}_{0}$ maximizes this first term, thus ignoring the second term. It is an M-estimator with


\begin{equation*}
m\left(\mathbf{w}_{t} ; \boldsymbol{\theta}\right)=\log f\left(y_{t} \mid \mathbf{x}_{t} ; \boldsymbol{\theta}\right), \text { that is, } Q_{n}(\boldsymbol{\theta})=\frac{1}{n} \sum_{t=1}^{n} \log f\left(y_{t} \mid \mathbf{x}_{t} ; \boldsymbol{\theta}\right) \tag{7.1.11}
\end{equation*}


The second term on the right-hand side of (7.1.10) is the average log marginal likelihood. It does not depend on $\boldsymbol{\theta}$, so the conditional ML estimator of $\boldsymbol{\theta}_{0}$ is\\
numerically the same as the (joint or full) ML estimator that maximizes the average joint likelihood (7.1.10).

Now suppose that $\boldsymbol{\theta}_{0}$ and $\boldsymbol{\psi}_{0}$ are functionally related. For example, $\boldsymbol{\theta}_{0}$ and $\boldsymbol{\psi}_{0}$ might be partitioned as $\boldsymbol{\theta}_{0}=\left(\boldsymbol{\alpha}_{0}^{\prime}, \boldsymbol{\beta}_{0}^{\prime}\right)^{\prime}, \boldsymbol{\psi}_{0}=\left(\boldsymbol{\beta}_{0}^{\prime}, \boldsymbol{\gamma}_{0}^{\prime}\right)^{\prime}$. Then the full (i.e., unconditional) ML and conditional ML estimators are no longer numerically equal. It is intuitively clear that the conditional ML misses information that could have been obtained from the marginal likelihood. Indeed it can be shown that the conditional ML estimator of $\boldsymbol{\theta}_{0}$ is less efficient than the full ML estimator obtained from maximizing the joint log likelihood (7.1.10) (see, e.g., Gourieroux and Monfort, 1995, Section 7.5.3). In most applications, this loss of efficiency is unavoidable because we do not know, or do not wish to specify, the parametric form, $f\left(\mathbf{x}_{t} ; \boldsymbol{\psi}\right)$, of the marginal distribution.

Here are two examples of conditional ML (other examples will be in the next chapter).

Example 7.2 (Linear regression model with normal errors): In Chapter 2, we considered the linear regression model with conditional homoskedasticity. Assume further that the error term is normally distributed, so


\begin{equation*}
\left\{y_{t}, \mathbf{x}_{t}\right\} \text { is i.i.d., } y_{t}=\mathbf{x}_{t}^{\prime} \boldsymbol{\beta}_{0}+\varepsilon_{t}, \quad \varepsilon_{t} \mid \mathbf{x}_{t} \sim N\left(0, \sigma_{0}^{2}\right) . \tag{7.1.12}
\end{equation*}


The conditional likelihood of $y_{t}$ given $\mathbf{x}_{t}, f\left(\mathbf{y}_{t} \mid \mathbf{x}_{t} ; \boldsymbol{\theta}\right)$, is the density function of $N\left(\mathbf{x}_{t}^{\prime} \boldsymbol{\beta}, \sigma^{2}\right)$. So, with $\boldsymbol{\theta}=\left(\boldsymbol{\beta}^{\prime}, \sigma^{2}\right)^{\prime}$ and $\mathbf{w}_{t}=\left(y_{t}, \mathbf{x}_{t}^{\prime}\right)^{\prime}$, the $m$ function is


\begin{align*}
m\left(\mathbf{w}_{t} ; \boldsymbol{\theta}\right) & =\log f\left(y_{t} \mid \mathbf{x}_{t} ; \boldsymbol{\beta}, \sigma^{2}\right) \\
& =\log \left(\frac{1}{\sqrt{2 \pi \sigma^{2}}} \exp \left[-\frac{\left(y_{t}-\mathbf{x}_{t}^{\prime} \boldsymbol{\beta}\right)^{2}}{2 \sigma^{2}}\right]\right) \\
& =-\frac{1}{2} \log (2 \pi)-\log (\sigma)-\frac{1}{2}\left(\frac{y_{t}-\mathbf{x}_{t}^{\prime} \boldsymbol{\beta}}{\sigma}\right)^{2} \tag{7.1.13}
\end{align*}


The parameter space $\Theta$ is $\mathbb{R}^{K} \times \mathbb{R}_{++}$, where $K$ is the dimension of $\boldsymbol{\beta}$ and $\mathbb{R}_{++}$is the set of positive real numbers reflecting the a priori restriction that $\sigma_{0}^{2}>0$. As was verified in Chapter 1, the ML estimator of $\boldsymbol{\beta}_{0}$ is the OLS estimator and the ML estimator of $\sigma_{0}^{2}$ is the sum of squared residuals divided by the sample size $n$.

Example 7.3 (Probit): In the probit model, a scalar dependent variable $y_{t}$ is a binary variable, $y_{t} \in\{0,1\}$. For example, $y_{t}=1$ if the wife chooses to enter the labor force and $y_{t}=0$ otherwise, and $\mathbf{x}_{t}$ might include the husband's\\
income. The conditional probability of $y_{t}$ given a vector of regressors $\mathbf{x}_{t}$ is given by

\[
\left\{\begin{array}{l}
f\left(y_{t}=1 \mid \mathbf{x}_{t} ; \boldsymbol{\theta}_{0}\right)=\Phi\left(\mathbf{x}_{t}^{\prime} \boldsymbol{\theta}_{0}\right)  \tag{7.1.14}\\
f\left(y_{t}=0 \mid \mathbf{x}_{t} ; \boldsymbol{\theta}_{0}\right)=1-\Phi\left(\mathbf{x}_{t}^{\prime} \boldsymbol{\theta}_{0}\right)
\end{array}\right.
\]

where $\Phi(\cdot)$ is the cumulative density function of the standard normal distribution. This can be written compactly as

$$
f\left(y_{t} \mid \mathbf{x}_{t} ; \boldsymbol{\theta}_{0}\right)=\Phi\left(\mathbf{x}_{t}^{\prime} \boldsymbol{\theta}_{0}\right)^{y_{t}}\left[1-\Phi\left(\mathbf{x}_{t}^{\prime} \boldsymbol{\theta}_{0}\right)\right]^{1-y_{t}}
$$

If there is no a priori restriction on $\boldsymbol{\theta}_{0}$, the parameter space $\Theta$ is $\mathbb{R}^{p}$. The ML estimator of $\boldsymbol{\theta}_{0}$ is an M-estimator with the $m$ function given by


\begin{equation*}
m\left(\mathbf{w}_{t} ; \boldsymbol{\theta}\right)=\log f\left(y_{t} \mid \mathbf{x}_{t} ; \boldsymbol{\theta}\right)=y_{t} \log \Phi\left(\mathbf{x}_{t}^{\prime} \boldsymbol{\theta}\right)+\left(1-y_{t}\right) \log \left[1-\Phi\left(\mathbf{x}_{t}^{\prime} \boldsymbol{\theta}\right)\right] \tag{7.1.15}
\end{equation*}


\section*{Invariance of ML}
The joint ML and conditional ML estimators have a number of desirable properties. One of them, which holds true in finite samples, is invariance. To describe invariance for an extremum estimator in general, consider reparameterizing the model by a mapping or function $\lambda=\boldsymbol{\tau}(\boldsymbol{\theta})$ defined on $\Theta$. Let $\Lambda$ be the range of the mapping:


\begin{equation*}
\Lambda \equiv \tau(\Theta) \equiv\{\lambda \mid \lambda=\tau(\theta) \text { for some } \theta \in \Theta\} \tag{7.1.16}
\end{equation*}


By definition, for every $\boldsymbol{\lambda}$ in $\Lambda$, there exists at least one $\boldsymbol{\theta}$ such that $\boldsymbol{\lambda}=\boldsymbol{\tau}(\boldsymbol{\theta})$. The mapping $\tau: \Theta \rightarrow \Lambda$ is called a reparameterization if it is one-to-one; namely, for every $\boldsymbol{\lambda}$ in $\Lambda$, there is only one $\boldsymbol{\theta}$ in $\Theta$ such that $\boldsymbol{\lambda}=\boldsymbol{\tau}(\boldsymbol{\theta})$. This unique $\boldsymbol{\theta}$ is therefore a function of $\lambda$. This mapping from $\Lambda$ to $\Theta$ is called the inverse of $\tau$ and denoted $\boldsymbol{\tau}^{-1}$. We say that an extremum estimator $\hat{\boldsymbol{\theta}}$ is invariant to reparameterization $\boldsymbol{\tau}$ if the extremum estimator for the reparameterized model is $\tau(\hat{\boldsymbol{\theta}})$. Let $\widetilde{Q}_{n}(\lambda)$ be the objective function associated with the reparameterized model. An extremum estimator is invariant if and only if


\begin{equation*}
\widetilde{Q}_{n}(\lambda)=Q_{n}\left(\tau^{-1}(\lambda)\right) \text { for all } \lambda \in \Lambda \tag{7.1.17}
\end{equation*}


To see this, let $\hat{\boldsymbol{\lambda}} \equiv \boldsymbol{\tau}(\hat{\boldsymbol{\theta}})$. For any $\boldsymbol{\lambda} \in \Lambda$, we have $\widetilde{Q}_{n}(\boldsymbol{\lambda})=Q_{n}\left(\boldsymbol{\tau}^{-1}(\boldsymbol{\lambda})\right) \leq Q_{n}(\hat{\boldsymbol{\theta}})$ because $\hat{\boldsymbol{\theta}}$ maximizes $Q_{n}(\boldsymbol{\theta})$ on $\Theta$ and $\boldsymbol{\tau}^{-1}(\lambda)$ is in $\Theta$. But $Q_{n}(\hat{\boldsymbol{\theta}})=Q_{n}\left(\boldsymbol{\tau}^{-1}(\hat{\boldsymbol{\lambda}})\right)= \widetilde{Q}_{n}(\hat{\lambda})$. So $\widetilde{Q}_{n}(\lambda) \leq \widetilde{Q}_{n}(\hat{\lambda})$ for all $\lambda \in \Lambda$.

ML (be it conditional or joint) is invariant to reparameterization because the likelihood after reparameterization is $\tilde{f}(. ; \lambda)=f\left(. ; \boldsymbol{\tau}^{-1}(\lambda)\right)$, which produces the objective functions that satisfy (7.1.17). Can there be an estimator that is not invariant? The answer is yes; Review Question 5 will ask you to verify that GMM is not invariant.

\section*{Nonlinear Least Squares (NLS)}
In NLS, as in conditional ML, the observation vector $\mathbf{w}_{t}$ is partitioned into two groups, $y_{t}$ and $\mathbf{x}_{t}: \mathbf{w}_{t}=\left(y_{t}, \mathbf{x}_{t}^{\prime}\right)^{\prime}$. The model in NLS is a set of stochastic processes $\left\{y_{t}, \mathbf{x}_{t}\right\}$ such that the conditional mean $\mathrm{E}\left(y_{t} \mid \mathbf{x}_{t}\right)$, which is a function of $\mathbf{x}_{t}$, is a member of the parametric family of functions $\varphi\left(\mathbf{x}_{t} ; \boldsymbol{\theta}\right), \boldsymbol{\theta} \in \Theta$. The functional form of $\varphi(\cdot ; \cdot)$ is known. If $\boldsymbol{\theta}_{0}$ is the true parameter value, then $\mathrm{E}\left(y_{t} \mid \mathbf{x}_{t}\right)=\varphi\left(\mathbf{x}_{t} ; \boldsymbol{\theta}_{0}\right)$ for the true DGP $\left\{y_{t}, \mathbf{x}_{t}\right\}$. If we define $\varepsilon_{t} \equiv y_{t}-\mathrm{E}\left(y_{t} \mid \mathbf{x}_{t}\right)$, then the correctly specified model can be written as


\begin{equation*}
y_{t}=\varphi\left(\mathbf{x}_{t} ; \boldsymbol{\theta}_{0}\right)+\varepsilon_{t}, \mathbf{E}\left(\varepsilon_{t} \mid \mathbf{x}_{t}\right)=0, \boldsymbol{\theta}_{0} \in \Theta . \tag{7.1.18}
\end{equation*}


The most widely used estimation method to estimate $\boldsymbol{\theta}_{0}$ here is least squares, which is to minimize the sum of squared residuals. Therefore, an NLS estimator, which is least squares applied to the model just described, is an M-estimator with


\begin{equation*}
m\left(\mathbf{w}_{t} ; \boldsymbol{\theta}\right)=-\left[y_{t}-\varphi\left(\mathbf{x}_{t} ; \boldsymbol{\theta}\right)\right]^{2}, \text { that is, } Q_{n}(\boldsymbol{\theta})=-\frac{1}{n} \sum_{t=1}^{n}\left[y_{t}-\varphi\left(\mathbf{x}_{t} ; \boldsymbol{\theta}\right)\right]^{2} . \tag{7.1.19}
\end{equation*}


Here maximization of $Q_{n}(\boldsymbol{\theta})$ is the same as minimization of the sum of squared residuals.

Example 7.4 (CES production function with additive errors): As an illustration, consider the CES production function


\begin{equation*}
Q_{t}=A_{0} \cdot\left[\delta_{0} K_{t}^{-\rho_{0}}+\left(1-\delta_{0}\right) L_{t}^{-\rho_{0}}\right]^{-1 / \rho_{0}}+\varepsilon_{t}, \tag{7.1.20}
\end{equation*}


where $Q_{t}$ is output in period $t, K_{t}$ is the capital stock, $L_{t}$ is labor input, and $\varepsilon_{t}$ is an additive shock to the production function with $\mathrm{E}\left(\varepsilon_{t} \mid K_{t}, L_{t}\right)=0$. This is a conditional mean model (7.1.18) with $y_{t}=Q_{t}, \mathbf{x}_{t}=\left(K_{t}, L_{t}\right)^{\prime}, \boldsymbol{\theta}_{0}= \left(A_{0}, \delta_{0}, \rho_{0}\right)^{\prime}, \varphi\left(\mathbf{x}_{t} ; \boldsymbol{\theta}\right)=A \cdot\left[\delta K_{t}^{-\rho}+(1-\delta) L_{t}^{-\rho}\right]^{-1 / \rho}$. The usual properties for production functions (of monotonicity and concavity) are satisfied if $A> 0,0<\delta<1,-1<\rho$, which determines the parameter space $\Theta$.

\section*{Linear and Nonlinear GMM}
In Chapter 3, we applied GMM to linear equations. The linear equation to be estimated was


\begin{equation*}
y_{t}=\mathbf{z}_{t}^{\prime} \boldsymbol{\theta}_{0}+\varepsilon_{t} . \tag{7.1.21}
\end{equation*}


With $\mathbf{x}_{t}$ as the vector of instruments, the orthogonality (zero-mean) conditions were


\begin{equation*}
\mathrm{E}\left[\mathbf{x}_{t} \cdot\left(y_{t}-\mathbf{z}_{t}^{\prime} \boldsymbol{\theta}_{0}\right)\right]=\mathbf{0} . \tag{7.1.22}
\end{equation*}


The correctly specified model here is a set of ergodic stationary processes $\mathbf{w}_{t}= \left(y_{t}, \mathbf{z}_{t}^{\prime}, \mathbf{x}_{t}^{\prime}\right)^{\prime}$ such that these zero-mean conditions hold for $\boldsymbol{\theta}_{0}$ in $\Theta$. The linear GMM estimator of $\boldsymbol{\theta}_{0}$ is a GMM estimator with the $\mathbf{g}$ function in the GMM objective function (7.1.3) given by


\begin{equation*}
\mathbf{g}\left(\mathbf{w}_{t} ; \boldsymbol{\theta}\right) \equiv \mathbf{x}_{t} \cdot\left(y_{t}-\mathbf{z}_{t}^{\prime} \boldsymbol{\theta}\right)=\mathbf{x}_{t} \cdot y_{t}-\mathbf{x}_{t} \mathbf{z}_{t}^{\prime} \boldsymbol{\theta} \tag{7.1.23}
\end{equation*}


GMM can be readily applied to nonlinear equations. Suppose that the estimation equation is nonlinear and also implicit in $y_{t}$ :


\begin{equation*}
a\left(y_{t}, \mathbf{z}_{t} ; \boldsymbol{\theta}_{0}\right)=\varepsilon_{t} . \tag{7.1.24}
\end{equation*}


Just set $\mathbf{g}\left(\mathbf{w}_{t} ; \boldsymbol{\theta}\right)=\mathbf{x}_{t} \cdot a\left(y_{t}, \mathbf{z}_{t} ; \boldsymbol{\theta}\right)$. This estimator is called a generalized nonlinear instrumental variables estimator. ${ }^{3}$ It is still a special case of GMM because the $\mathbf{g}$ function here can be written as a product of the vector of instruments and the error term. The properties of GMM estimators to be developed in this chapter do not rely on this special structure.

As an example of nonlinear GMM, consider the nonlinear Euler equation of Hansen and Singleton (1982).

Example 7.5 (Nonlinear consumption Euler equation): The Euler equation for the household optimization problem, standard in macroeconomics, is


\begin{equation*}
\mathrm{E}\left[\left.R_{t+1} \frac{\beta_{0} u^{\prime}\left(c_{t+1}\right)}{u^{\prime}\left(c_{t}\right)}\right|_{I_{t}}\right]=1, \tag{7.1.25}
\end{equation*}


where $R_{t+1}$ is the gross ex-post rate of return (1 plus the rate of return) from the asset in question, $c_{t}$ is consumption, $\beta_{0}$ is the discounting factor, $u^{\prime}(c)$ is the marginal utility, and $I_{t}$ is the information available at date $t$. Assume for

\footnotetext{${ }^{3}$ It is more general than the usual nonlinear instrumental variables estimator because conditional homoskedasticity is not assumed.
}
the utility function that $u(c)=c^{1-\alpha_{0}} /\left(1-\alpha_{0}\right)$. Then $u^{\prime}(c)=c^{-\alpha_{0}}$ and the Euler equation becomes

\begin{align*}
& \mathrm{E}\left[\left.a\left(\frac{c_{t+1}}{c_{t}}, R_{t+1} ; \alpha_{0}, \beta_{0}\right) \right\rvert\, I_{t}\right]=0 \\
& \quad \text { with } a\left(\frac{c_{t+1}}{c_{t}}, R_{t+1} ; \alpha_{0}, \beta_{0}\right) \equiv R_{t+1} \cdot \beta_{0} \cdot\left(\frac{c_{t+1}}{c_{t}}\right)^{-\alpha_{0}}-1 \tag{7.1.26}
\end{align*}


If $\mathbf{x}_{t}$ is a vector of variables whose values are known at date $t$, then $\mathbf{x}_{t} \in I_{t}$ and it follows from (7.1.26) that


\begin{equation*}
\mathrm{E}\left[\left.\mathbf{x}_{t} \cdot a\left(\frac{c_{t+1}}{c_{t}}, R_{t+1} ; \alpha_{0}, \beta_{0}\right) \right\rvert\, I_{t}\right]=\mathbf{0} \tag{7.1.27}
\end{equation*}


Taking the unconditional expectation of both sides and using the Law of Total Expectations (that $\left.\mathrm{E}\left[\mathrm{E}\left(x \mid I_{t}\right)\right]=\mathrm{E}(x)\right)$, we obtain the orthogonality (zeromean) conditions $\mathrm{E}\left[\mathbf{g}\left(\mathbf{w}_{t} ; \boldsymbol{\theta}_{0}\right)\right]=\mathbf{0}$ where

\[
\mathbf{g}\left(\mathbf{w}_{t} ; \boldsymbol{\theta}_{0}\right)=\mathbf{x}_{t} \cdot a\left(\frac{c_{t+1}}{c_{t}}, R_{t+1} ; \alpha_{0}, \beta_{0}\right) \text { with } \mathbf{w}_{t}=\left[\begin{array}{c}
\mathbf{x}_{t}  \tag{7.1.28}\\
c_{t+1} / c_{t} \\
R_{t+1}
\end{array}\right], \boldsymbol{\theta}_{0}=\left[\begin{array}{l}
\beta_{0} \\
\alpha_{0}
\end{array}\right] .
\]

\section*{QUESTIONS FOR REVIEW}
\begin{enumerate}
  \item (Estimating probit model by NLS) For the probit model, verify that $\mathrm{E}\left(y_{t} \mid\right. \left.\mathbf{x}_{t}\right)=\Phi\left(\mathbf{x}_{t}^{\prime} \boldsymbol{\theta}_{0}\right)$. Consider applying least squares to this model. Specify the $m$ function in the M-estimator objective function (7.1.2).
\end{enumerate}

2 . (Estimating probit model by GMM) For the probit model, verify that the orthogonality conditions $\mathrm{E}\left[\mathbf{x}_{t} \cdot\left(y_{t}-\Phi\left(\mathbf{x}_{t}^{\prime} \boldsymbol{\theta}_{0}\right)\right)\right]=\mathbf{0}$ hold. Specify the $\mathbf{g}$ function in the GMM objective function (7.1.3).\\
3. (Estimation by ML of a GMM model?) Suppose that $\left\{y_{t}, \mathbf{z}_{t}, \mathbf{x}_{t}\right\}$ is i.i.d. in the correctly specified linear model described in the last subsection. Can you estimate $\boldsymbol{\theta}_{0}$ by ML? [Answer: No, because the model does not specify the parametric form for the density of $\mathbf{w}_{t}$.]\\
4. ( $\mathrm{E}\left(\varepsilon_{t} \mid \mathbf{x}_{t}\right)=0$ vs. $\mathrm{E}\left(\varepsilon_{t} \cdot \mathbf{x}_{t}\right)=\mathbf{0}$ in NLS) Consider the quadratic NLS model where $\varphi\left(x_{t} ; \boldsymbol{\theta}\right)=\theta_{0}+\theta_{1} x_{t}+\theta_{2} x_{t}^{2}$. Suppose that $\mathrm{E}\left(\varepsilon_{t} x_{t}\right)=0$ but $\mathrm{E}\left(\varepsilon_{t} x_{t}^{2}\right) \neq 0$. Does $\mathrm{E}\left(\varepsilon_{t} \mid x_{t}\right)=0$ hold? [Answer: No. If $\mathrm{E}\left(\varepsilon_{t} \mid x_{t}\right)=0$, then $\mathrm{E}\left(\varepsilon_{t} h\left(x_{t}\right)\right)=0$\\
for any (measurable) function $h$.] Would the NLS estimator of $\boldsymbol{\theta}_{0}$ be consistent? [Answer: No.]\\
5. (Lack of invariance in GMM) Consider the linear model $y_{t}=\theta_{0} z_{t}+\varepsilon_{t}$ where $\theta_{0}$ and $z_{t}$ are scalars. The orthogonality conditions are $\mathrm{E}\left[\mathbf{x}_{t} \cdot\left(y_{t}-\theta_{0} z_{t}\right)\right]=\mathbf{0}$. Write down the GMM objective function $Q_{n}(\theta)$. Assume that $\Theta=\mathbb{R}_{++}$(i.e., $\theta_{0}>0$ ) and consider the reparameterization $\lambda=1 / \theta$. With this reparameterization, the linear equation can be rewritten as $z_{t}=\lambda_{0} y_{t}-\lambda_{0} \varepsilon_{t}$ and the orthogonality conditions can be rewritten as $\mathrm{E}\left[\mathbf{x}_{t} \cdot\left(z_{t}-\lambda_{0} y_{t}\right)\right]=\mathbf{0}$. Let $\widetilde{Q}_{n}(\lambda)$ be the GMM objective functions associated with this set of orthogonality conditions. Is it the case that $Q_{n}(\theta)=\widetilde{Q}_{n}(1 / \theta)$ or $\widetilde{Q}_{n}(\lambda)=Q_{n}(1 / \lambda)$ when $\widehat{\mathbf{W}}$ is the same in the two objective functions? [Answer: No.]

\subsection*{7.2 Consistency}
In this section, we first present a set of sufficient conditions under which an extremum estimator is consistent. Those conditions will then be specialized to NLS, ML, and GMM.

\section*{Two Consistency Theorems for Extremum Estimators}
The objective function $Q_{n}(\cdot)$ is a random function because for each $\theta$ its value $Q_{n}(\boldsymbol{\theta})$, being dependent on the data $\left(\mathbf{w}_{1}, \ldots, \mathbf{w}_{n}\right)$, is a random variable. The basic idea for the consistency of an extremum estimator is that if $Q_{n}(\boldsymbol{\theta})$ converges in probability to $Q_{0}(\boldsymbol{\theta})$, and the true parameter $\boldsymbol{\theta}_{0}$ solves the "limit problem" of maximizing the limit function $Q_{0}(\boldsymbol{\theta})$, then the limit of the maximum $\hat{\boldsymbol{\theta}}$ should be $\boldsymbol{\theta}_{0}$. Sufficient conditions for the maximum of the limit to be the limit of the maximum are that the convergence in probability be "uniform" (in the sense made precise in a moment) and that the parameter space $\Theta$ be compact. As already mentioned, the parameter space is not compact in most applications. We therefore provide an alternative set of consistency conditions, which does not require $\Theta$ to be compact.

To state the first consistency theorem for extremum estimators, we need to make precise what we mean by convergence being "uniform." If you are already familiar from calculus with the notion of uniform convergence of a sequence of functions, you will recognize that the definition about to be given is a natural extension to a sequence of random functions, $Q_{n}(\cdot)(n=1,2, \ldots)$. Pointwise convergence in probability of $Q_{n}(\cdot)$ to some nonrandom function $Q_{0}(\cdot)$ simply means $\operatorname{plim}_{n \rightarrow \infty} Q_{n}(\boldsymbol{\theta})=Q_{0}(\boldsymbol{\theta})$ for all $\boldsymbol{\theta}$, namely, that the sequence of random\\
variables $\left|Q_{n}(\boldsymbol{\theta})-Q_{0}(\boldsymbol{\theta})\right|(n=1,2, \ldots)$ converges in probability to 0 for each $\theta$. Uniform convergence in probability is stronger. The convergence has to occur "uniformly" over the parameter space $\Theta$ in the following sense:


\begin{equation*}
\sup _{\boldsymbol{\theta} \in \Theta}\left|Q_{n}(\boldsymbol{\theta})-Q_{0}(\boldsymbol{\theta})\right| \underset{\mathrm{p}}{\rightarrow} 0 \text { as } n \rightarrow \infty . \tag{7.2.1}
\end{equation*}


As in the case of convergence in probability, uniform convergence in probability can be extended readily to sequences of vector random functions or matrix random functions (by viewing a matrix as a vector whose elements have been rearranged), by requiring uniform convergence for each element: a sequence of vector random functions $\left\{\mathbf{h}_{n}(\cdot)\right\}$ converges uniformly in probability to a nonrandom function $\mathbf{h}_{0}(\cdot)$ if each element converges uniformly. This element-by-element convergence is equivalent to convergence in the norm:


\begin{equation*}
\sup _{\boldsymbol{\theta} \in \Theta}\left\|\mathbf{h}_{n}(\boldsymbol{\theta})-\mathbf{h}_{0}(\boldsymbol{\theta})\right\| \underset{\mathrm{p}}{\rightarrow} 0 \text { as } n \rightarrow \infty \text {, } \tag{7.2.2}
\end{equation*}


where $\|$.$\| is the usual Euclidean norm. { }^{4}$\\
In the following statement of our first consistency theorem, the conditions ensuring that $\hat{\boldsymbol{\theta}}$ is well defined are labeled (i)-(iii), while those ensuring consistency are (a) and (b).

Proposition 7.1 (Consistency with compact parameter space): Suppose that (i) $\Theta$ is a compact subset of $\mathbb{R}^{p}$, (ii) $Q_{n}(\boldsymbol{\theta})$ is continuous in $\boldsymbol{\theta}$ for any data $\left(\mathbf{w}_{1}, \ldots\right.$, $\mathbf{w}_{n}$ ), and (iii) $Q_{n}(\boldsymbol{\theta})$ is a measurable function of the data for all $\boldsymbol{\theta}$ in $\Theta$ (so by Lemma 7.1 the extremum estimator $\hat{\boldsymbol{\theta}}$ defined in (7.1.1) is a well-defined random variable). If there is a function $Q_{0}(\boldsymbol{\theta})$ such that\\
(a) (identification) $\quad Q_{0}(\boldsymbol{\theta})$ is uniquely maximized on $\Theta$ at $\boldsymbol{\theta}_{0} \in \Theta$,\\
(b) (uniform convergence) $Q_{n}(\cdot)$ converges uniformly in probability to $Q_{0}(\cdot)$, then $\hat{\boldsymbol{\theta}} \rightarrow_{\mathrm{p}} \boldsymbol{\theta}_{0}$.

We will not prove this result; see, e.g., Amemiya (1985, Theorem 4.1.1) for a proof. For estimates such as ML, NLS, and GMM, we will state later the identification condition (a) in more primitive terms so that the condition can be interpreted more intuitively.

\footnotetext{${ }^{4}$ The Euclidean norm of a $p$-dimensional vector $\mathbf{x}=\left(x_{1}, x_{2}, \ldots, x_{p}\right)^{\prime}$, denoted $\|\mathbf{x}\|$, is defined as $\sqrt{x_{1}^{2}+x_{2}^{2}+\cdots+x_{p}^{2}}$
}The theorem that does away with the compactness of $\Theta$ is

Proposition 7.2 (Consistency without compactness): Suppose that (i) the true parameter vector $\boldsymbol{\theta}_{0}$ is an element of the interior of a convex parameter space $\Theta \left(\subset \mathbb{R}^{p}\right)$, (ii) $Q_{n}(\boldsymbol{\theta})$ is concave ${ }^{5}$ over the parameter space for any data $\left(\mathbf{w}_{1}, \ldots, \mathbf{w}_{n}\right)$, and (iii) $Q_{n}(\boldsymbol{\theta})$ is a measurable function of the data for all $\boldsymbol{\theta}$ in $\Theta$. Let $\hat{\boldsymbol{\theta}}$ be the extremum estimator defined by (7.1.1) (wait for a moment to learn about its existence). If there is a function $Q_{0}(\boldsymbol{\theta})$ such that\\
(a) (identification) $Q_{0}(\boldsymbol{\theta})$ is uniquely maximized on $\Theta$ at $\boldsymbol{\theta}_{0} \in \Theta$,\\
(b) (pointwise convergence) $Q_{n}(\theta)$ converges in probability to $Q_{0}(\theta)$ for all $\boldsymbol{\theta} \in \Theta$,\\
then, as $n \rightarrow \infty, \hat{\boldsymbol{\theta}}$ exists with probability approaching 1 and $\hat{\boldsymbol{\theta}} \rightarrow_{\mathrm{p}} \boldsymbol{\theta}_{0}$.\\
See Newey and McFadden (1994, pp. 2133-2134) for a proof. In most applications, $\Theta$ is an open convex set (as in Examples 7.1-7.4), so condition (i) is satisfied. The convergence of $Q_{n}(\boldsymbol{\theta})$ to $Q_{0}(\boldsymbol{\theta})$ is required to be only pointwise, which is easier to verify in applications. The price of these niceties is the requirement that the objective function $Q_{n}(\boldsymbol{\theta})$ be concave for any data (under which pointwise convergence implies uniform convergence), but in many applications this condition is satisfied. Below we will verify that concavity is satisfied for ML for the linear regression model (after a reparameterization), probit ML, and linear GMM.

\section*{Consistency of M-Estimators}
For an M-estimator, the objective function is given by (7.1.2). If $\left\{\mathbf{w}_{t}\right\}$ is ergodic stationary, the Ergodic Theorem implies pointwise convergence for $Q_{n}(\boldsymbol{\theta})$ : $Q_{n}(\boldsymbol{\theta})$ converges in probability pointwise for each $\boldsymbol{\theta} \in \Theta$ to $Q_{0}(\boldsymbol{\theta})$ given by


\begin{equation*}
Q_{0}(\boldsymbol{\theta})=\mathrm{E}\left[m\left(\mathbf{w}_{t} ; \boldsymbol{\theta}\right)\right] \tag{7.2.3}
\end{equation*}


(provided, of course, that $\mathrm{E}\left[m\left(\mathbf{w}_{t} ; \boldsymbol{\theta}\right)\right]$ exists and is finite). To apply the first consistency result (Proposition 7.1) to M-estimators, we need to show that the convergence of $\frac{1}{n} \sum_{t=1}^{n} m\left(\mathbf{w}_{t} ;.\right)$ to $\mathrm{E}\left[m\left(\mathbf{w}_{t} ;.\right)\right]$ is uniform. A standard method to prove uniform convergence in probability, which exploits the fact that the expression is

\footnotetext{${ }^{5}$ A scalar function $h(\mathbf{x})$ is said to be concave if $\alpha h\left(\mathbf{x}_{1}\right)+(1-\alpha) h\left(\mathbf{x}_{2}\right) \leq h\left(\alpha \mathbf{x}_{1}+(1-\alpha) \mathbf{x}_{2}\right)$ for all $\mathbf{x}_{1}, \mathbf{x}_{2}, 0 \leq \alpha \leq 1$. So a linear function is concave. Some authors use the term "globally concave" to distinguish it from the concept of "local concavity" that the Hessian evaluated at the parameter vector in question is negative semidefinite. If a function is concave, then it is continuous. So condition (ii) in this proposition is stronger than condition (ii) in the previous proposition.
}
a sample mean, is the uniform law of large numbers. Its version for ergodic stationary processes (stated as Lemma 2.4 in Newey and McFadden, 1994, and as Theorem A.2.2 in White, 1994) is

Lemma 7.2 (Uniform law of large numbers): Let $\left\{\mathbf{w}_{t}\right\}$ be an ergodic stationary process. Suppose that (i) the set $\Theta$ is compact, (ii) $m\left(\mathbf{w}_{t} ; \boldsymbol{\theta}\right)$ is continuous in $\boldsymbol{\theta}$ for all $\mathbf{w}_{t}$, and (iii) $m\left(\mathbf{w}_{t} ; \boldsymbol{\theta}\right)$ is measurable in $\mathbf{w}_{t}$ for all $\boldsymbol{\theta}$ in $\Theta$. Suppose, in addition, the\\
dominance condition: there exists a function $d\left(\mathbf{w}_{t}\right)$ (sometimes called the "dominating function") such that $\left|m\left(\mathbf{w}_{t} ; \boldsymbol{\theta}\right)\right| \leq d\left(\mathbf{w}_{t}\right)$ for all $\boldsymbol{\theta} \in \Theta$ and $\mathrm{E}\left[d\left(\mathbf{w}_{t}\right)\right]<\infty$.

Then $\frac{1}{n} \sum_{t=1}^{n} m\left(\mathbf{w}_{t} ;.\right)$ converges uniformly in probability to $\mathrm{E}\left[m\left(\mathbf{w}_{t} ;.\right)\right]$ over $\Theta$. Moreover, $\mathrm{E}\left[m\left(\mathbf{w}_{t} ; \boldsymbol{\theta}\right)\right]$ is a continuous function of $\boldsymbol{\theta}$.

Since by construction $\left|m\left(\mathbf{w}_{t} ; \boldsymbol{\theta}\right)\right| \leq \sup _{\boldsymbol{\theta} \in \Theta}\left|m\left(\mathbf{w}_{t} ; \boldsymbol{\theta}\right)\right|$ for all $\boldsymbol{\theta} \in \Theta$, we can use $\sup _{\boldsymbol{\theta} \in \Theta}\left|m\left(\mathbf{w}_{t} ; \boldsymbol{\theta}\right)\right|$ for $d\left(\mathbf{w}_{t}\right)$, so the above dominance condition can be stated simply as $\mathrm{E}\left[\sup _{\boldsymbol{\theta} \in \Theta}\left|m\left(\mathbf{w}_{t} ; \boldsymbol{\theta}\right)\right|\right]<\infty$. The Uniform Law of Large Numbers can be extended immediately to vector random functions as follows:

Let $\left\{\mathbf{w}_{l}\right\}$ be an ergodic stationary process. Suppose that (i) the parameter space $\Theta$ is compact, (ii) $\mathbf{h}\left(\mathbf{w}_{t} ; \boldsymbol{\theta}\right)$ is continuous in $\boldsymbol{\theta}$ for all $\mathbf{w}_{t}$, and (iii) $\mathbf{h}\left(\mathbf{w}_{t} ; \boldsymbol{\theta}\right)$ is measurable in $\mathbf{w}_{t}$ for all $\boldsymbol{\theta}$ in $\Theta$. Suppose, in addition, the

$$
\text { dominance condition: } \mathrm{E}\left[\sup _{\boldsymbol{\theta} \in \Theta}\left\|\mathbf{h}\left(\mathbf{w}_{t} ; \boldsymbol{\theta}\right)\right\|\right]<\infty \text {. }
$$

Then $\mathrm{E}\left[\mathbf{h}\left(\mathbf{w}_{t} ; \boldsymbol{\theta}\right)\right]$ is a continuous function of $\boldsymbol{\theta}$ and $\frac{1}{n} \sum_{t=1}^{n} \mathbf{h}\left(\mathbf{w}_{t} ;.\right)$ converges uniformly in probability to $\mathrm{E}\left[\mathbf{h}\left(\mathbf{w}_{i} ;.\right)\right]$ over $\Theta$.

Just by combining this lemma with the first consistency theorem for extremum estimators, and noting that $Q_{n}(\boldsymbol{\theta})$ is continuous if $m\left(\mathbf{w}_{t} ; \boldsymbol{\theta}\right)$ is continuous in $\boldsymbol{\theta}$ and that $Q_{n}(\boldsymbol{\theta})$ is a measurable function of the data if $m\left(\mathbf{w}_{t} ; \boldsymbol{\theta}\right)$ is measurable in $\mathbf{w}_{t}$, we obtain

\section*{Proposition 7.3 (Consistency of M-estimators with compact parameter space):}
Let $\left\{\mathbf{w}_{t}\right\}$ be ergodic stationary. Suppose that (i) the parameter space $\Theta$ is a compact subset of $\mathbb{R}^{p}$, (ii) $m\left(\mathbf{w}_{t} ; \boldsymbol{\theta}\right)$ is continuous in $\boldsymbol{\theta}$ for all $\mathbf{w}_{t}$, and (iii) $m\left(\mathbf{w}_{t} ; \boldsymbol{\theta}\right)$ is measurable in $\mathbf{w}_{t}$ for all $\boldsymbol{\theta}$ in $\Theta$. Let $\hat{\boldsymbol{\theta}}$ be the M-estimator defined by (7.1.1) and (7.1.2). Suppose, further, that\\
(a) (identification) $\mathrm{E}\left[m\left(\mathbf{w}_{i} ; \boldsymbol{\theta}\right)\right]$ is uniquely maximized on $\Theta$ at $\boldsymbol{\theta}_{0} \in \Theta$,\\
(b) (dominance) $\mathrm{E}\left[\sup _{\boldsymbol{\theta} \in \Theta}\left|m\left(\mathbf{w}_{t} ; \boldsymbol{\theta}\right)\right|\right]<\infty$.

Then $\hat{\boldsymbol{\theta}} \rightarrow_{\mathrm{p}} \boldsymbol{\theta}_{0}$.

It is even more straightforward to specialize the second consistency theorem for extremum estimators to M-estimators. The objective function $Q_{n}(\theta)$ is concave if $m\left(\mathbf{w}_{t} ; \boldsymbol{\theta}\right)$ is concave in $\boldsymbol{\theta}$. ${ }^{6}$ As noted above, the pointwise convergence of $Q_{n}(\boldsymbol{\theta})$ to $Q_{0}(\boldsymbol{\theta})\left(=\mathrm{E}\left[m\left(\mathbf{w}_{t} ; \boldsymbol{\theta}\right)\right]\right)$ is assured by the Ergodic Theorem if $\mathrm{E}\left[m\left(\mathbf{w}_{t} ; \boldsymbol{\theta}\right)\right]$ exists and is finite for all $\theta$. Thus

Proposition 7.4 (Consistency of M-estimators without compactness): Let $\left\{\mathbf{w}_{t}\right\}$ be ergodic stationary. Suppose that (i) the true parameter vector $\boldsymbol{\theta}_{0}$ is an element of the interior of a convex parameter space $\Theta\left(\subset \mathbb{R}^{p}\right)$, (ii) $m\left(\mathbf{w}_{t} ; \boldsymbol{\theta}\right)$ is concave over the parameter space for all $\mathbf{w}_{t}$, and (iii) $m\left(\mathbf{w}_{t} ; \boldsymbol{\theta}\right)$ is measurable in $\mathbf{w}_{t}$ for all $\boldsymbol{\theta}$ in $\Theta$. Let $\hat{\theta}$ be the M-estimator defined by (7.1.1) and (7.1.2). Suppose, further, that\\
(a) (identification) $\mathrm{E}\left[m\left(\mathbf{w}_{t} ; \boldsymbol{\theta}\right)\right]$ is uniquely maximized on $\Theta$ at $\boldsymbol{\theta}_{0} \in \Theta$,\\
(b) $\mathrm{E}\left[\left|m\left(\mathbf{w}_{t} ; \boldsymbol{\theta}\right)\right|\right]<\infty$ (i.e., $\mathrm{E}\left[m\left(\mathbf{w}_{t} ; \boldsymbol{\theta}\right)\right]$ exists and is finite) for all $\boldsymbol{\theta} \in \Theta$.

Then, as $n \rightarrow \infty, \hat{\boldsymbol{\theta}}$ exists with probability approaching 1 and $\hat{\boldsymbol{\theta}} \rightarrow_{\mathrm{p}} \boldsymbol{\theta}_{0}$.

The following example shows that probit conditional ML satisfies the concavity condition.

Example 7.6 (Concavity of the probit likelihood function): The $m$ function for probit is given in (7.1.15) where $\Phi(\cdot)$ is the cumulative density function of the standard normal distribution. Since the normal distribution is symmetric, $\Phi(-v)=1-\Phi(v)$ and the $m$ function can be rewritten as


\begin{equation*}
m\left(\mathbf{w}_{t} ; \boldsymbol{\theta}\right)=y_{t} \log \Phi\left(\mathbf{x}_{t}^{\prime} \boldsymbol{\theta}\right)+\left(1-y_{t}\right) \log \Phi\left(-\mathbf{x}_{t}^{\prime} \boldsymbol{\theta}\right) . \tag{7.2.4}
\end{equation*}


Concavity of $m$ is implied if both $\log \Phi\left(\mathbf{x}_{t}^{\prime} \boldsymbol{\theta}\right)$ and $\log \Phi\left(-\mathbf{x}_{t}^{\prime} \boldsymbol{\theta}\right)$ are concave in $\boldsymbol{\theta}$. Since $\mathbf{x}_{t}^{\prime} \boldsymbol{\theta}$ is linear in $\boldsymbol{\theta}$, it suffices to show concavity of $\log \Phi(v)$. The first derivative of $\log \Phi(v)$ is


\begin{equation*}
\frac{\mathrm{d} \log \Phi(v)}{\mathrm{d} v}=\frac{\phi(v)}{\Phi(v)}, \tag{7.2.5}
\end{equation*}


\footnotetext{${ }^{6}$ A fact from calculus: if $f(\mathbf{x})$ and $g(\mathbf{x})$ are concave functions, then so is $f(\mathbf{x})+g(\mathbf{x})$.
}
where $\phi(v)\left(=\Phi^{\prime}(v)\right)$ is the density function of $N(0,1)$. It is well known that $\phi(v) / \Phi(v)$ is a decreasing function. So $\log \Phi(v)$ is (strictly) concave.

\section*{Concavity after Reparameterization}
Proposition 7.4 can be extended easily to estimators for objective functions that are concave after reparameterization. Suppose that all the conditions of the proposition except concavity are satisfied and that there is a continuous one-to-one mapping $\boldsymbol{\tau}(\boldsymbol{\theta}): \Theta \rightarrow \Lambda \equiv \boldsymbol{\tau}(\Theta)$ such that

$$
\widetilde{m}\left(\mathbf{w}_{t} ; \lambda\right) \equiv m\left(\mathbf{w}_{t} ; \boldsymbol{\tau}^{-1}(\lambda)\right)
$$

is concave in $\lambda$ and $\Lambda=\tau(\Theta)$ is a convex set. Let

$$
\widetilde{Q}_{n}(\lambda) \equiv \frac{1}{n} \sum_{t=1}^{n} \widetilde{m}\left(\mathbf{w}_{t} ; \lambda\right)
$$

be the objective function after this reparameterization. Clearly, all the assumptions of Proposition 7.4 are satisfied for $\lambda, \Lambda$, and $\tilde{m}\left(\mathbf{w}_{r} ; \lambda\right): \lambda_{0} \equiv \boldsymbol{\tau}\left(\boldsymbol{\theta}_{0}\right)$ is an interior point of $\Lambda, \mathrm{E}\left[\widetilde{m}\left(\mathbf{w}_{t} ; \lambda\right)\right]$ is uniquely maximized at $\lambda_{0}$, and $\mathrm{E}\left[\left|\widetilde{m}\left(\mathbf{w}_{t} ; \lambda\right)\right|\right]<\infty$ for all $\lambda$ in $\Lambda$. Thus for sufficiently large $n$, there exists an M-estimator $\hat{\lambda}$ such that $\operatorname{plim}_{n \rightarrow \infty} \hat{\lambda}=\lambda_{0}$. Since $\widetilde{Q}_{n}(\lambda)$ satisfies (7.1.17) for the one-to-one mapping $\tau$, $\hat{\boldsymbol{\theta}} \equiv \boldsymbol{\tau}^{-1}(\hat{\boldsymbol{\lambda}})$ maximizes the original objective function $Q_{n}(\boldsymbol{\theta})$. The estimator $\hat{\boldsymbol{\theta}}$ is consistent because


\begin{align*}
\operatorname{plim}_{n \rightarrow \infty} \hat{\boldsymbol{\theta}} & =\operatorname{plim}_{n \rightarrow \infty} \boldsymbol{\tau}^{-1}(\hat{\boldsymbol{\lambda}}) \\
& =\boldsymbol{\tau}^{-1}(\operatorname{plim} \hat{\boldsymbol{\lambda}}) \quad \text { (since } \boldsymbol{\tau}^{-1} \text { is continuous) } \\
& =\boldsymbol{\tau}^{-1}\left(\boldsymbol{\lambda}_{0}\right)=\boldsymbol{\theta}_{0} . \tag{7.2.6}
\end{align*}


Example 7.7 (Concavity of linear regression likelihood function): In the conditional ML estimation of the linear regression model, $m\left(\mathbf{w}_{t} ; \boldsymbol{\theta}\right)$ is given in (7.1.13) with $\boldsymbol{\theta}=\left(\boldsymbol{\beta}^{\prime}, \sigma^{2}\right)^{\prime}$ and $\mathbf{w}_{t}=\left(y_{t}, \mathbf{x}_{t}^{\prime}\right)^{\prime}$. This function is not necessarily concave, ${ }^{7}$ but consider a reparameterization: $\gamma=1 / \sigma$ and $\boldsymbol{\delta}=\boldsymbol{\beta} / \sigma$. It is a continuous one-to-one mapping between $\Theta=\mathbb{R}^{\boldsymbol{K}} \times \mathbb{R}_{++}$and $\Lambda=\mathbb{R}^{K} \times \mathbb{R}_{++}$. The new parameter space $\Lambda$ is convex. With $\lambda=\left(\boldsymbol{\delta}^{\prime}, \gamma\right)^{\prime}$,

\footnotetext{${ }^{7}$ For example, the second partial derivative of $m$ with respect to $\sigma^{2}$ may or may not be non-negative depending on the value of $y_{t}-\mathbf{x}_{t}^{\prime} \boldsymbol{\beta}$. Recall that a necessary condition for a twice differentiable function $h\left(x_{1}, \ldots, x_{p}\right)$ to be concave is that $\partial^{2} h /\left(\partial x_{j}\right)^{2}$ be non-negative for all $j=1,2, \ldots, p$.
}
the $m$ function after reparameterization becomes

\begin{equation*}
\widetilde{m}\left(\mathbf{w}_{t} ; \lambda\right)=-\frac{1}{2} \log (2 \pi)+\log (\gamma)-\frac{1}{2}\left(\gamma y_{t}-\mathbf{x}_{t}^{\prime} \boldsymbol{\delta}\right)^{2} . \tag{7.2.7}
\end{equation*}


Since $\log (\gamma)$ is concave in $\gamma$, it is concave in $\lambda$. Since $v \equiv \gamma y_{t}-\mathbf{x}_{t}^{\prime} \boldsymbol{\delta}$ is linear in $\boldsymbol{\lambda}$ and $-\frac{1}{2} v^{2}$ is concave, $-\frac{1}{2}\left(\gamma y_{t}-\mathbf{x}_{t}^{\prime} \boldsymbol{\delta}\right)^{2}$ is concave in $\boldsymbol{\lambda}$. So $\widetilde{m}$ above, which is the sum of two concave functions (plus a constant), is concave in $\boldsymbol{\lambda}$ for any $\mathbf{w}_{t}$. Therefore, provided that all the conditions of Proposition 7.4 besides concavity are satisfied for $\boldsymbol{\theta}, \Theta$, and $m\left(\mathbf{w}_{t} ; \boldsymbol{\theta}\right)$, the conditional ML estimator $\hat{\boldsymbol{\theta}}$ is consistent.

\section*{Identification in NLS and ML}
The identification condition (a) (that $\mathrm{E}\left[m\left(\mathbf{w}_{t} ; \boldsymbol{\theta}\right)\right]$ be uniquely maximized at $\boldsymbol{\theta}_{0}$ ) in these consistency theorems for M -estimators can be stated in more primitive terms for NLS and ML.

\section*{NLS}
Consider first NLS, where $m\left(\mathbf{w}_{t} ; \boldsymbol{\theta}\right)=-\left(y_{t}-\varphi\left(\mathbf{x}_{t} ; \boldsymbol{\theta}\right)\right)^{2}$ and $\mathrm{E}\left(y_{t} \mid \mathbf{x}_{t}\right)=\varphi\left(\mathbf{x}_{t} ; \boldsymbol{\theta}_{0}\right)$. We have shown in Section 2.9 that the conditional expectation minimizes the mean squared error, that is, for any (measurable) function $h\left(\mathbf{x}_{t}\right)$,


\begin{align*}
\text { mean squared error } & \equiv \mathrm{E}\left[\left\{y_{t}-h\left(\mathbf{x}_{t}\right)\right\}^{2}\right] \\
& =\mathrm{E}\left[\left\{y_{t}-\mathrm{E}\left(y_{t} \mid \mathbf{x}_{t}\right)\right\}^{2}\right]+\mathrm{E}\left[\left\{\mathrm{E}\left(y_{t} \mid \mathbf{x}_{t}\right)-h\left(\mathbf{x}_{t}\right)\right\}^{2}\right] \\
& \geq \mathrm{E}\left[\left\{y_{t}-\mathrm{E}\left(y_{t} \mid \mathbf{x}_{t}\right)\right\}^{2}\right] \tag{7.2.8}
\end{align*}


This is just reproducing (2.9.4) in the current notation. Furthermore, the conditional expectation is essentially the unique minimizer in the following sense.

Being functions of $\mathbf{x}_{t}, h\left(\mathbf{x}_{t}\right)$ and $\mathrm{E}\left(y_{t} \mid \mathbf{x}_{t}\right)$ are random variables. For two random variables $X$ and $Y$, let $X \neq Y$ mean $\operatorname{Prob}(X \neq Y)>0$ (i.e., $\operatorname{Prob}(X=Y)<1$ ) and let $X=Y$ mean $\operatorname{Prob}(X \neq Y)=0$. (If $X$ and $Y$ differ only for events whose probability is zero, we do not need to care about the difference.) Therefore, $h\left(\mathbf{x}_{t}\right) \neq \mathrm{E}\left(y_{t} \mid \mathbf{x}_{t}\right)$ means $\operatorname{Prob}\left(h\left(\mathbf{x}_{t}\right) \neq \mathrm{E}\left(y_{t} \mid \mathbf{x}_{t}\right)\right)>0$. If $h\left(\mathbf{x}_{t}\right) \neq \mathrm{E}\left(y_{t} \mid \mathbf{x}_{t}\right)$ in this sense, then $\left\{\mathrm{E}\left(y_{t} \mid \mathbf{x}_{t}\right)-h\left(\mathbf{x}_{t}\right)\right\}^{2}>0$ with positive probability. Consequently, $\mathrm{E}\left[\left\{\mathrm{E}\left(y_{t} \mid \mathbf{x}_{t}\right)-h\left(\mathbf{x}_{t}\right)\right\}^{2}\right]$ in (7.2.8) is positive and the mean squared error is strictly greater than $\mathrm{E}\left[\left\{y_{t}-\mathrm{E}\left(y_{t} \mid \mathbf{x}_{t}\right)\right\}^{2}\right]$.

Just setting $h\left(\mathbf{x}_{t}\right)=\varphi\left(\mathbf{x}_{t} ; \boldsymbol{\theta}\right)$ and noting that $\varphi\left(\mathbf{x}_{t} ; \boldsymbol{\theta}_{0}\right)=\mathrm{E}\left(y_{t} \mid \mathbf{x}_{t}\right)$, we conclude that

\[
\begin{cases}\mathrm{E}\left[\left\{y_{t}-\varphi\left(\mathbf{x}_{t} ; \boldsymbol{\theta}\right)\right\}^{2}\right]>\mathrm{E}\left[\left\{y_{t}-\varphi\left(\mathbf{x}_{t} ; \boldsymbol{\theta}_{0}\right)\right\}^{2}\right] & \text { if } \varphi\left(\mathbf{x}_{t} ; \boldsymbol{\theta}\right) \neq \varphi\left(\mathbf{x}_{t} ; \boldsymbol{\theta}_{0}\right)  \tag{7.2.9}\\ \mathrm{E}\left[\left\{y_{t}-\varphi\left(\mathbf{x}_{t} ; \boldsymbol{\theta}\right)\right\}^{2}\right]=\mathrm{E}\left[\left\{y_{t}-\varphi\left(\mathbf{x}_{t} ; \boldsymbol{\theta}_{0}\right)\right\}^{2}\right] & \text { if } \varphi\left(\mathbf{x}_{t} ; \boldsymbol{\theta}\right)=\varphi\left(\mathbf{x}_{t} ; \boldsymbol{\theta}_{0}\right)\end{cases}
\]

Therefore, the identification condition - that $\mathrm{E}\left[m\left(\mathbf{x}_{t} ; \boldsymbol{\theta}\right)\right]=-\mathrm{E}\left[\left\{y_{t}-\varphi\left(\mathbf{x}_{t} ; \boldsymbol{\theta}\right)\right\}^{2}\right]$ be maximized (i.e., $\mathrm{E}\left[\left\{y_{t}-\varphi\left(\mathbf{x}_{t} ; \boldsymbol{\theta}\right)\right\}^{2}\right]$ be minimized) uniquely at $\boldsymbol{\theta}_{0}$-holds if

$$
\text { conditional mean identification: } \varphi\left(\mathbf{x}_{t} ; \boldsymbol{\theta}\right) \neq \varphi\left(\mathbf{x}_{t} ; \boldsymbol{\theta}_{0}\right) \text { for all } \boldsymbol{\theta} \neq \boldsymbol{\theta}_{0} \text {. (7.2.10) }
$$

ML\\
For ML, the role played by (7.2.8) for NLS is played by the Kullback-Leibler information inequality. Here we examine identification for conditional ML (doing the same for unconditional or joint ML is more straightforward). The hypothetical conditional density $f\left(y_{t} \mid \mathbf{x}_{t} ; \boldsymbol{\theta}\right)$, being a function of ( $y_{t}, \mathbf{x}_{t}$ ), is a random variable for any $\boldsymbol{\theta}$. The expected value of the random variable $\log f\left(y_{t} \mid \mathbf{x}_{t} ; \boldsymbol{\theta}\right)$ can be written as ${ }^{8}$


\begin{equation*}
\mathrm{E}\left[\log f\left(y_{t} \mid \mathbf{x}_{t} ; \boldsymbol{\theta}\right)\right]=\int \log f\left(y_{t} \mid \mathbf{x}_{t} ; \boldsymbol{\theta}\right) f\left(y_{t}, \mathbf{x}_{t} ; \boldsymbol{\theta}_{0}, \boldsymbol{\psi}_{0}\right) \mathrm{d} y_{t} \mathrm{~d} \mathbf{x}_{t} \tag{7.2.11}
\end{equation*}


Note well that the expected value is with respect to the true joint density $f\left(y_{t}, \mathbf{x}_{t}\right.$; $\boldsymbol{\theta}_{0}, \boldsymbol{\psi}_{0}$ ). The Kullback-Leibler information inequality about density functions (adapted to conditional densities) states that

\[
\begin{cases}\mathrm{E}\left[\log f\left(y_{t} \mid \mathbf{x}_{t} ; \boldsymbol{\theta}\right)\right]<\mathrm{E}\left[\log f\left(y_{t} \mid \mathbf{x}_{t} ; \boldsymbol{\theta}_{0}\right)\right] & \text { if } f\left(y_{t} \mid \mathbf{x}_{t} ; \boldsymbol{\theta}\right) \neq f\left(y_{t} \mid \mathbf{x}_{t} ; \boldsymbol{\theta}_{0}\right)  \tag{7.2.12}\\ \mathrm{E}\left[\log f\left(y_{t} \mid \mathbf{x}_{t} ; \boldsymbol{\theta}\right)\right]=\mathrm{E}\left[\log f\left(y_{t} \mid \mathbf{x}_{t} ; \boldsymbol{\theta}_{0}\right)\right] & \text { if } f\left(y_{t} \mid \mathbf{x}_{t} ; \boldsymbol{\theta}\right)=f\left(y_{t} \mid \mathbf{x}_{t} ; \boldsymbol{\theta}_{0}\right)\end{cases}
\]

(provided, of course, that $\mathrm{E}\left[\log f\left(y_{t} \mid \mathbf{x}_{t} ; \boldsymbol{\theta}\right)\right]$ exists and is finite). (See Analytical Exercise 1 for derivation.) It immediately follows that the identification condition (that $\mathrm{E}\left[\log f\left(y_{t} \mid \mathbf{x}_{t} ; \boldsymbol{\theta}\right)\right]$ be maximized uniquely at $\boldsymbol{\theta}_{0}$ ) is satisfied if


\begin{align*}
& \text { conditional density identification: } \begin{aligned}
& f\left(y_{t} \mid \mathbf{x}_{t} ; \boldsymbol{\theta}\right) \neq f\left(y_{t} \mid \mathbf{x}_{t} ; \boldsymbol{\theta}_{0}\right) \\
& \text { for all } \boldsymbol{\theta} \neq \boldsymbol{\theta}_{0} .
\end{aligned}
\end{align*}


\footnotetext{${ }^{8}$ If $f\left(y_{t} \mid \mathbf{x}_{t} ; \boldsymbol{\theta}\right)=0$, then its log cannot be defined. If this bothers you, assume $f\left(y_{t} \mid \mathbf{x}_{t} ; \boldsymbol{\theta}\right)$ is positive for all $y_{t}, \mathbf{x}_{t}$, and $\theta$. See White (1994, p. 9) for an argument that avoids such a simplifying assumption.
}We can turn the two consistency theorems for M-estimators, Propositions 7.3 and 7.4, into ones for conditional ML, with $\mathbf{w}_{t}=\left(y_{t}, \mathbf{x}_{t}^{\prime}\right)^{\prime}$ and $m\left(\mathbf{w}_{t} ; \boldsymbol{\theta}\right)=\log f$. $\left(y_{t} \mid \mathbf{x}_{t} ; \boldsymbol{\theta}\right)$. Replacing the identification condition for M -estimators by conditional density identification, we obtain the following two self-contained statements about consistency of ML.

Proposition 7.5 (Consistency of conditional ML with compact parameter space): Let $\left\{y_{t}, \mathbf{x}_{t}\right\}$ be ergodic stationary with conditional density $f\left(y_{t} \mid \mathbf{x}_{t} ; \boldsymbol{\theta}_{0}\right)$ and let $\hat{\boldsymbol{\theta}}$ be the conditional ML estimator, which maximizes the average log conditional likelihood (derived under the assumption that $\left\{y_{t}, \mathbf{x}_{t}\right\}$ is i.i.d.):

$$
\hat{\boldsymbol{\theta}}=\underset{\boldsymbol{\theta} \in \Theta}{\operatorname{argmax}} \frac{1}{n} \sum_{t=1}^{n} \log f\left(y_{t} \mid \mathbf{x}_{t} ; \boldsymbol{\theta}\right)
$$

Suppose the model is correctly specified so that $\boldsymbol{\theta}_{0}$ is in $\Theta$. Suppose that (i) the parameter space $\Theta$ is a compact subset of $\mathbb{R}^{p}$, (ii) $f\left(y_{t} \mid \mathbf{x}_{t} ; \boldsymbol{\theta}\right)$ is continuous in $\boldsymbol{\theta}$ for all ( $y_{t}, \mathbf{x}_{t}$ ), and (iii) $f\left(y_{t} \mid \mathbf{x}_{t} ; \boldsymbol{\theta}\right)$ is measurable in ( $y_{t}, \mathbf{x}_{t}$ ) for all $\boldsymbol{\theta} \in \Theta$ (so by Lemma $7.1 \hat{\boldsymbol{\theta}}$ is a well-defined random variable). Suppose, further, that\\
(a) (identification) $\operatorname{Prob}\left[f\left(y_{t} \mid \mathbf{x}_{t} ; \boldsymbol{\theta}\right) \neq f\left(y_{t} \mid \mathbf{x}_{t} ; \boldsymbol{\theta}_{0}\right)\right]>0$ for all $\boldsymbol{\theta} \neq \boldsymbol{\theta}_{0}$ in $\Theta$,\\
(b) (dominance) $\mathrm{E}\left[\sup _{\theta \in \Theta}\left|\log f\left(y_{t} \mid \mathbf{x}_{t} ; \boldsymbol{\theta}\right)\right|\right]<\infty$ (note: the expectation is over $y_{t}$ and $\mathbf{x}_{t}$ ).

Then $\hat{\boldsymbol{\theta}} \rightarrow_{\mathrm{p}} \boldsymbol{\theta}_{0}$.

Proposition 7.6 (Consistency of conditional ML without compactness): Let $\left\{y_{t}, \mathbf{x}_{t}\right\}$ be ergodic stationary with conditional density $f\left(y_{t} \mid \mathbf{x}_{t} ; \boldsymbol{\theta}_{0}\right)$ and let $\hat{\boldsymbol{\theta}}$ be the conditional ML estimator, which maximizes the average log conditional likelihood (derived under the assumption that $\left\{y_{t}, \mathbf{x}_{t}\right\}$ is i.i.d.):

$$
\hat{\boldsymbol{\theta}}=\underset{\boldsymbol{\theta} \in \Theta}{\operatorname{argmax}} \frac{1}{n} \sum_{t=1}^{n} \log f\left(y_{t} \mid \mathbf{x}_{t} ; \boldsymbol{\theta}\right)
$$

Suppose the model is correctly specified so that $\boldsymbol{\theta}_{0}$ is in $\Theta$. Suppose that (i) the true parameter vector $\boldsymbol{\theta}_{0}$ is an element of the interior of a convex parameter space $\Theta\left(\subset \mathbb{R}^{p}\right)$, (ii) $\log f\left(y_{t} \mid \mathbf{x}_{t} ; \boldsymbol{\theta}\right)$ is concave in $\boldsymbol{\theta}$ for all $\left(y_{t}, \mathbf{x}_{t}\right)$, and (iii) $\log f\left(y_{t} \mid\right. \mathbf{x}_{t} ; \boldsymbol{\theta}$ ) is measurable in ( $y_{t}, \mathbf{x}_{t}$ ) for all $\boldsymbol{\theta}$ in $\Theta$. (For sufficiently large $n, \hat{\boldsymbol{\theta}}$ is well defined; see a few lines below.) Suppose, further, that\\
(a) (identification) $\operatorname{Prob}\left[f\left(y_{t} \mid \mathbf{x}_{t} ; \boldsymbol{\theta}\right) \neq f\left(y_{t} \mid \mathbf{x}_{t} ; \boldsymbol{\theta}_{0}\right)\right]>0$ for all $\boldsymbol{\theta} \neq \boldsymbol{\theta}_{0}$ in $\Theta$,\\
(b) $\mathrm{E}\left[\left|\log f\left(y_{t} \mid \mathbf{x}_{t} ; \boldsymbol{\theta}\right)\right|\right]<\infty$ (i.e., $\mathrm{E}\left[\log f\left(y_{t} \mid \mathbf{x}_{t} ; \boldsymbol{\theta}\right)\right]$ exists and is finite) for all $\boldsymbol{\theta} \in \Theta$ (note: the expectation is over $y_{t}$ and $\mathbf{x}_{t}$ ).

Then, as $n \rightarrow \infty, \hat{\boldsymbol{\theta}}$ exists with probability approaching 1 and $\hat{\boldsymbol{\theta}} \rightarrow_{\mathrm{p}} \boldsymbol{\theta}_{0}$.

There are two remarks worth making.

\begin{itemize}
  \item The Kullback-Leibler information inequality (7.2.12) holds for unconditional densities as well (the proof is easier for unconditional densities). Therefore, Propositions 7.5 and 7.6 hold for unconditional ML; just replace $f\left(y_{t} \mid \mathbf{x}_{t} ; \boldsymbol{\theta}\right)$ by $f\left(\mathbf{w}_{t} ; \boldsymbol{\theta}\right)$.
  \item The noteworthy feature of these two ML consistency theorems is that, despite the fact that the log likelihood is derived under the i.i.d. assumption for $\left\{y_{t}, \mathbf{x}_{t}\right\}$, consistency is assured even if $\left\{y_{t}, \mathbf{x}_{t}\right\}$ is not i.i.d. but merely ergodic stationary. This is because the ML estimator is an M-estimator. An ML estimator that maximizes a likelihood function different from the model's likelihood function is called a quasi-ML estimator or QML estimator. Put differently, a QML estimator maximizes the likelihood function of a misspecified model. Therefore, the "maximum likelihood" estimators in Propositions 7.5 and 7.6 are QML estimators if $\left\{y_{t}, \mathbf{x}_{t}\right\}$ is ergodic stationary but not i.i.d. The propositions say that those particular QML estimators are consistent.
\end{itemize}

The identification condition in conditional ML can be stated in even more primitive terms for specific examples.

Example 7.8 (Identification in linear regression model): For the linear regression model with normal errors, the log conditional likelihood for observation $t$ is

$$
\log f\left(y_{t} \mid \mathbf{x}_{t} ; \boldsymbol{\theta}\right)=-\frac{1}{2} \log (2 \pi)-\frac{1}{2} \log \left(\sigma^{2}\right)-\frac{1}{2 \sigma^{2}}\left(y_{t}-\mathbf{x}_{t}^{\prime} \boldsymbol{\beta}\right)^{2} .
$$

Since the log function is strictly monotonic, conditional density identification is equivalent to the condition that $\log f\left(y_{t} \mid \mathbf{x}_{t} ; \boldsymbol{\theta}\right) \neq \log f\left(y_{t} \mid \mathbf{x}_{t} ; \boldsymbol{\theta}_{0}\right)$ with positive probability if $\boldsymbol{\theta} \neq \boldsymbol{\theta}_{0}$. This condition is satisfied if $\mathrm{E}\left(\mathbf{x}_{t} \mathbf{x}_{t}^{\prime}\right)$ is nonsingular (and hence positive definite). To see why, let $\boldsymbol{\theta}=\left(\boldsymbol{\beta}^{\prime}, \sigma^{2}\right)^{\prime}$ and $\boldsymbol{\theta}_{0}=\left(\boldsymbol{\beta}_{0}^{\prime}, \sigma_{0}^{2}\right)^{\prime}$. Clearly, $\log f\left(y_{t} \mid \mathbf{x}_{t} ; \boldsymbol{\theta}\right) \neq \log f\left(y_{t} \mid \mathbf{x}_{t} ; \boldsymbol{\theta}_{0}\right)$ with positive probability if $\sigma^{2} \neq \sigma_{0}^{2}$. So consider the case where $\sigma^{2}=\sigma_{0}^{2}$ but $\boldsymbol{\beta} \neq \boldsymbol{\beta}_{0}$. Then

$$
\mathrm{E}\left[\left\{\mathbf{x}_{t}^{\prime} \boldsymbol{\beta}-\mathbf{x}_{t}^{\prime} \boldsymbol{\beta}_{0}\right\}^{2}\right]=\mathrm{E}\left[\left\{\mathbf{x}_{t}^{\prime}\left(\boldsymbol{\beta}-\boldsymbol{\beta}_{0}\right)\right\}^{2}\right]=\left(\boldsymbol{\beta}-\boldsymbol{\beta}_{0}\right)^{\prime} \mathrm{E}\left(\mathbf{x}_{t} \mathbf{x}_{t}^{\prime}\right)\left(\boldsymbol{\beta}-\boldsymbol{\beta}_{0}\right)>0 .
$$

Hence, $\mathbf{x}_{t}^{\prime} \boldsymbol{\beta} \neq \mathbf{x}_{t}^{\prime} \boldsymbol{\beta}_{0}$ with positive probability; if $\mathbf{x}_{t}^{\prime} \boldsymbol{\beta}=\mathbf{x}_{t}^{\prime} \boldsymbol{\beta}_{0}$ with probability 1 (i.e., if $\mathbf{x}_{t}^{\prime} \boldsymbol{\beta} \neq \mathbf{x}_{t}^{\prime} \boldsymbol{\beta}_{0}$ with probability zero), then $\mathrm{E}\left[\left\{\mathbf{x}_{t}^{\prime} \boldsymbol{\beta}-\mathbf{x}_{t}^{\prime} \boldsymbol{\beta}_{0}\right\}^{2}\right]$ will be zero. Thus, $y_{t}-\mathbf{x}_{t}^{\prime} \boldsymbol{\beta} \neq y_{t}-\mathbf{x}_{t}^{\prime} \boldsymbol{\beta}_{0}$ with positive probability, implying that the identification condition is satisfied. The same condition about $\mathrm{E}\left(\mathbf{x}_{t} \mathbf{x}_{t}^{\prime}\right)$ also implies condition (b) in Proposition 7.6 because

$$
\begin{aligned}
\mathrm{E}\left[\left(y_{t}-\mathbf{x}_{t}^{\prime} \boldsymbol{\beta}\right)^{2}\right]= & \mathrm{E}\left[\left\{\varepsilon_{t}+\mathbf{x}_{t}^{\prime}\left(\boldsymbol{\beta}_{0}-\boldsymbol{\beta}\right)\right\}^{2}\right] \quad\left(\text { since } y_{t}=\mathbf{x}_{t}^{\prime} \boldsymbol{\beta}_{0}+\varepsilon_{t}\right) \\
= & \mathrm{E}\left(\varepsilon_{t}^{2}\right)+\left(\boldsymbol{\beta}_{0}-\boldsymbol{\beta}\right)^{\prime} \mathrm{E}\left(\mathbf{x}_{t} \mathbf{x}_{t}^{\prime}\right)\left(\boldsymbol{\beta}_{0}-\boldsymbol{\beta}\right) \\
& \left(\text { since } \mathrm{E}\left(\varepsilon_{t} \cdot \mathbf{x}_{t}\right)=\mathbf{0}\right) .
\end{aligned}
$$

The last term is finite if $\mathrm{E}\left(\mathbf{x}_{t} \mathbf{x}_{t}^{\prime}\right)$ is. These results, together with the fact noted earlier that the log likelihood is concave after reparameterization, imply that the conditional ML estimator of the linear regression model with normal errors (derived under the i.i.d. assumption for $\left\{y_{t}, \mathbf{x}_{t}\right\}$ ) is consistent if $\left\{y_{t}, \mathbf{x}_{t}\right\}$ is ergodic stationary and $\mathrm{E}\left(\mathbf{x}_{t} \mathbf{x}_{t}^{\prime}\right)$ is nonsingular.

Example 7.9 (Identification in probit): The same conclusion just derived in the previous example for the linear model also holds for the probit model, whose conditional likelihood for observation $t$ is

$$
f\left(y_{t} \mid \mathbf{x}_{t} ; \boldsymbol{\theta}\right)=\Phi\left(\mathbf{x}_{t}^{\prime} \boldsymbol{\theta}\right)^{y_{t}} \Phi\left(-\mathbf{x}_{t}^{\prime} \boldsymbol{\theta}\right)^{1-y_{t}} .
$$

The same argument used in Example 7.8 implies that, under the nonsingularity of $\mathrm{E}\left(\mathbf{x}_{t} \mathbf{x}_{t}^{\prime}\right), \mathbf{x}_{t}^{\prime} \boldsymbol{\theta} \neq \mathbf{x}_{t}^{\prime} \boldsymbol{\theta}_{0}$ with positive probability if $\boldsymbol{\theta} \neq \boldsymbol{\theta}_{0}$. But since $\Phi(v)$ is strictly monotonic, $\mathbf{x}_{t}^{\prime} \boldsymbol{\theta} \neq \mathbf{x}_{t}^{\prime} \boldsymbol{\theta}_{0}$ with positive probability implies $\Phi\left(\mathbf{x}_{t}^{\prime} \boldsymbol{\theta}\right) \neq \Phi\left(\mathbf{x}_{t}^{\prime} \boldsymbol{\theta}_{0}\right)$ and $\Phi\left(-\mathbf{x}_{t}^{\prime} \boldsymbol{\theta}\right) \neq \Phi\left(-\mathbf{x}_{t}^{\prime} \boldsymbol{\theta}_{0}\right)$ with positive probability. Hence, the identification condition (a) in Proposition 7.6 is satisfied if $\mathbf{E}\left(\mathbf{x}_{t} \mathbf{x}_{t}^{\prime}\right)$ is nonsingular.

The nonsingularity of $\mathrm{E}\left(\mathbf{x}_{t} \mathbf{x}_{t}^{\prime}\right)$ also implies that condition (b) of Proposition 7.6 is satisfied for probit. It is easy to verify that


\begin{equation*}
|\log \Phi(v)| \leq|\log \Phi(0)|+|v|+|v|^{2} \text { for all } v . \tag{7.2.14}
\end{equation*}


Combining this bound for $|\log \Phi(v)|$ and the fact that $y_{t}$ and $1-y_{t}$ are less than or equal to 1 in absolute value, it is easy to show that $\mathrm{E}\left[\mid \log f\left(y_{t} \mid\right.\right. \left.\left.\mathbf{x}_{t} ; \boldsymbol{\theta}\right) \mid\right]<\infty$ if $\mathrm{E}\left(\mathbf{x}_{t}, \mathbf{x}_{t}^{\prime}\right)$ exists and is finite (a condition implied by the non-\\
singularity of $\left.\mathrm{E}\left(\mathbf{x}_{t} \mathbf{x}_{t}^{\prime}\right)\right) .{ }^{9}$ We therefore conclude that the probit ML estimator (whose log likelihood is derived under the i.i.d. assumption for $\left\{y_{t}, \mathbf{x}_{t}\right\}$ ) is consistent if $\left\{y_{t}, \mathbf{x}_{t}\right\}$ is ergodic stationary and $\mathrm{E}\left(\mathbf{x}_{t} \mathbf{x}_{t}^{\prime}\right)$ is nonsingular.

\section*{Consistency of GMM}
We now turn to GMM and specialize Proposition 7.1 to GMM by verifying the conditions of the proposition under a set of appropriate sufficient conditions. The GMM objective function is given by (7.1.3). Clearly, the continuity of $Q_{n}(\boldsymbol{\theta})$ in $\boldsymbol{\theta}$ is satisfied if $\mathbf{g}\left(\mathbf{w}_{t} ; \boldsymbol{\theta}\right)$ is continuous in $\boldsymbol{\theta}$ for all $\mathbf{w}_{t}$. If $\left\{\mathbf{w}_{t}\right\}$ is ergodic stationary, then $\mathbf{g}_{n}(\boldsymbol{\theta})\left(\equiv \frac{1}{n} \sum_{t=1}^{n} \mathbf{g}\left(\mathbf{w}_{t} ; \boldsymbol{\theta}\right)\right) \rightarrow_{\mathrm{p}} \mathrm{E}\left[\mathbf{g}\left(\mathbf{w}_{t} ; \boldsymbol{\theta}\right)\right]$. So the limit function $Q_{0}(\boldsymbol{\theta})$ is


\begin{equation*}
Q_{0}(\boldsymbol{\theta})=-\frac{1}{2} \mathrm{E}\left[\mathbf{g}\left(\mathbf{w}_{t} ; \boldsymbol{\theta}\right)\right]^{\prime} \mathrm{W} \mathrm{E}\left[\mathbf{g}\left(\mathbf{w}_{t} ; \boldsymbol{\theta}\right)\right] \tag{7.2.15}
\end{equation*}


This function is nonpositive if $\mathbf{W}(\equiv \operatorname{plim} \widehat{\mathbf{W}})$ is positive definite. It has a maximum of zero at $\boldsymbol{\theta}_{0}$ because $\mathbf{E}\left[\mathbf{g}\left(\mathbf{w}_{t} ; \boldsymbol{\theta}_{0}\right)\right]=\mathbf{0}$ by the orthogonality conditions. Thus, the identification condition (that $Q_{0}(\boldsymbol{\theta})$ be uniquely maximized at $\boldsymbol{\theta}_{0}$ ) in Proposition 7.1 is satisfied if $\mathrm{E}\left[\mathbf{g}\left(\mathbf{w}_{t} ; \boldsymbol{\theta}\right)\right] \neq \mathbf{0}$ for $\boldsymbol{\theta} \neq \boldsymbol{\theta}_{0}$. Regarding the uniform convergence of $Q_{n}(\boldsymbol{\theta})$ to $Q_{0}(\boldsymbol{\theta})$, it is not hard to show that the condition is satisfied if $\mathbf{g}_{n}(\cdot)$ converges uniformly to $\mathbf{E}\left[\mathbf{g}\left(\mathbf{w}_{t} ; \cdot\right)\right]$ (see Newey and McFadden, 1994, p. 2132). The multivariate version of the Uniform Law of Large Numbers stated right below Lemma 7.2, in turn, provides a sufficient condition for the uniform convergence. So we have proved

Proposition 7.7 (Consistency of GMM with compact parameter space): Let $\left\{\mathbf{w}_{t}\right\}$ be ergodic stationary and let $\hat{\boldsymbol{\theta}}$ be the GMM estimator defined by

$$
\hat{\boldsymbol{\theta}}=\underset{\boldsymbol{\theta} \in \Theta}{\operatorname{argmin}}\left[\frac{1}{n} \sum_{t=1}^{n} \mathbf{g}\left(\mathbf{w}_{t} ; \boldsymbol{\theta}\right)\right]^{\prime} \widehat{\mathbf{W}}\left[\frac{1}{n} \sum_{t=1}^{n} \mathbf{g}\left(\mathbf{w}_{t} ; \boldsymbol{\theta}\right)\right],
$$

where the symmetric matrix $\widehat{\mathbf{W}}$ is assumed to converge in probability to some\\
${ }^{9}$ Only for interested readers, here is a proof:

$$
\begin{aligned}
\left|\log f\left(y_{t} \mid \mathbf{x}_{t} ; \boldsymbol{\theta}\right)\right| & \leq\left|y_{t}\right|\left|\log \Phi\left(\mathbf{x}_{t}^{\prime} \boldsymbol{\theta}\right)\right|+\left|1-y_{t}\right|\left|\log \Phi\left(-\mathbf{x}_{t}^{\prime} \boldsymbol{\theta}\right)\right| \\
& \leq\left|\log \Phi\left(\mathbf{x}_{t}^{\prime} \boldsymbol{\theta}\right)\right|+\left|\log \Phi\left(-\mathbf{x}_{t}^{\prime} \boldsymbol{\theta}\right)\right| \quad\left(\text { since }\left|y_{t}\right| \leq 1 \text { and }\left|1-y_{t}\right| \leq 1\right) \\
& \leq 2 \times\left[|\log \Phi(0)|+\left|\mathbf{x}_{t}^{\prime} \boldsymbol{\theta}\right|+\left|\mathbf{x}_{t}^{\prime} \boldsymbol{\theta}\right|^{2}\right] \quad(\text { by }(7.2 .14)) \\
& \leq 2 \times\left[|\log \Phi(0)|+\left\|\mathbf{x}_{t}\right\| \cdot\|\boldsymbol{\theta}\|+\left\|\mathbf{x}_{t}\right\|^{2} \cdot\|\boldsymbol{\theta}\|^{2}\right] .
\end{aligned}
$$

The last inequality is due to the Cauchy-Schwartz inequality that $\left|\mathbf{x}^{\prime} \mathbf{y}\right| \leq\|\mathbf{x}\|\|\mathbf{y}\|$ for any two conformable vectors $\mathbf{x}$ and $\mathbf{y}$. Also, $\mathrm{E}\left(\left\|\mathbf{x}_{t}\right\|^{2}\right)$, which is the sum of $\mathrm{E}\left(x_{i t}^{2}\right)$ over $i$, is finite because the nonsingularity of $\mathrm{E}\left(\mathbf{x}_{t} \mathbf{x}_{t}^{\prime}\right)$ implies $\mathrm{E}\left(x_{i t}^{2}\right)<\infty$ for all $i$. If $\mathrm{E}\left(\left\|\mathbf{x}_{t}\right\|^{2}\right)<\infty$, then $\mathrm{E}\left(\left\|\mathbf{x}_{t}\right\|\right)<\infty$.\\
symmetric positive definite matrix W. Suppose that the model is correctly specified in that $\mathrm{E}\left[\mathbf{g}\left(\mathbf{w}_{t} ; \boldsymbol{\theta}_{0}\right)\right]=\mathbf{0}$ (i.e., the orthogonality conditions hold) for $\boldsymbol{\theta}_{0} \in \Theta$. Suppose that (i) the parameter space $\Theta$ is a compact subset of $\mathbb{R}^{p}$, (ii) $\mathbf{g}\left(\mathbf{w}_{t} ; \boldsymbol{\theta}\right)$ is continuous in $\boldsymbol{\theta}$ for all $\mathbf{w}_{t}$, and (iii) $\mathbf{g}\left(\mathbf{w}_{t} ; \boldsymbol{\theta}\right)$ is measurable in $\mathbf{w}_{t}$ for all $\boldsymbol{\theta}$ in $\Theta$ (so $\hat{\boldsymbol{\theta}}$ is a well-defined random variable by Lemma 7.1). Suppose, further, that\\
(a) (identification) $\mathrm{E}\left[\mathbf{g}\left(\mathbf{w}_{t} ; \boldsymbol{\theta}\right)\right] \neq \mathbf{0}$ for all $\boldsymbol{\theta} \neq \boldsymbol{\theta}_{0}$ in $\Theta$,\\
(b) (dominance) $\mathrm{E}\left[\sup _{\theta \in \Theta}\left\|\mathbf{g}\left(\mathbf{w}_{t} ; \boldsymbol{\theta}\right)\right\|\right]<\infty$.

Then $\hat{\boldsymbol{\theta}} \rightarrow_{\mathrm{p}} \boldsymbol{\theta}_{0}$.

If $Q_{n}(\boldsymbol{\theta})$ is concave, then the compactness of $\Theta$, the continuity of $\mathbf{g}\left(\mathbf{w}_{t} ; \boldsymbol{\theta}\right)$, and the dominance condition can be replaced by the condition that $\mathrm{E}\left[\mathbf{g}\left(\mathbf{w}_{t} ; \boldsymbol{\theta}\right)\right]$ exist and be finite for all $\theta$. We will not state this result as a proposition because there are very few applications in which $Q_{n}(\boldsymbol{\theta})$ is concave. For example, the GMM objective function for the Euler equation model in Example 7.5 is not concave. Just about the only well-known case is one in which $\mathbf{g}\left(\mathbf{w}_{t} ; \boldsymbol{\theta}\right)$ is linear in $\boldsymbol{\theta}$, but the linear case has been treated rather thoroughly in Chapter 3. When $\mathbf{g}\left(\mathbf{w}_{t} ; \boldsymbol{\theta}\right)$ is nonlinear in $\boldsymbol{\theta}$, specifying primitive conditions for (a) (identification) and (b) (dominance) in Proposition 7.7 is generally quite difficult. So in most applications those conditions are simply assumed.

\section*{QUESTIONS FOR REVIEW}
\begin{enumerate}
  \item (Necessity of density identification) We have shown for conditional ML that the conditional density identification (7.2.13) is sufficient for condition (a) (that $\mathrm{E}\left[\log f\left(y_{t} \mid \mathbf{x}_{t} ; \boldsymbol{\theta}\right)\right]$ be maximized uniquely at $\left.\boldsymbol{\theta}_{0}\right)$. That it is also a necessary condition is clear from the Kullback-Leibler information inequality. Without using this inequality, show the necessity of conditional mean identification. Hint: Show that, if (7.2.13) were false, then (a) would not hold. Suppose that there exists a $\boldsymbol{\theta}_{1} \neq \boldsymbol{\theta}_{0}$ such that $\boldsymbol{\theta}_{1} \in \Theta$ and $f\left(y_{t} \mid \mathbf{x}_{t} ; \boldsymbol{\theta}_{1}\right)=f\left(y_{t} \mid \mathbf{x}_{t} ; \boldsymbol{\theta}_{0}\right)$. Then $\mathrm{E}\left[\log f\left(y_{t} \mid \mathbf{x}_{t} ; \boldsymbol{\theta}_{1}\right)\right]=\mathrm{E}\left[\log f\left(y_{t} \mid \mathbf{x}_{t} ; \boldsymbol{\theta}_{0}\right)\right]$.
  \item (Necessity of conditional mean identification) Without using (7.2.8), show that conditional mean identification is necessary as well as sufficient for identification in NLS.
  \item (Identification in NLS) Suppose that the $\varphi$ function in NLS is linear: $\varphi\left(\mathbf{x}_{t} ; \boldsymbol{\theta}\right)=\mathbf{x}_{t}^{\prime} \boldsymbol{\theta}$. Verify that identification is satisfied if $\mathrm{E}\left(\mathbf{x}_{t} \mathbf{x}_{t}^{\prime}\right)$ is nonsingular.
  \item (Identification in NLS estimation of probit) In the probit model, $\mathrm{E}\left(y_{t} \mid \mathbf{x}_{t}\right)= \Phi\left(\mathbf{x}_{t}^{\prime} \boldsymbol{\theta}_{0}\right)$. Consider estimating $\boldsymbol{\theta}_{0}$ by least squares. Verify that conditional mean identification is satisfied if $\mathrm{E}\left(\mathbf{x}_{t} \mathbf{x}_{t}^{\prime}\right)$ is nonsingular.
  \item (Quasi-ML) Suppose in Proposition 7.5 that $\left\{y_{t}, \mathbf{x}_{t}\right\}$ is i.i.d., so the likelihood function itself is not misspecified. But suppose that the true parameter vector $\boldsymbol{\theta}_{0}$ is not included in the assumed parameter space $\Theta$. So in this sense the estimator $\hat{\boldsymbol{\theta}}$ is a quasi-ML estimator. Drop condition (a), and suppose instead that there is a unique $\boldsymbol{\theta}^{*}$ in $\Theta$ that maximizes $\mathrm{E}\left[\log f\left(y_{t} \mid \mathbf{x}_{t} ; \boldsymbol{\theta}\right)\right]$. Show that the QML estimator $\hat{\boldsymbol{\theta}}$ is consistent for $\boldsymbol{\theta}^{*}$. Hint: Replace the $\boldsymbol{\theta}_{0}$ in Proposition 7.3 by $\boldsymbol{\theta}^{*}$. In Propositions 7.1-7.4, $\boldsymbol{\theta}_{0}$ does not need to be the "true" parameter vector.
  \item (Identification in linear GMM) Consider the linear GMM model of Chapter 3. Verify that the rank condition for identification derived in Chapter 3 is equivalent to the identification condition in Proposition 7.7.
  \item (Concavity of the objective function in linear GMM) Verify that the GMM objective function in linear GMM is concave in $\boldsymbol{\theta}$. Hint: The Hessian (the matrix of second derivatives) of $Q_{n}(\boldsymbol{\theta})$ is $-\mathbf{S}_{\mathbf{x z}}^{\prime} \widehat{\mathbf{W}} \mathbf{S}_{\mathbf{x z}}$ where $\mathbf{S}_{\mathbf{x z}} \equiv \frac{1}{n} \sum_{t=1}^{n} \mathbf{x}_{t} \mathbf{z}_{t}^{\prime}$. A twice continuously differentiable function is concave if and only if its Hessian is everywhere negative semidefinite.
  \item (GMM identification with singular $\mathbf{W}$ ) For $\mathbf{W}$, if we require it to be merely positive semidefinite, not necessarily positive definite, what would be the identification condition in Proposition 7.7? Hint: Suppose W is positive semidefinite. Show that if $\mathrm{E}\left[\mathbf{g}\left(\mathbf{w}_{t} ; \boldsymbol{\theta}_{0}\right)\right]=\mathbf{0}$ and $\mathbf{W} \mathrm{E}\left[\mathbf{g}\left(\mathbf{w}_{t} ; \boldsymbol{\theta}\right)\right] \neq \mathbf{0}$ for $\boldsymbol{\theta} \neq \boldsymbol{\theta}_{0}$, then $Q_{0}(\boldsymbol{\theta})=-\frac{1}{2} \mathrm{E}\left[\mathbf{g}\left(\mathbf{w}_{t} ; \boldsymbol{\theta}\right)\right]^{\prime} \mathbf{W} \mathrm{E}\left[\mathbf{g}\left(\mathbf{w}_{t} ; \boldsymbol{\theta}\right)\right]$ has a unique maximum at $\boldsymbol{\theta}_{0}$.
\end{enumerate}

\subsection*{7.3 Asymptotic Normality}
The basic tool in deriving the asymptotic distribution of an extremum estimator is the Mean Value Theorem from calculus. How it is used depends on whether the estimator is an M-estimator or a GMM estimator. Therefore, the proof of asymptotic normality will be done separately for M-estimators (including ML as a special case) and GMM estimators. This is followed by a brief remark on the relative efficiency of GMM and ML. By way of summing up, we will point out at the end of\\
the section that the Taylor expansion for the sampling error has a structure shared by both M-estimators and GMM.

\section*{Asymptotic Normality of M-Estimators}
Reproducing (7.1.2), the objective function for M-estimators is


\begin{equation*}
Q_{n}(\boldsymbol{\theta})=\frac{1}{n} \sum_{t=1}^{n} m\left(\mathbf{w}_{t} ; \boldsymbol{\theta}\right) \tag{7.3.1}
\end{equation*}


It will be convenient to give symbols to the gradient (vector of first derivatives) and the Hessian (matrix of second derivatives) of the $m$ function as


\begin{align*}
\underset{(p \times 1)}{\mathbf{s}\left(\mathbf{w}_{t} ; \boldsymbol{\theta}\right)} & =\frac{\partial m\left(\mathbf{w}_{t} ; \boldsymbol{\theta}\right)}{\partial \boldsymbol{\theta}}  \tag{7.3.2}\\
\underset{(p \times p)}{\mathbf{H}\left(\mathbf{w}_{t} ; \boldsymbol{\theta}\right)} & =\frac{\partial \mathbf{s}\left(\mathbf{w}_{t} ; \boldsymbol{\theta}\right)}{\partial \boldsymbol{\theta}^{\prime}}=\frac{\partial^{2} m\left(\mathbf{w}_{t} ; \boldsymbol{\theta}\right)}{\partial \boldsymbol{\theta} \partial \boldsymbol{\theta}^{\prime}} \tag{7.3.3}
\end{align*}


In analogy to ML, $\mathbf{s}\left(\mathbf{w}_{t} ; \boldsymbol{\theta}\right)$ will be referred to as the score vector for observation $\boldsymbol{t}$. This $\mathbf{s}\left(\mathbf{w}_{t} ; \boldsymbol{\theta}\right)$ should not be confused with the score in the usual sense, which is the gradient of the objective function $Q_{n}(\boldsymbol{\theta})$. The score in the latter sense will be denoted $\mathbf{s}_{n}(\boldsymbol{\theta})$ later in this chapter. The same applies to the Hessian: $\mathbf{H}\left(\mathbf{w}_{t} ; \boldsymbol{\theta}\right)$ will be referred to as the Hessian for observation $t$, and the Hessian of $Q_{n}(\boldsymbol{\theta})$ will be denoted $\mathbf{H}_{n}(\boldsymbol{\theta})$.

The goal of this subsection is the asymptotic normality of $\hat{\boldsymbol{\theta}}$ described in (7.3.10) below. In the process of deriving it, we will make a number of assumptions, which will be collected in Proposition 7.8 below. Assume that $m\left(\mathbf{w}_{t} ; \boldsymbol{\theta}\right)$ is differentiable in $\boldsymbol{\theta}$ and that $\hat{\boldsymbol{\theta}}$ is in the interior of $\Theta$. So $\hat{\boldsymbol{\theta}}$, being the interior solution to the problem of maximizing $Q_{n}(\boldsymbol{\theta})$, satisfies the first-order conditions


\begin{align*}
\underset{(p \times 1)}{\mathbf{0}} & =\frac{\partial Q_{n}(\hat{\boldsymbol{\theta}})}{\partial \boldsymbol{\theta}} \\
& =\frac{1}{n} \sum_{t=1}^{n} \mathbf{s}\left(\mathbf{w}_{t} ; \hat{\boldsymbol{\theta}}\right) \quad(\text { by (7.3.1) and (7.3.2)) } \tag{7.3.4}
\end{align*}


We now use the following result from calculus:

Mean Value Theorem: Let $\mathbf{h}: \mathbb{R}^{p} \rightarrow \mathbb{R}^{q}$ be continuously differentiable. Then $\mathbf{h}(\mathbf{x})$ admits the mean value expansion

$$
\underset{(q \times 1)}{\mathbf{h}(\mathbf{x})}=\underset{(q \times 1)}{\mathbf{h}\left(\mathbf{x}_{0}\right)}+\underset{(q \times p)}{\frac{\partial \mathbf{h}(\overline{\mathbf{x}})}{\partial \mathbf{x}^{\prime}}} \underset{(p \times 1)}{\left(\mathbf{x}-\mathbf{x}_{0}\right)},
$$

where $\overline{\mathbf{x}}$ is a mean value lying between $\mathbf{x}$ and $\mathbf{x}_{0} .^{10}$

Setting $q=p, \mathbf{x}=\hat{\boldsymbol{\theta}}, \mathbf{x}_{0}=\boldsymbol{\theta}_{0}$, and $\mathbf{h}(\cdot)=\frac{\partial Q_{n}(\cdot)}{\partial \boldsymbol{\theta}}$ in the Mean Value Theorem, we obtain the following mean value expansion:


\begin{align*}
\frac{\partial Q_{n}(\hat{\boldsymbol{\theta}})}{\partial \boldsymbol{\theta}} & =\frac{\partial Q_{n}\left(\boldsymbol{\theta}_{0}\right)}{\partial \boldsymbol{\theta}}+\frac{\partial^{2} Q_{n}(\overline{\boldsymbol{\theta}})}{\partial \boldsymbol{\theta} \partial \boldsymbol{\theta}^{\prime}}\left(\hat{\boldsymbol{\theta}}-\boldsymbol{\theta}_{0}\right) \\
& =\frac{1}{n} \sum_{t=1}^{n} \mathbf{s}\left(\mathbf{w}_{t} ; \boldsymbol{\theta}_{0}\right)+\left[\frac{1}{n} \sum_{t=1}^{n} \mathbf{H}\left(\mathbf{w}_{t} ; \overline{\boldsymbol{\theta}}\right)\right] \underset{(p \times p)}{\left(\hat{\boldsymbol{\theta}}-\boldsymbol{\theta}_{0}\right)} \quad(\text { by }(7.3 .1)-(7.3 .3)) \tag{7.3.5}
\end{align*}


where $\overline{\boldsymbol{\theta}}$ is a mean value that lies between $\hat{\boldsymbol{\theta}}$ and $\boldsymbol{\theta}_{0}$. The continuous differentiability requirement of the Mean Value Theorem is satisfied if $m\left(\mathbf{w}_{t} ; \boldsymbol{\theta}\right)$ is twice continuously differentiable with respect to $\theta$. Combining this equation with the first-order condition above, we obtain


\begin{equation*}
\underset{(p \times 1)}{\mathbf{0}}=\frac{1}{n} \sum_{t=1}^{n} \mathbf{s}\left(\mathbf{w}_{t} ; \boldsymbol{\theta}_{0}\right)+\left[\frac{1}{n} \sum_{t=1}^{n} \mathbf{H}\left(\mathbf{w}_{t} ; \overline{\boldsymbol{\theta}}\right)\right]\left(\hat{\boldsymbol{\theta}}-\boldsymbol{\theta}_{0}\right) . \tag{7.3.6}
\end{equation*}


Assuming that $\frac{1}{n} \sum_{t=1}^{n} \mathbf{H}\left(\mathbf{w}_{t} ; \overline{\boldsymbol{\theta}}\right)$ is nonsingular, this equation can be solved for $\hat{\boldsymbol{\theta}}-\boldsymbol{\theta}_{0}$ to yield


\begin{equation*}
\sqrt{n}\left(\hat{\boldsymbol{\theta}}-\boldsymbol{\theta}_{0}\right)=-\left[\frac{1}{n} \sum_{t=1}^{n} \mathbf{H}\left(\mathbf{w}_{t} ; \overline{\boldsymbol{\theta}}\right)\right]^{-1} \frac{1}{\sqrt{n}} \sum_{t=1}^{n} \mathbf{s}\left(\mathbf{w}_{t} ; \boldsymbol{\theta}_{0}\right) . \tag{7.3.7}
\end{equation*}


This expression for ( $\sqrt{n}$ times) the sampling error will be referred to as the mean value expansion for the sampling error. Note that the score vector $\mathbf{s}\left(\mathbf{w}_{t} ; \boldsymbol{\theta}\right)$ is evaluated at the true parameter value $\boldsymbol{\theta}_{0}$.

Now, since $\overline{\boldsymbol{\theta}}$ lies between $\boldsymbol{\theta}_{0}$ and $\hat{\boldsymbol{\theta}}, \overline{\boldsymbol{\theta}}$ is consistent for $\boldsymbol{\theta}_{0}$ if $\hat{\boldsymbol{\theta}}$ is. If $\left\{\mathbf{w}_{t}\right\}$ is ergodic stationary, it is natural to conjecture that


\begin{equation*}
\frac{1}{n} \sum_{t=1}^{n} \mathbf{H}\left(\mathbf{w}_{t} ; \overline{\boldsymbol{\theta}}\right) \underset{\mathrm{p}}{\rightarrow} \mathrm{E}\left[\mathbf{H}\left(\mathbf{w}_{t} ; \boldsymbol{\theta}_{0}\right)\right] . \tag{7.3.8}
\end{equation*}


\footnotetext{${ }^{10}$ The Mean Value Theorem only applies to individual elements of $\mathbf{h}$, so that $\overline{\mathbf{x}}$ actually differs from element to element of the vector equation. This complication does not affect the discussion in the text.
}The ergodic stationarity of $\mathbf{w}_{t}$ and consistency of $\overline{\boldsymbol{\theta}}$ alone, however, are not enough to ensure this; some technical condition needs to be assumed. One such technical condition is the uniform convergence of $\frac{1}{n} \sum_{t=1}^{n} \mathbf{H}\left(\mathbf{w}_{t} ; \cdot\right)$ to $\mathrm{E}\left[\mathbf{H}\left(\mathbf{w}_{t} ; \cdot\right)\right]$ in a neighborhood of $\boldsymbol{\theta}_{0} .{ }^{11}$ But by the Uniform Convergence Theorem, the uniform convergence of $\frac{1}{n} \sum_{t=1}^{n} \mathbf{H}\left(\mathbf{w}_{t} ; \cdot\right)$ is satisfied if the following dominance condition is satisfied for the Hessian: for some neighborhood $\mathcal{N}$ of $\boldsymbol{\theta}_{0}, \mathrm{E}\left[\sup _{\boldsymbol{\theta} \in \mathcal{N}}\left\|\mathbf{H}\left(\mathbf{w}_{t} ; \boldsymbol{\theta}\right)\right\|\right]<\infty$ (this is condition (4) below). ${ }^{12}$ This is a sufficient condition for (7.3.8); if you can directly verify (7.3.8), then there is no need to verify the dominance condition for asymptotic normality.

Finally, if (7.3.8) holds and if


\begin{equation*}
\frac{1}{\sqrt{n}} \sum_{t=1}^{n} \mathbf{s}\left(\mathbf{w}_{t} ; \boldsymbol{\theta}_{0}\right) \underset{\mathrm{d}}{\rightarrow} N(\mathbf{0}, \boldsymbol{\Sigma}), \tag{7.3.9}
\end{equation*}


then by the Slutzky theorem (see Lemma 2.4(c)) we have


\begin{equation*}
\sqrt{n}\left(\hat{\boldsymbol{\theta}}-\boldsymbol{\theta}_{0}\right) \underset{\mathrm{d}}{\rightarrow} N\left(\boldsymbol{0},\left(\mathrm{E}\left[\mathbf{H}\left(\mathbf{w}_{t} ; \boldsymbol{\theta}_{0}\right)\right]\right)^{-1} \boldsymbol{\Sigma}\left(\mathrm{E}\left[\mathbf{H}\left(\mathbf{w}_{t} ; \boldsymbol{\theta}_{0}\right)\right]\right)^{-1}\right) . \tag{7.3.10}
\end{equation*}


Collecting the assumptions we have made so far, we have

Proposition 7.8 (Asymptotic normality of M-estimators): Suppose that the conditions of either Proposition 7.3 or Proposition 7.4 are satisfied, so that $\left\{\mathbf{w}_{t}\right\}$ is ergodic stationary and the M-estimator $\hat{\boldsymbol{\theta}}$ defined by (7.1.1) and (7.1.2) is consistent. Suppose, further, that\\
(1) $\boldsymbol{\theta}_{0}$ is in the interior of $\Theta$,\\
(2) $m\left(\mathbf{w}_{t} ; \boldsymbol{\theta}\right)$ is twice continuously differentiable in $\boldsymbol{\theta}$ for any $\mathbf{w}_{t}$,\\
(3) $\frac{1}{\sqrt{n}} \sum_{t=1}^{n} \mathbf{s}\left(\mathbf{w}_{t} ; \boldsymbol{\theta}_{0}\right) \rightarrow_{\mathrm{d}} N(\mathbf{0}, \mathbf{\Sigma}), \mathbf{\Sigma}$ positive definite, where $\mathbf{s}\left(\mathbf{w}_{t} ; \boldsymbol{\theta}\right)$ is defined in (7.3.2),

\footnotetext{${ }^{11}$ This is a consequence of the following result (see, e.g., Theorem 4.1.5 of Amemiya, 1985):\\
Suppose a sequence of random functions $h_{n}(\boldsymbol{\theta})$ converges to a nonstochastic function $h_{0}(\boldsymbol{\theta})$ uniformly in $\boldsymbol{\theta}$ over an open neighborhood of $\boldsymbol{\theta}_{0}$. Then $\operatorname{plim}_{n \rightarrow \infty} h_{n}(\tilde{\boldsymbol{\theta}})=h_{0}\left(\boldsymbol{\theta}_{0}\right)$ if plim $\tilde{\boldsymbol{\theta}}=\boldsymbol{\theta}_{0}$ and $h_{0}(\boldsymbol{\theta})$ is continuous at $\boldsymbol{\theta}_{0}$.

Set $h_{n}(\theta)$ to $\frac{1}{n} \sum_{t=1}^{n} \mathbf{H}\left(\mathbf{w}_{t} ; \cdot\right)$ and $h_{0}(\theta)$ to $\mathrm{E}\left[\mathbf{H}\left(\mathbf{w}_{t} ; \theta\right)\right]$. Continuity of $\mathrm{E}\left[\mathbf{H}\left(\mathbf{w}_{t} ; \theta\right)\right]$ will be assured by the Uniform Convergence Theorem. The uniform convergence needs to be only over a neighborhood of $\boldsymbol{\theta}_{0}$, not over the entire parameter space, because $\tilde{\theta}$ is close to $\theta_{0}$ for $n$ sufficiently large.\\
${ }^{12}$ A matrix can be viewed as a vector whose elements have been rearranged. So the Euclidean norm of a matrix, such as $\left\|\mathbf{H}\left(\mathbf{w}_{t} ; \boldsymbol{\theta}\right)\right\|$, is the square root of the sum of squares of all its elements.
}
(4) (local dominance condition on the Hessian) for some neighborhood $\mathcal{N}$ of $\theta_{0}$,
$$
\mathrm{E}\left[\sup _{\boldsymbol{\theta} \in \mathcal{N}}\left\|\mathbf{H}\left(\mathbf{w}_{t} ; \boldsymbol{\theta}\right)\right\|\right]<\infty,
$$
so that for any consistent estimator $\tilde{\boldsymbol{\theta}}, \frac{1}{n} \sum_{t=1}^{n} \mathbf{H}\left(\mathbf{w}_{t} ; \tilde{\boldsymbol{\theta}}\right) \rightarrow_{\mathrm{p}} \mathrm{E}\left[\mathbf{H}\left(\mathbf{w}_{t} ; \boldsymbol{\theta}_{0}\right)\right]$, where $\mathbf{H}\left(\mathbf{w}_{t} ; \boldsymbol{\theta}\right)$ is defined in (7.3.3),\\
(5) $\mathrm{E}\left[\mathrm{H}\left(\mathrm{w}_{t}, \boldsymbol{\theta}_{0}\right)\right]$ is nonsingular.

Then $\hat{\boldsymbol{\theta}}$ is asymptotically normal with

$$
\operatorname{Avar}(\hat{\boldsymbol{\theta}})=\left(\mathrm{E}\left[\mathbf{H}\left(\mathbf{w}_{t} ; \boldsymbol{\theta}_{0}\right)\right]\right)^{-1} \boldsymbol{\Sigma}\left(\mathrm{E}\left[\mathbf{H}\left(\mathbf{w}_{t} ; \boldsymbol{\theta}_{0}\right)\right]\right)^{-1} .
$$

(This is Theorem 4.1.3 of Amemiya (1985) adapted to M-estimators.) Two remarks are in order.

\begin{itemize}
  \item Of the assumptions we have made in the derivation of asymptotic normality, the following are not listed in the proposition: (i) $\hat{\boldsymbol{\theta}}$ is an interior point, and (ii) $\frac{1}{n} \sum_{t=1}^{n} \mathbf{H}\left(\mathbf{w}_{t} ; \overline{\boldsymbol{\theta}}\right)$ is nonsingular. It is intuitively clear that these conditions hold because $\hat{\boldsymbol{\theta}}$ converges in probability to an interior point $\boldsymbol{\theta}_{0}$, and $\frac{1}{n} \sum_{t=1}^{n} \mathbf{H}\left(\mathbf{w}_{t} ; \overline{\boldsymbol{\theta}}\right)$ converges in probability to a nonsingular matrix $\mathrm{E}\left[\mathbf{H}\left(\mathbf{w}_{t} ; \boldsymbol{\theta}_{0}\right)\right]$. See Newey and McFadden (1994, p. 2152) for a rigorous proof. This sort of technicality will be ignored in the rest of this chapter.
  \item If $\mathbf{w}_{t}$ is ergodic stationary, then so is $\mathbf{s}\left(\mathbf{w}_{t} ; \boldsymbol{\theta}_{0}\right)$ and the matrix $\boldsymbol{\Sigma}$ is the long run variance matrix of $\left\{\mathbf{s}\left(\mathbf{w}_{i} ; \boldsymbol{\theta}_{0}\right)\right\} .{ }^{13}$ A sufficient condition for (3) is Gordin's condition introduced in Section 6.5. So condition (3) in the proposition can be replaced by Gordin's condition on $\left\{\mathbf{s}\left(\mathbf{w}_{t} ; \boldsymbol{\theta}_{0}\right)\right\}$. It is satisfied, for example, if $\mathbf{w}_{t}$ is i.i.d. and $\mathrm{E}\left[\mathbf{s}\left(\mathbf{w}_{t} ; \boldsymbol{\theta}_{0}\right)\right]=\mathbf{0}$. The assumption that $\boldsymbol{\Sigma}$ is positive definite is not really needed for the conclusion of the proposition, but we might as well assume it here because in virtually all applications it is satisfied (or assumed) and also because it will be required in the discussion of hypothesis testing later in this chapter.
\end{itemize}

\section*{Consistent Asymptotic Variance Estimation}
To use this asymptotic result for hypothesis testing, we need a consistent estimate of

$$
\operatorname{Avar}(\hat{\boldsymbol{\theta}})=\left(\mathrm{E}\left[\mathbf{H}\left(\mathbf{w}_{t} ; \boldsymbol{\theta}_{0}\right)\right]\right)^{-1} \boldsymbol{\Sigma}\left(\mathrm{E}\left[\mathbf{H}\left(\mathbf{w}_{t} ; \boldsymbol{\theta}_{0}\right)\right]\right)^{-1} .
$$

${ }^{13}$ The long-run variance was introduced in Section 6.5.

Since $\hat{\boldsymbol{\theta}} \rightarrow_{\mathrm{p}} \boldsymbol{\theta}_{0}$, condition (4) of Proposition 7.8 implies that

$$
\frac{1}{n} \sum_{t=1}^{n} \mathbf{H}\left(\mathbf{w}_{t} ; \hat{\boldsymbol{\theta}}\right) \underset{\mathrm{p}}{\rightarrow} \mathrm{E}\left[\mathbf{H}\left(\mathbf{w}_{t} ; \boldsymbol{\theta}_{0}\right)\right] .
$$

Therefore, provided that there is available a consistent estimator $\widehat{\mathbf{\Sigma}}$ of $\mathbf{\Sigma}$,


\begin{equation*}
\widehat{\operatorname{Avar}(\hat{\boldsymbol{\theta}})}=\left\{\frac{1}{n} \sum_{t=1}^{n} \mathbf{H}\left(\mathbf{w}_{t} ; \hat{\boldsymbol{\theta}}\right)\right\}^{-1} \widehat{\boldsymbol{\Sigma}}\left\{\frac{1}{n} \sum_{t=1}^{n} \mathbf{H}\left(\mathbf{w}_{t} ; \hat{\boldsymbol{\theta}}\right)\right\}^{-1} \tag{7.3.11}
\end{equation*}


is a consistent estimator of the asymptotic variance matrix. To obtain $\widehat{\mathbf{\Sigma}}$, the methods introduced in Section 6.6, such as the VARHAC, can be applied to the estimated series $\left\{\mathbf{s}\left(\mathbf{w}_{t} ; \hat{\boldsymbol{\theta}}\right)\right\}$ under some suitable technical conditions.

\section*{Asymptotic Normality of Conditional ML}
We now specialize this asymptotic normality result to ML for the case of i.i.d. data. In ML, where the $m$ function is a log (conditional) likelihood, the first-order conditions (7.3.4) are called the likelihood equations. The fact that the $m$ function is a log (conditional) likelihood leads to a simplification of the asymptotic variance matrix of $\hat{\boldsymbol{\theta}}$. We show this for conditional ML; doing the same for unconditional or joint ML is easier.

Consider the two equalities in condition (3) of Proposition 7.9 below. The second equality is called the information matrix equality (not to be confused with the Kullback-Leibler information inequality). It is so called because $\mathrm{E}\left[\mathbf{s}\left(\mathbf{w}_{t} ; \boldsymbol{\theta}_{0}\right) \mathbf{s}\left(\mathbf{w}_{t}\right.\right.$; $\left.\theta_{0}\right)^{\prime}$ ] is the information matrix for observation $t$. It is left as Analytical Exercise 2 to derive these two equalities under some technical conditions permitting the interchange of differentiation and integration. Here, we just note the implication of these two equalities for Proposition 7.8. If $\mathbf{w}_{t}$ is i.i.d., the Lindeberg-Levy CLT and the first equality (that $\mathrm{E}\left[\mathbf{s}\left(\mathbf{w}_{t} ; \boldsymbol{\theta}_{0}\right)\right]=\mathbf{0}$ ) imply


\begin{equation*}
\frac{1}{\sqrt{n}} \sum_{t=1}^{n} \mathbf{s}\left(\mathbf{w}_{t} ; \boldsymbol{\theta}_{0}\right) \rightarrow N(\mathbf{0}, \mathbf{\Sigma}) \text { where } \mathbf{\Sigma}=\mathrm{E}\left[\mathbf{s}\left(\mathbf{w}_{t} ; \boldsymbol{\theta}_{0}\right) \mathbf{s}\left(\mathbf{w}_{t} ; \boldsymbol{\theta}_{0}\right)^{\prime}\right] . \tag{7.3.12}
\end{equation*}


This and the information matrix equality imply that $\operatorname{Avar}(\hat{\boldsymbol{\theta}})$ in Proposition 7.8 is simplified in two ways as


\begin{equation*}
\operatorname{Avar}(\hat{\boldsymbol{\theta}})=-\left\{\mathrm{E}\left[\mathbf{H}\left(\mathbf{w}_{t} ; \boldsymbol{\theta}_{0}\right)\right]\right\}^{-1}=\left\{\mathrm{E}\left[\mathbf{s}\left(\mathbf{w}_{t} ; \boldsymbol{\theta}_{0}\right) \mathbf{s}\left(\mathbf{w}_{t} ; \boldsymbol{\theta}_{0}\right)^{\prime}\right]\right\}^{-1} . \tag{7.3.13}
\end{equation*}


Thus, we have proved

Proposition 7.9 (Asymptotic normality of conditional ML): Let $\mathbf{w}_{t}\left(\equiv\left(y_{t}, \mathbf{x}_{t}^{\prime}\right)^{\prime}\right)$ be i.i.d. Suppose the conditions of either Proposition 7.5 or Proposition 7.6 are satisfied, so that $\hat{\boldsymbol{\theta}} \rightarrow_{\mathrm{p}} \boldsymbol{\theta}_{0}$. Suppose, in addition, that\\
(1) $\boldsymbol{\theta}_{0}$ is in the interior of $\Theta$,\\
(2) $f\left(y_{t} \mid \mathbf{x}_{t} ; \boldsymbol{\theta}\right)$ is twice continuously differentiable in $\boldsymbol{\theta}$ for all $\left(y_{t}, \mathbf{x}_{t}\right)$,\\
(3) $\mathrm{E}\left[\mathbf{s}\left(\mathbf{w}_{t} ; \boldsymbol{\theta}_{0}\right)\right]=\mathbf{0}$ and $-\mathrm{E}\left[\mathbf{H}\left(\mathbf{w}_{t} ; \boldsymbol{\theta}_{0}\right)\right]=\mathrm{E}\left[\mathbf{s}\left(\mathbf{w}_{t} ; \boldsymbol{\theta}_{0}\right) \mathbf{s}\left(\mathbf{w}_{t} ; \boldsymbol{\theta}_{0}\right)^{\prime}\right]$, where $\mathbf{s}$ and $\mathbf{H}$ functions are defined in (7.3.2) and (7.3.3),\\
(4) (local dominance condition on the Hessian) for some neighborhood $\mathcal{N}$ of $\boldsymbol{\theta}_{0}$,

$$
\mathrm{E}\left[\sup _{\boldsymbol{\theta} \in \mathcal{N}}\left\|\mathbf{H}\left(\mathbf{w}_{t} ; \boldsymbol{\theta}\right)\right\|\right]<\infty,
$$

so that for any consistent estimator $\tilde{\boldsymbol{\theta}}, \frac{1}{n} \sum_{t=1}^{n} \mathbf{H}\left(\mathbf{w}_{t} ; \tilde{\boldsymbol{\theta}}\right) \rightarrow_{\mathrm{p}} \mathrm{E}\left[\mathbf{H}\left(\mathbf{w}_{t} ; \boldsymbol{\theta}_{0}\right)\right]$,\\
(5) $\mathrm{E}\left[\mathbf{H}\left(\mathbf{w}_{t} ; \boldsymbol{\theta}_{0}\right)\right]$ is nonsingular.

Then $\hat{\boldsymbol{\theta}}$ is asymptotically normal with $\operatorname{Avar}(\hat{\boldsymbol{\theta}})$ given by (7.3.13).

Two remarks about the proposition:

\begin{itemize}
  \item Condition (3) is stated as an assumption here because its derivation requires interchange of integration and differentiation (see Analytical Exercise 2). A technical condition under which the interchange is legal is readily available from calculus (see, e.g., Lemma 3.6 of Newey and McFadden, 1994). So condition (3) could be replaced by such a technical condition on $f\left(y_{t} \mid \mathbf{x}_{t} ; \boldsymbol{\theta}\right)$. We do not do it here because in most applications condition (3) can be verified directly.
  \item It should be clear from the above discussion how this proposition can be adapted to unconditional ML: simply replace $f\left(y_{t} \mid \mathbf{x}_{t} ; \boldsymbol{\theta}\right)$ by $f\left(\mathbf{w}_{t} ; \boldsymbol{\theta}\right)$.
\end{itemize}

To use this asymptotic result for hypothesis testing, we need a consistent estimate of $\operatorname{Avar}(\hat{\boldsymbol{\theta}})$. Noting that it can be written as $-\left\{\mathrm{E}\left[\mathbf{H}\left(\mathbf{w}_{t} ; \boldsymbol{\theta}_{0}\right)\right]\right\}^{-1}$, a natural estimator is


\begin{equation*}
\text { first estimator of } \operatorname{Avar}(\hat{\boldsymbol{\theta}})=-\left\{\frac{1}{n} \sum_{t=1}^{n} \mathbf{H}\left(\mathbf{w}_{t} ; \hat{\boldsymbol{\theta}}\right)\right\}^{-1} . \tag{7.3.14}
\end{equation*}


Since $\hat{\boldsymbol{\theta}} \rightarrow_{\mathrm{p}} \boldsymbol{\theta}_{0}$, this estimator is consistent by condition (4) of Proposition 7.9.

Another estimator, based on the relation $\operatorname{Avar}(\hat{\boldsymbol{\theta}})=\left\{\mathrm{E}\left[\mathbf{s}\left(\mathbf{w}_{t} ; \boldsymbol{\theta}_{0}\right) \mathbf{s}\left(\mathbf{w}_{t} ; \boldsymbol{\theta}_{0}\right)^{\prime}\right]\right\}^{-1}$, is


\begin{equation*}
\text { second estimator of } \operatorname{Avar}(\hat{\boldsymbol{\theta}})=\left\{\frac{1}{n} \sum_{t=1}^{n} \mathbf{s}\left(\mathbf{w}_{t} ; \hat{\boldsymbol{\theta}}\right) \mathbf{s}\left(\mathbf{w}_{t} ; \hat{\boldsymbol{\theta}}\right)^{\prime}\right\}^{-1} . \tag{7.3.15}
\end{equation*}


To insure consistency for this estimator, we need to make an assumption in addition to the hypothesis of Proposition 7.9 (see, e.g., Theorem 4.4 of Newey and McFadden, 1994). We will not show it here because in many applications the consistency of this second estimator can be verified directly.

There is no compelling reason for preferring one estimator of $\operatorname{Avar}(\hat{\boldsymbol{\theta}})$ over the other. The second estimator, which requires only the first derivative of the log likelihood, is easier to compute than the first. This can be an important consideration when it is impossible to calculate the derivatives of the density function analytically and some numerical method must be used. On the other hand, in at least some cases the first provides a better finite-sample approximation of the asymptotic distribution (see, e.g., Davidson and MacKinnon, 1993).

\section*{Two Examples}
To get a better grasp of the asymptotic normality results about conditional ML, it is useful to see how the results can be applied to familiar examples. We consider two such examples.

Example 7.10 (Asymptotic normality in linear regression model): To reproduce the log conditional density for observation $t$ for the linear regression model with normal errors,


\begin{equation*}
\log f\left(y_{t} \mid \mathbf{x}_{t} ; \boldsymbol{\beta}, \sigma^{2}\right)=-\frac{1}{2} \log (2 \pi)-\frac{1}{2} \log \left(\sigma^{2}\right)-\frac{\left(y_{t}-\mathbf{x}_{t}^{\prime} \boldsymbol{\beta}\right)^{2}}{2 \sigma^{2}} . \tag{7.3.16}
\end{equation*}


With $\Theta=\mathbb{R}^{K} \times \mathbb{R}_{++}$(where $K$ is the dimension of $\boldsymbol{\beta}$ ), condition (1) of Proposition 7.9 is satisfied. Condition (2) is obviously satisfied. To verify (3), a routine calculation yields: ${ }^{14}$


\begin{gather*}
\mathbf{s}\left(\mathbf{w}_{t} ; \boldsymbol{\theta}\right)=\left[\begin{array}{c}
\frac{1}{\sigma^{2}} \mathbf{x}_{t} \cdot \hat{\varepsilon}_{t} \\
-\frac{1}{2 \sigma^{2}}+\frac{1}{2 \sigma^{4}} \hat{\varepsilon}_{t}^{2}
\end{array}\right], \mathbf{H}\left(\mathbf{w}_{t} ; \boldsymbol{\theta}\right)=\left[\begin{array}{cc}
-\frac{1}{\sigma^{2}} \mathbf{x}_{t} \mathbf{x}_{t}^{\prime} & -\frac{1}{\sigma^{4}} \mathbf{x}_{t} \cdot \hat{\varepsilon}_{t} \\
-\frac{1}{\sigma^{4}} \mathbf{x}_{t}^{\prime} \cdot \hat{\varepsilon}_{t} & \frac{1}{2 \sigma^{4}}-\frac{1}{\sigma^{6}} \hat{\varepsilon}_{t}^{2}
\end{array}\right], \\
\mathbf{s}\left(\mathbf{w}_{t} ; \boldsymbol{\theta}\right) \mathbf{s}\left(\mathbf{w}_{t} ; \boldsymbol{\theta}\right)^{\prime}=\left[\begin{array}{cc}
\frac{1}{\sigma^{4}} \mathbf{x}_{t} \mathbf{x}_{t}^{\prime} \hat{\varepsilon}_{t}^{2} & -\frac{1}{2 \sigma^{4}} \mathbf{x}_{t} \cdot \hat{\varepsilon}_{t}+\frac{1}{2 \sigma^{6}} \mathbf{x}_{t} \cdot \hat{\varepsilon}_{t}^{3} \\
-\frac{1}{2 \sigma^{4}} \mathbf{x}_{t}^{\prime} \cdot \hat{\varepsilon}_{t}+\frac{1}{2 \sigma^{6}} \mathbf{x}_{t}^{\prime} \cdot \hat{\varepsilon}_{t}^{3} & \frac{1}{4 \sigma^{4}}-\frac{1}{2 \sigma^{6}} \hat{\varepsilon}_{t}^{2}+\frac{1}{4 \sigma^{8}} \hat{\varepsilon}_{t}^{4}
\end{array}\right], \tag{7.3.17}
\end{gather*}


\footnotetext{${ }^{14}$ For the parameter $\sigma^{2}$ the differentiation is with respect to $\sigma^{2}$ rather than $\sigma$.
}
where $\mathbf{w}_{t}=\left(y_{t}, \mathbf{x}_{t}^{\prime}\right)^{\prime}, \boldsymbol{\theta}=\left(\boldsymbol{\beta}^{\prime}, \sigma^{2}\right)^{\prime}$ and $\hat{\varepsilon}_{t} \equiv y_{t}-\mathbf{x}_{t}^{\prime} \boldsymbol{\beta}$ (which should be distinguished from $\left.\varepsilon_{t} \equiv y_{t}-\mathbf{x}_{t}^{\prime} \boldsymbol{\beta}_{0}\right)$. So for $\boldsymbol{\theta}=\boldsymbol{\theta}_{0}$ the $\hat{\varepsilon}_{t}$ in these expressions can be replaced by $\varepsilon_{t}$. In the linear regression model, $\mathrm{E}\left(\varepsilon_{t} \mid \mathbf{x}_{t}\right)=0$. Also, since $\varepsilon_{t}$ is $N\left(0, \sigma_{0}^{2}\right)$, we have $\mathrm{E}\left(\varepsilon_{t}^{3}\right)=0$ and $\mathrm{E}\left(\varepsilon_{t}^{4}\right)=3 \sigma_{0}^{4}$. Using these relations, it is easy to verify (3). In particular,
\[
-\mathrm{E}\left[\mathbf{H}\left(\mathbf{w}_{t} ; \boldsymbol{\theta}_{0}\right)\right]=\mathrm{E}\left[\mathbf{s}\left(\mathbf{w}_{t} ; \boldsymbol{\theta}_{0}\right) \mathbf{s}\left(\mathbf{w}_{t} ; \boldsymbol{\theta}_{0}\right)^{\prime}\right]=\left[\begin{array}{cc}
\frac{1}{\sigma_{0}^{2}} \mathrm{E}\left(\mathbf{x}_{t} \mathbf{x}_{t}^{\prime}\right) & \mathbf{0}  \tag{7.3.18}\\
\boldsymbol{0}^{\prime} & \frac{1}{2 \sigma_{0}^{4}}
\end{array}\right] .
\]

If $\mathrm{E}\left(\mathbf{x}_{t} \mathbf{x}_{t}^{\prime}\right)$ is nonsingular, then $\mathrm{E}\left[\mathbf{H}\left(\mathbf{w}_{t} ; \boldsymbol{\theta}_{0}\right)\right]$ is nonsingular and (5) is satisfied. Regarding condition (4), let $\tilde{\varepsilon}_{t} \equiv y_{t}-\mathbf{x}_{t}^{\prime} \widetilde{\boldsymbol{\beta}}$ for some consistent estimator $\widetilde{\boldsymbol{\beta}}$ and $\tilde{\sigma}^{2}$. Condition (7.3.8) in this example is


\begin{gather*}
{\left[\begin{array}{cc}
-\frac{1}{\tilde{\sigma}^{2}} \frac{1}{n} \sum_{t=1}^{n} \mathbf{x}_{t} \mathbf{x}_{t}^{\prime} & -\frac{1}{\left(\tilde{\sigma}^{2}\right)^{2}} \frac{1}{n} \sum_{t=1}^{n} \mathbf{x}_{t} \cdot \tilde{\varepsilon}_{t} \\
-\frac{1}{\left(\tilde{\sigma}^{2}\right)^{2}} \frac{1}{n} \sum_{t=1}^{n} \mathbf{x}_{t}^{\prime} \cdot \tilde{\varepsilon}_{t} & \frac{1}{2\left(\tilde{\sigma}^{2}\right)^{2}}-\frac{1}{\left(\tilde{\sigma}^{2}\right)^{3}} \frac{1}{n} \sum_{t=1}^{n} \tilde{\varepsilon}_{t}^{2}
\end{array}\right]} \\
\overrightarrow{\mathrm{p}}\left[\begin{array}{cc}
-\frac{1}{\sigma_{0}^{2}} \mathrm{E}\left(\mathbf{x}_{t} \mathbf{x}_{t}^{\prime}\right) & \mathbf{0} \\
\mathbf{0}^{\prime} & -\frac{1}{2 \sigma_{0}^{4}}
\end{array}\right] \tag{7.3.19}
\end{gather*}


It is straightforward to show this, given that $\tilde{\varepsilon}_{t}=\varepsilon_{t}-\mathbf{x}_{t}^{\prime}\left(\widetilde{\boldsymbol{\beta}}-\boldsymbol{\beta}_{0}\right)$. We conclude that all the conditions of Proposition 7.9 are satisfied if $\mathrm{E}\left(\mathbf{x}_{t} \mathbf{x}_{t}^{\prime}\right)$ is nonsingular.

Example 7.11 (Asymptotic normality of probit ML): To reproduce the log conditional likelihood of the probit model,


\begin{equation*}
\log f\left(y_{t} \mid \mathbf{x}_{t} ; \boldsymbol{\theta}\right)=y_{t} \log \Phi\left(\mathbf{x}_{t}^{\prime} \boldsymbol{\theta}\right)+\left(1-y_{t}\right) \log \left[1-\Phi\left(\mathbf{x}_{t}^{\prime} \boldsymbol{\theta}\right)\right] \tag{7.3.20}
\end{equation*}


With $\Theta=\mathbb{R}^{p}$, condition (1) of Proposition 7.9 is satisfied. Condition (2) is obviously satisfied. To verify (3), a tedious calculation yields


\begin{gather*}
\mathbf{s}\left(\mathbf{w}_{t} ; \boldsymbol{\theta}\right)=\frac{\left(y_{t}-\Phi\left(\mathbf{x}_{t}^{\prime} \boldsymbol{\theta}\right)\right) \phi\left(\mathbf{x}_{t}^{\prime} \boldsymbol{\theta}\right)}{\left(1-\Phi\left(\mathbf{x}_{t}^{\prime} \boldsymbol{\theta}\right)\right) \Phi\left(\mathbf{x}_{t}^{\prime} \boldsymbol{\theta}\right)} \mathbf{x}_{t}  \tag{7.3.21}\\
\mathbf{H}\left(\mathbf{w}_{t} ; \boldsymbol{\theta}\right)=\left\{-\left[\frac{y_{t}-\Phi\left(\mathbf{x}_{t}^{\prime} \boldsymbol{\theta}\right)}{\Phi\left(\mathbf{x}_{t}^{\prime} \boldsymbol{\theta}\right)\left(1-\Phi\left(\mathbf{x}_{t}^{\prime} \boldsymbol{\theta}\right)\right)}\right]^{2}\left[\phi\left(\mathbf{x}_{t}^{\prime} \boldsymbol{\theta}\right)\right]^{2}\right. \\
\left.+\left[\frac{y_{t}-\Phi\left(\mathbf{x}_{t}^{\prime} \boldsymbol{\theta}\right)}{\Phi\left(\mathbf{x}_{t}^{\prime} \boldsymbol{\theta}\right)\left(1-\Phi\left(\mathbf{x}_{t}^{\prime} \boldsymbol{\theta}\right)\right)}\right] \phi^{\prime}\left(\mathbf{x}_{t}^{\prime} \boldsymbol{\theta}\right)\right\} \mathbf{x}_{t} \mathbf{x}_{t}^{\prime} \tag{7.3.22}
\end{gather*}


(To derive (7.3.22), use the fact that $y_{t}^{2}=y_{t}$.) Here, $\Phi(\cdot)$ is the cumulative\\
density function of $N(0,1)$ and $\phi(\cdot) \equiv \Phi^{\prime}(\cdot)$ is its density function. It is then easy to prove the conditional version of (3):


\begin{gather*}
\mathrm{E}\left[\mathbf{s}\left(\mathbf{w}_{t} ; \boldsymbol{\theta}_{0}\right) \mid \mathbf{x}_{t}\right]=\mathbf{0} \text { and } \\
-\mathrm{E}\left[\mathbf{H}\left(\mathbf{w}_{t} ; \boldsymbol{\theta}_{0}\right) \mid \mathbf{x}_{t}\right]=\mathrm{E}\left[\mathbf{s}\left(\mathbf{w}_{t} ; \boldsymbol{\theta}_{0}\right) \mathbf{s}\left(\mathbf{w}_{t} ; \boldsymbol{\theta}_{0}\right)^{\prime} \mid \mathbf{x}_{t}\right] . \tag{7.3.23}
\end{gather*}


So (3) is satisfied by the Law of Total Expectations. In particular, for probit, it is easy to show that


\begin{equation*}
-\mathrm{E}\left[\mathbf{H}\left(\mathbf{w}_{t} ; \boldsymbol{\theta}_{0}\right)\right]=\mathrm{E}\left[\mathbf{s}\left(\mathbf{w}_{t} ; \boldsymbol{\theta}_{0}\right) \mathbf{s}\left(\mathbf{w}_{t} ; \boldsymbol{\theta}_{0}\right)^{\prime}\right]=\mathrm{E}\left[\lambda\left(\mathbf{x}_{t}^{\prime} \boldsymbol{\theta}_{0}\right) \lambda\left(-\mathbf{x}_{t}^{\prime} \boldsymbol{\theta}_{0}\right) \mathbf{x}_{t} \mathbf{x}_{t}^{\prime}\right], \tag{7.3.24}
\end{equation*}


where


\begin{equation*}
\lambda(v) \equiv \frac{\phi(v)}{1-\Phi(v)} \tag{7.3.25}
\end{equation*}


is called the inverse Mill's ratio or the hazard for $N(0,1)$. It can be shown that the term in braces in (7.3.22) is between 0 and 2 (see Newey and McFadden, 1994, p. 2147). So


\begin{equation*}
\left\|\mathbf{H}\left(\mathbf{w}_{t} ; \boldsymbol{\theta}\right)\right\| \leq 2\left\|\mathbf{x}_{t} \mathbf{x}_{t}^{\prime}\right\| . \tag{7.3.26}
\end{equation*}


The Euclidean norm $\left\|\mathbf{x}_{t} \mathbf{x}_{t}^{\prime}\right\|$ is the square root of the sum of squares of the elements of $\mathbf{x}_{t} \mathbf{x}_{t}^{\prime}$. Therefore, the expectation of $\left\|\mathbf{x}_{t} \mathbf{x}_{t}^{\prime}\right\|^{2}$ and hence that of $\left\|\mathbf{x}_{t} \mathbf{x}_{t}^{\prime}\right\|$ are finite if $\mathrm{E}\left(\mathbf{x}_{t} \mathbf{x}_{t}^{\prime}\right)$ (exists and) is finite. So the local dominance condition on the Hessian (condition (4)) is satisfied if $\mathrm{E}\left(\mathbf{x}_{t} \mathbf{x}_{t}^{\prime}\right)$ is nonsingular. It can be shown that condition (5) is satisfied if $\mathrm{E}\left(\mathbf{x}_{t} \mathbf{x}_{t}^{\prime}\right)$ is nonsingular (see footnote 29 on p. 2147 of Newey and McFadden, 1994). We conclude, again, that all the conditions of Proposition 7.9 are satisfied if $\mathrm{E}\left(\mathbf{x}_{t} \mathbf{x}_{t}^{\prime}\right)$ is nonsingular.

\section*{Asymptotic Normality of GMM}
Now turn to GMM. The GMM objective function is


\begin{equation*}
Q_{n}(\boldsymbol{\theta})=-\frac{1}{2} \mathbf{g}_{n}(\boldsymbol{\theta})^{\prime} \widehat{\mathbf{W}} \mathbf{g}_{n}(\boldsymbol{\theta}) \text { with } \underset{(K \times 1)}{\mathbf{g}_{n}(\boldsymbol{\theta})} \equiv \frac{1}{n} \sum_{t=1}^{n} \mathbf{g}\left(\mathbf{w}_{t} ; \boldsymbol{\theta}\right) \tag{7.3.27}
\end{equation*}


As in the case of M-estimators, we apply the Mean Value Theorem to the first-order condition for maximization, but the theorem will be applied to $\mathbf{g}_{n}(\boldsymbol{\theta})$, not its first derivative. Thus, unlike in the case of M-estimators, the objective function needs to be continuously differentiable only once, not twice. This reflects the fact that the sample average enters the objective function differently for the case of GMM.

Assuming that $\hat{\boldsymbol{\theta}}$ is an interior point, the first-order condition for maximization of the objective function is


\begin{equation*}
\underset{(p \times 1)}{\mathbf{0}}=\frac{\partial Q_{n}(\hat{\boldsymbol{\theta}})}{\partial \underset{(p \times 1)}{\partial \boldsymbol{\theta}}}=-\underset{(p \times K)}{\mathbf{G}_{n}(\hat{\boldsymbol{\theta}})^{\prime}} \underset{(K \times K)}{\widehat{\mathbf{W}}} \underset{(K \times 1)}{\mathbf{g}_{n}(\hat{\boldsymbol{\theta}})}, \tag{7.3.28}
\end{equation*}


where $\mathbf{G}_{n}(\boldsymbol{\theta})$ is the Jacobian of $\mathbf{g}_{n}(\boldsymbol{\theta})$ :


\begin{equation*}
\underset{(K \times p)}{\mathbf{G}_{n}(\boldsymbol{\theta})} \equiv \frac{\partial \mathbf{g}_{n}(\boldsymbol{\theta})}{\partial \boldsymbol{\theta}^{\prime}} . \tag{7.3.29}
\end{equation*}


Now apply the Mean Value Theorem to $\mathbf{g}_{n}(\boldsymbol{\theta})$, not to $\partial Q_{n}(\boldsymbol{\theta}) / \partial \boldsymbol{\theta}$ as in M-estimation, to obtain the mean value expansion


\begin{equation*}
\underset{(K \times 1)}{\mathbf{g}_{n}(\hat{\boldsymbol{\theta}})}=\underset{(K \times 1)}{\mathbf{g}_{n}\left(\boldsymbol{\theta}_{0}\right)}+\underset{(K \times p)}{\mathbf{G}_{n}(\overline{\boldsymbol{\theta}})} \underset{(p \times 1)}{\left(\hat{\boldsymbol{\theta}}-\boldsymbol{\theta}_{0}\right)} . \tag{7.3.30}
\end{equation*}


Substituting this into the first-order condition above, we obtain


\begin{equation*}
\underset{(p \times 1)}{\mathbf{0}}=-\underset{(p \times K)}{\mathbf{G}_{n}(\hat{\boldsymbol{\theta}})^{\prime}} \underset{(K \times K)}{\widehat{\mathbf{W}}} \underset{(K \times 1)}{\mathbf{g}_{n}\left(\boldsymbol{\theta}_{0}\right)}-\underset{(p \times K)}{\mathbf{G}_{n}(\hat{\boldsymbol{\theta}})^{\prime}} \underset{(K \times K)}{\widehat{\mathbf{W}}} \underset{(K \times p)}{\mathbf{G}_{n}(\overline{\boldsymbol{\theta}})} \underset{(p \times 1)}{\left(\hat{\boldsymbol{\theta}}-\boldsymbol{\theta}_{0}\right)} . \tag{7.3.31}
\end{equation*}


Solving this for $\left(\hat{\boldsymbol{\theta}}-\boldsymbol{\theta}_{0}\right)$ and noting that $\mathbf{g}_{n}(\boldsymbol{\theta})=\frac{1}{n} \sum_{t=1}^{n} \mathbf{g}\left(\mathbf{w}_{t} ; \boldsymbol{\theta}\right)$, we obtain


\begin{equation*}
\underset{(p \times 1)}{\sqrt{n}\left(\hat{\boldsymbol{\theta}}-\boldsymbol{\theta}_{0}\right)}=-\left[\underset{(p \times K)}{\mathbf{G}_{n}(\hat{\boldsymbol{\theta}})^{\prime}} \underset{(K \times K)}{\widehat{\mathbf{W}}} \underset{(K \times p)}{\mathbf{G}_{n}(\overline{\boldsymbol{\theta}})}\right]^{-1} \underset{(p \times K)}{\mathbf{G}_{n}(\hat{\boldsymbol{\theta}})^{\prime}} \underset{(K \times K)}{\widehat{\mathbf{W}}} \frac{1}{\sqrt{n}} \sum_{t=1}^{n} \mathbf{g}\left(\mathbf{w}_{t} ; \boldsymbol{\theta}_{0}\right) \tag{7.3.32}
\end{equation*}


In this expression, the Jacobian $\mathbf{G}_{n}(\boldsymbol{\theta})$ of $\mathbf{g}_{n}(\boldsymbol{\theta})$ is evaluated at two different points, $\hat{\boldsymbol{\theta}}$ and $\overline{\boldsymbol{\theta}}$. Since $\mathbf{g}_{n}(\boldsymbol{\theta})=\frac{1}{n} \sum_{t=1}^{n} \mathbf{g}\left(\mathbf{w}_{t} ; \boldsymbol{\theta}\right)$, the Jacobian at any given estimator $\tilde{\boldsymbol{\theta}}$ can be written as


\begin{equation*}
\mathbf{G}_{n}(\tilde{\boldsymbol{\theta}})=\frac{1}{n} \sum_{t=1}^{n} \frac{\partial \mathbf{g}\left(\mathbf{w}_{t} ; \tilde{\boldsymbol{\theta}}\right)}{\partial \boldsymbol{\theta}^{\prime}} \tag{7.3.33}
\end{equation*}


If $\mathbf{w}_{t}$ is ergodic stationary and $\tilde{\boldsymbol{\theta}}$ is consistent, it is natural to conjecture that this expression converges to $\mathrm{E}\left[\partial \mathbf{g}\left(\mathbf{w}_{t} ; \boldsymbol{\theta}_{0}\right) / \partial \boldsymbol{\theta}^{\prime}\right]$. As was true for M-estimators, this conjecture is true if $\partial \mathbf{g}\left(\mathbf{w}_{t} ; \boldsymbol{\theta}\right) / \partial \boldsymbol{\theta}^{\prime}$ satisfies the suitable dominance condition specified in condition (4) below. It should now be clear that the GMM equivalent of Proposition 7.8 is

Proposition 7.10 (Asymptotic normality of GMM): Suppose that the conditions of Proposition 7.7 are satisfied, so that $\left\{\mathbf{w}_{t}\right\}$ is ergodic stationary, $\widehat{\mathbf{W}}(K \times K)$ converges in probability to a symmetric positive definite matrix $\mathbf{W}$, and the $p$-dimensional GMM estimator $\hat{\boldsymbol{\theta}}$ is consistent. Suppose, further, that\\
(I) $\boldsymbol{\theta}_{0}$ is in the interior of $\Theta$,\\
(2) $\mathbf{g}\left(\mathbf{w}_{t} ; \boldsymbol{\theta}\right)$ is continuously differentiable in $\boldsymbol{\theta}$ for any $\mathbf{w}_{t}$,\\
( $K \times 1$ )\\
(3) $\frac{1}{\sqrt{n}} \sum_{t=1}^{n} \mathbf{g}\left(\mathbf{w}_{t} ; \boldsymbol{\theta}_{0}\right) \rightarrow_{\mathrm{d}} N(\mathbf{0}, \underset{(K \times K)}{\mathbf{S}}), \mathbf{S}$ positive definite,\\
(4) (local dominance condition on $\frac{\partial \mathbf{g}\left(\mathbf{w}_{r} ; \boldsymbol{\theta}\right)}{\partial \boldsymbol{\theta}^{\prime}}$ ) for some neighborhood $\mathcal{N}$ of $\boldsymbol{\theta}_{0}$,

$$
\mathrm{E}\left[\sup _{\boldsymbol{\theta} \in \mathcal{N}}\left\|\frac{\partial \mathbf{g}\left(\mathbf{w}_{t} ; \boldsymbol{\theta}\right)}{\partial \boldsymbol{\theta}^{\prime}}\right\|\right]<\infty,
$$

so that for any consistent estimator $\tilde{\boldsymbol{\theta}}, \frac{1}{n} \sum_{t=1}^{n} \frac{\partial \mathrm{~g}\left(\mathrm{w}_{t} ; \tilde{\boldsymbol{\theta}}\right)}{\partial \boldsymbol{\theta}^{\prime}}(K \times p),{ }_{\mathrm{p}} \mathrm{E}\left[\frac{\partial \mathrm{g}\left(\mathrm{w}_{t} ; \theta_{0}\right)}{\partial \theta^{\prime}}\right]$,\\
(5) $\mathrm{E}\left[\frac{\partial \mathrm{g}\left(\mathrm{w}_{t} ; \theta_{0}\right)}{\partial \theta^{\prime}}\right]$ is of full column rank.

Then the following results hold:\\
(a) (asymptotic normality) $\hat{\boldsymbol{\theta}}$ is asymptotically normal with

$$
\underset{(p \times p)}{\operatorname{Avar}(\hat{\boldsymbol{\theta}})}=\left(\mathbf{G}^{\prime} \mathbf{W G}\right)^{-1} \mathbf{G}^{\prime} \mathbf{W S} \mathbf{W G}\left(\mathbf{G}^{\prime} \mathbf{W G}\right)^{-1},
$$

where $\underset{(K \times p)}{\mathbf{G}} \equiv \mathrm{E}\left[\frac{\partial \mathbf{g}\left(\mathbf{w}_{i} ; \boldsymbol{\theta}_{0}\right)}{\partial \boldsymbol{\theta}^{\prime}}\right]$,\\
(b) (consistent asymptotic variance estimation) Suppose there is available a consistent estimate $\widehat{\mathbf{S}}$ of $\mathbf{S}$. The asymptotic variance is consistently estimated by


\begin{gather*}
\widehat{\operatorname{Avar}(\hat{\boldsymbol{\theta}})}=\left(\widehat{\mathbf{G}}^{\prime} \widehat{\mathbf{W}} \widehat{\mathbf{G}}\right)^{-1} \widehat{\mathbf{G}}^{\prime} \widehat{\mathbf{W}} \widehat{\mathbf{S}} \widehat{\mathbf{W}} \widehat{\mathbf{G}}\left(\widehat{\mathbf{G}}^{\prime} \widehat{\mathbf{W}} \widehat{\mathbf{G}}\right)^{-1},  \tag{7.3.34}\\
\text { where } \underset{(K \times p)}{\widehat{\mathbf{G}}} \equiv \mathbf{G}_{n}(\hat{\boldsymbol{\theta}})=\frac{1}{n} \sum_{t=1}^{n} \frac{\partial \mathbf{g}\left(\mathbf{w}_{t}, \hat{\boldsymbol{\theta}}\right)}{\partial \boldsymbol{\theta}^{\prime}} .
\end{gather*}


The assumption that $\mathbf{S}$ is positive definite is not really needed for the conclusion of the proposition, but we might as well assume it here because in virtually all applications it is satisfied (or assumed) and also because it will be required in the discussion of hypothesis testing later in this chapter.

It is useful to note the analogy between linear and nonlinear GMM by comparing this proposition to Proposition 3.1. In linear GMM, $\mathbf{g}_{n}(\boldsymbol{\theta})$ is given by


\begin{equation*}
\mathbf{g}_{n}(\boldsymbol{\theta})=\left(\frac{1}{n} \sum_{t=1}^{n} \mathbf{x}_{t} \cdot y_{t}\right)-\left(\frac{1}{n} \sum_{t=1}^{n} \mathbf{x}_{t} \mathbf{z}_{t}^{\prime}\right) \boldsymbol{\theta} . \tag{7.3.35}
\end{equation*}


So $\widehat{\mathbf{G}} \equiv \mathbf{G}_{n}(\hat{\boldsymbol{\theta}})$ in Proposition 7.10 reduces to $-\left(\frac{1}{n} \sum \mathbf{X}_{t} \mathbf{Z}_{t}^{\prime}\right)$, and $\mathbf{G}$ reduces to $-E\left(\mathbf{x}_{t} \mathbf{z}_{t}^{\prime}\right)$. The thrust of Proposition 7.10 is that the asymptotic distribution of the nonlinear GMM estimator can be obtained by taking the linear approximation of $\mathbf{g}_{n}(\boldsymbol{\theta})$ around the true parameter value.

The analogy to linear GMM extends further:

\begin{itemize}
  \item (Efficient nonlinear GMM) Using the current notation, the matrix inequality (3.5.11) in Section 3.5 is
\end{itemize}


\begin{equation*}
\left(\mathbf{G}^{\prime} \mathbf{W G}\right)^{-1} \mathbf{G}^{\prime} \mathbf{W S} \mathbf{W G}\left(\mathbf{G}^{\prime} \mathbf{W G}\right)^{-1} \geq\left(\mathbf{G}^{\prime} \mathbf{S}^{-1} \mathbf{G}\right)^{-1}, \tag{7.3.36}
\end{equation*}


which says that the lower bound for $\operatorname{Avar}(\hat{\theta})$ is $\left(\mathbf{G}^{\prime} \mathbf{S}^{-1} \mathbf{G}\right)^{-1}$. The lower bound is achieved if $\widehat{\mathbf{W}}$ satisfies the efficiency condition that $\operatorname{plim}_{n \rightarrow \infty} \widehat{\mathbf{W}}=\mathbf{S}^{-1}$, namely, that $\widehat{\mathbf{W}}=\widehat{\mathbf{S}}^{-1}$ for some consistent estimator of $\mathbf{S}$.

\begin{itemize}
  \item ( $J$ statistic for testing overidentifying restrictions) Suppose that the conditions of Proposition 7.10 are satisfied and that $\widehat{\mathbf{W}}=\widehat{\mathbf{S}}^{-1}$ so that $\widehat{\mathbf{W}}$ satisfies the efficiency condition: $\operatorname{plim} \widehat{\mathbf{W}}=\mathbf{S}^{-1}$. Then the minimized distance $J \equiv n \mathbf{g}_{n}(\hat{\boldsymbol{\theta}})^{\prime} \widehat{\mathbf{S}}^{-1} \mathbf{g}_{n}(\hat{\boldsymbol{\theta}})\left(=-2 n Q_{n}(\hat{\boldsymbol{\theta}})\right)$ is asymptotically chi-squared with $K-p$ degrees of freedom. The proof of this result is essentially the same as in the linear case; see Newey and McFadden (1994, Section 9.5) for a proof.
  \item (Estimation of $\mathbf{S}$ ) $\mathbf{S}$ is the long-run variance of $\left\{\mathbf{g}\left(\mathbf{w}_{t} ; \boldsymbol{\theta}_{0}\right)\right\}$. Under some suitable technical conditions, the methods introduced in Section 6.6-such as the VARHAC - can be applied to the estimated series $\left\{\mathbf{g}\left(\mathbf{w}_{t} ; \hat{\boldsymbol{\theta}}\right)\right\}$ to obtain $\widehat{\mathbf{S}}$. In particular, if $\mathbf{w}_{t}$ is i.i.d., then $\mathbf{S}=\mathrm{E}\left[\mathbf{g}\left(\mathbf{w}_{t} ; \boldsymbol{\theta}_{0}\right) \mathbf{g}\left(\mathbf{w}_{t} ; \boldsymbol{\theta}_{0}\right)^{\prime}\right]$ and the sample second moment of $\mathbf{g}\left(\mathbf{w}_{i} ; \hat{\boldsymbol{\theta}}\right)$ can serve as $\widehat{\mathbf{S}}$.
\end{itemize}

\section*{GMM versus ML}
The question we have just answered is: taking the orthogonality conditions as given, what is the optimal choice of the weighting matrix $\mathbf{W}$ ? A related, but distinct, question is: what is the optimal choice of orthogonality conditions? ${ }^{15}$ This question, too, has a clear answer, provided that the parametric form of the density of the data $\left(\mathbf{w}_{1}, \ldots, \mathbf{w}_{n}\right)$ is known. Here, we focus on the i.i.d. case with the

\footnotetext{${ }^{15}$ Yet another efficiency question concerns the optimal choice of instruments in generalized nonlinear instrumental estimation a special case of GMM where the $\mathbf{g}$ function is a product of the vector of instruments and the error term for an estimation equation). See Newey and McFadden (1994, Section 5.4) for a discussion of optimal instruments.
}
density of $\mathbf{w}_{t}$ given by $f\left(\mathbf{w}_{t} ; \boldsymbol{\theta}\right)$. Let $\hat{\boldsymbol{\theta}}$ be the GMM estimator associated with the orthogonality conditions $\mathrm{E}\left[\mathbf{g}\left(\mathbf{w}_{t} ; \boldsymbol{\theta}_{0}\right)\right]=\mathbf{0}$. Its asymptotic variance $\operatorname{Avar}(\hat{\boldsymbol{\theta}})$ is given in Proposition 7.10. It is not hard to show that

\begin{equation*}
\operatorname{Avar}(\hat{\boldsymbol{\theta}}) \geq \mathrm{E}\left[\mathbf{s}\left(\mathbf{w}_{t} ; \boldsymbol{\theta}_{0}\right) \mathbf{s}\left(\mathbf{w}_{t} ; \boldsymbol{\theta}_{0}\right)^{\prime}\right]^{-1} \quad \text { where } \quad \mathbf{s}\left(\mathbf{w}_{t} ; \boldsymbol{\theta}_{0}\right) \equiv \frac{\partial \log f\left(\mathbf{w}_{t} ; \boldsymbol{\theta}_{0}\right)}{\partial \boldsymbol{\theta}} \tag{7.3.37}
\end{equation*}


That is, the inverse of the information matrix $\mathrm{E}\left(\mathbf{s s}^{\prime}\right)$ is the lower bound for the asymptotic variance of GMM estimators. This matrix inequality holds under the conditions of Proposition 7.10 plus some technical condition on $f\left(\mathbf{w}_{t} ; \boldsymbol{\theta}\right)$ which allows the interchange of differentiation and integration. See Newey and McFadden (1994, Theorem 5.1) for a statement of those conditions and a proof. Asymptotic efficiency of ML over GMM estimators then follows from this result because the lower bound $\left[\mathrm{E}\left(\mathbf{s s}^{\prime}\right)\right]^{-1}$ is the asymptotic variance of the ML estimator. The superiority of ML, however, is not very surprising, given that ML exploits the knowledge of the parametric form $f\left(\mathbf{w}_{i} ; \boldsymbol{\theta}\right)$ of the density function while GMM does not.

As you will show in Review Question 4 below, GMM achieves this lower bound when the $\mathbf{g}$ function in the orthogonality conditions is the score for observation $t$ :


\begin{equation*}
\underset{(K \times 1)}{\mathbf{g}\left(\mathbf{w}_{t} ; \boldsymbol{\theta}\right)}=\underset{(p \times 1)}{\mathbf{s}\left(\mathbf{w}_{t} ; \boldsymbol{\theta}\right)} \equiv \frac{\partial \log f\left(\mathbf{w}_{t} ; \boldsymbol{\theta}\right)}{\partial \boldsymbol{\theta}} . \tag{7.3.38}
\end{equation*}


Therefore, the GMM estimator with optimal orthogonality conditions is asymptotically equivalent to ML. Actually, they are numerically equivalent, which can be shown as follows. Since $K$ (the number of orthogonality conditions) equals $p$ (the number of parameters) when $\mathbf{g}$ is chosen optimally as in (7.3.38), the GMM estimator $\hat{\boldsymbol{\theta}}$ should satisfy $\mathbf{g}_{n}(\hat{\boldsymbol{\theta}})=\mathbf{0}$, which under (7.3.38) can be written as


\begin{equation*}
\frac{1}{n} \sum_{t=1}^{n} \mathbf{s}\left(\mathbf{w}_{t} ; \hat{\boldsymbol{\theta}}\right)=\mathbf{0} . \tag{7.3.39}
\end{equation*}


This is none other than the likelihood equations (i.e., the first-order condition for ML). They have at least one solution because the ML estimator satisfies them. If they have only one solution, then $\hat{\boldsymbol{\theta}}$ is the ML estimator as well as the GMM estimator. If they have more than one solution, we need to choose one from them as the GMM estimator. Since one of the solutions is the ML estimator, we can simply choose it as the GMM estimator.

\section*{Expressing the Sampling Error in a Common Format}
To prepare for the next section's discussion, it is useful to develop a slightly different proof of asymptotic normality. We have used the Mean Value Theorem to derive the mean value expansion for the sampling error, which is (7.3.7) for M-estimators and (7.3.32) for GMM. We show in this subsection that they can be written in a common format, (7.3.43) below.

Consider M-estimators first. Noting that $\frac{\partial Q_{n}\left(\theta_{0}\right)}{\partial \theta}=\frac{1}{n} \sum_{t=1}^{n} \mathbf{s}\left(\mathbf{w}_{t} ; \boldsymbol{\theta}_{0}\right)$, the mean value expansion (7.3.7) can be written as


\begin{equation*}
\sqrt{n}\left(\hat{\boldsymbol{\theta}}-\boldsymbol{\theta}_{0}\right)=-\left[\frac{1}{n} \sum_{t=1}^{n} \mathbf{H}\left(\mathbf{w}_{t} ; \overline{\boldsymbol{\theta}}\right)\right]^{-1} \sqrt{n} \frac{\partial Q_{n}\left(\boldsymbol{\theta}_{0}\right)}{\partial \boldsymbol{\theta}} \tag{7.3.40}
\end{equation*}


By condition (4) of Proposition 7.8, $\frac{1}{n} \sum_{t=1}^{n} \mathbf{H}\left(\mathbf{w}_{t} ; \overline{\boldsymbol{\theta}}\right)$ converges in probability to some $p \times p$ symmetric matrix $\boldsymbol{\Psi}$ given by


\begin{equation*}
\underset{(p \times p)}{\boldsymbol{\Psi}}=\mathrm{E}\left[\mathbf{H}\left(\mathbf{w}_{t} ; \boldsymbol{\theta}_{0}\right)\right] \tag{7.3.41}
\end{equation*}


So rewrite (7.3.40) as


\begin{align*}
\sqrt{n}\left(\hat{\boldsymbol{\theta}}-\boldsymbol{\theta}_{0}\right)= & -\boldsymbol{\Psi}^{-1} \sqrt{n} \frac{\partial Q_{n}\left(\boldsymbol{\theta}_{0}\right)}{\partial \boldsymbol{\theta}} \\
& -\left\{\left[\frac{1}{n} \sum_{t=1}^{n} \mathbf{H}\left(\mathbf{w}_{t} ; \overline{\boldsymbol{\theta}}\right)\right]^{-1}-\boldsymbol{\Psi}^{-1}\right\} \sqrt{n} \frac{\partial Q_{n}\left(\boldsymbol{\theta}_{0}\right)}{\partial \boldsymbol{\theta}} . \tag{7.3.42}
\end{align*}


By construction, the term in braces converges in probability to zero. By condition (3) of Proposition 7.8, $\sqrt{n} \frac{\partial Q_{n}\left(\boldsymbol{\theta}_{0}\right)}{\partial \boldsymbol{\theta}}\left(=\frac{1}{\sqrt{n}} \sum_{t=1}^{n} \mathbf{s}\left(\mathbf{w}_{t} ; \boldsymbol{\theta}_{0}\right)\right)$ converges to a random variable. So the last term, which is the product of the term in braces and $\sqrt{n} \frac{\partial Q_{n}\left(\boldsymbol{\theta}_{0}\right)}{\partial \boldsymbol{\theta}}$, converges to zero in probability by Lemma 2.4(b). This fact can be written compactly as a Taylor expansion: ${ }^{16}$


\begin{equation*}
\sqrt{n}\left(\hat{\boldsymbol{\theta}}-\boldsymbol{\theta}_{0}\right)=-\underset{(p \times 1)}{\boldsymbol{\Psi}^{-1}} \sqrt{n \times p)} \frac{\partial Q_{n}\left(\boldsymbol{\theta}_{0}\right)}{\partial(p \times 1)}+\underset{(p \times 1)}{o_{p}} \tag{7.3.43}
\end{equation*}


where the term " $o_{p}$ " means "some random variable that converges to zero in probability." ${ }^{17}$ The exact expression for the $o_{p}$ term will depend on the context. Here, it equals the last term in (7.3.42). As you will see, all that matters for our argument

\footnotetext{${ }^{16}$ It is apt to call this equation a Taylor expansion rather than a mean value expansion because the matrix that multiplies $\sqrt{n} \frac{\partial Q_{n}\left(\theta_{0}\right)}{\partial \theta}$ is evaluated at $\theta_{0}$ rather than at $\bar{\theta}$ as in the mean value expansion.\\
${ }^{17}$ The " $o_{p}$ " notation was introduced in Section 2.1.
}
is that the $o_{p}$ term vanishes (converges to zero in probability), not its expression. Since the difference between $\sqrt{n}\left(\hat{\boldsymbol{\theta}}-\boldsymbol{\theta}_{0}\right)$ and $-\boldsymbol{\Psi}^{-1} \sqrt{n} \frac{\partial Q_{n}\left(\boldsymbol{\theta}_{0}\right)}{\partial \boldsymbol{\theta}}$ vanishes, the asymptotic distribution of $\sqrt{n}\left(\hat{\boldsymbol{\theta}}-\boldsymbol{\theta}_{0}\right)$ is the same as that of $-\boldsymbol{\Psi}^{-1} \sqrt{n} \frac{\partial Q_{n}\left(\boldsymbol{\theta}_{0}\right)}{\partial \boldsymbol{\theta}}$ by Lemma 2.4(a). It then follows that $\sqrt{n}\left(\hat{\boldsymbol{\theta}}-\boldsymbol{\theta}_{0}\right)$ converges to a normal distribution with mean zero and asymptotic variance given by

\begin{equation*}
\operatorname{Avar}(\hat{\boldsymbol{\theta}})=\boldsymbol{\Psi}^{-1} \boldsymbol{\Sigma} \boldsymbol{\Psi}^{-1} \text { where } \underset{(p \times p)}{\boldsymbol{\Sigma}} \equiv \operatorname{Avar}\left(\frac{\partial Q_{n}\left(\boldsymbol{\theta}_{0}\right)}{\partial \boldsymbol{\theta}}\right) \tag{7.3.44}
\end{equation*}


For M-estimators, $\sqrt{n} \frac{\partial Q_{n}\left(\boldsymbol{\theta}_{0}\right)}{\partial \boldsymbol{\theta}}=\frac{1}{\sqrt{n}} \sum_{t=1}^{n} \mathbf{s}\left(\mathbf{w}_{t} ; \boldsymbol{\theta}_{0}\right)$ and $\boldsymbol{\Sigma}$ is the long-run variance of $\mathbf{s}\left(\mathbf{w}_{t} ; \boldsymbol{\theta}_{0}\right)$. Setting $\boldsymbol{\Psi}=\mathrm{E}\left[\mathbf{H}\left(\mathbf{w}_{t} ; \boldsymbol{\theta}_{0}\right)\right]$ gives the expression for $\operatorname{Avar}(\hat{\boldsymbol{\theta}})$ in Proposition 7.8.

Next turn to GMM. The expression for $\sqrt{n} \frac{\partial Q_{n}\left(\boldsymbol{\theta}_{0}\right)}{\partial \boldsymbol{\theta}}$ is


\begin{equation*}
\sqrt{n} \frac{\partial Q_{n}\left(\boldsymbol{\theta}_{0}\right)}{\partial \boldsymbol{\theta}}=-\left[\mathbf{G}_{n}\left(\boldsymbol{\theta}_{0}\right)\right]^{\prime} \widehat{\mathbf{W}} \frac{1}{\sqrt{n}} \sum_{t=1}^{n} \mathbf{g}\left(\mathbf{w}_{t} ; \boldsymbol{\theta}_{0}\right) . \tag{7.3.45}
\end{equation*}


(Recall that $\mathbf{G}_{n}(\boldsymbol{\theta}) \equiv \frac{\partial \mathbf{g}_{n}(\boldsymbol{\theta})}{\partial \boldsymbol{\theta}^{\prime}}$.) Since $\mathbf{G}_{n}\left(\boldsymbol{\theta}_{0}\right)=\frac{1}{n} \sum \frac{\partial \mathbf{g}\left(\mathbf{w}_{i} ; \boldsymbol{\theta}_{0}\right)}{\partial \boldsymbol{\theta}^{\prime}} \rightarrow_{\mathbf{p}} \mathbf{G}\left(\equiv \mathrm{E}\left[\frac{\partial \mathbf{g}\left(\mathbf{w}_{r} ; \boldsymbol{\theta}_{0}\right)}{\partial \boldsymbol{\theta}^{\prime}}\right]\right)$, $\widehat{\mathbf{W}} \rightarrow_{\mathrm{p}} \mathbf{W}$, and $\frac{1}{\sqrt{n}} \sum_{t=1}^{n} \mathbf{g}\left(\mathbf{w}_{t} ; \boldsymbol{\theta}_{0}\right) \rightarrow_{\mathrm{d}} N(\mathbf{0}, \mathbf{S})$ under the conditions of Proposition 7.10, we have $\sqrt{n} \frac{\partial Q_{n}\left(\boldsymbol{\theta}_{0}\right)}{\partial \boldsymbol{\theta}}$ converging to a normal distribution with mean zero and


\begin{equation*}
\underset{(p \times p)}{\boldsymbol{\Sigma}} \equiv \operatorname{Avar}\left(\frac{\partial Q_{n}\left(\boldsymbol{\theta}_{0}\right)}{\partial \boldsymbol{\theta}}\right)=\underset{(p \times K)(K \times K)(K \times K)(K \times K)(K \times p)}{\mathbf{G}} \underset{(K)}{\mathbf{W}} \underset{(K)}{\mathbf{G}} . \tag{7.3.46}
\end{equation*}


Now rewrite the mean value expansion for the sampling error for GMM (7.3.32) as


\begin{gather*}
\sqrt{n}\left(\hat{\boldsymbol{\theta}}-\boldsymbol{\theta}_{0}\right)=-\mathbf{B}^{-1} \mathbf{c}, \quad \mathbf{B} \equiv-\mathbf{G}_{n}(\hat{\boldsymbol{\theta}})^{\prime} \widehat{\mathbf{W}} \mathbf{G}_{n}(\overline{\boldsymbol{\theta}}), \\
\mathbf{c} \equiv-\mathbf{G}_{n}(\hat{\boldsymbol{\theta}})^{\prime} \widehat{\mathbf{W}} \frac{1}{\sqrt{n}} \sum_{t=1}^{n} \mathbf{g}\left(\mathbf{w}_{t} ; \boldsymbol{\theta}_{0}\right) . \tag{7.3.47}
\end{gather*}


It is left as Review Question 5 below to show that this can be written as the Taylor expansion (7.3.43), with the symmetric matrix $\boldsymbol{\Psi}$ given by


\begin{equation*}
\underset{(p \times p)}{\boldsymbol{\Psi}}=-\underset{(p \times K)}{\mathbf{G}^{\prime}} \underset{(K \times K)}{\mathbf{W}} \underset{(K \times p)}{\mathbf{G}} . \tag{7.3.48}
\end{equation*}


Substituting (7.3.48) and (7.3.46) into (7.3.44), we obtain the GMM asymptotic variance indicated in Proposition 7.10.

Table 7.1 summarizes the discussion of this subsection by indicating how the substitution should be done to make the common formulas applicable to M-estimators and GMM.

\begin{table}[h]
\begin{center}
\captionsetup{labelformat=empty}
\caption{Table 7.1: Taylor Expansion for the Sampling Error}
\begin{tabular}{|l|l|l|}
\hline
\multicolumn{3}{|c|}{$\sqrt{n}\left(\hat{\boldsymbol{\theta}}-\boldsymbol{\theta}_{0}\right)=-\boldsymbol{\Psi}^{-1} \sqrt{n} \frac{\partial Q_{n}\left(\boldsymbol{\theta}_{0}\right)}{\partial \boldsymbol{\theta}}+o_{p}, \sqrt{n} \frac{\partial Q_{n}\left(\boldsymbol{\theta}_{0}\right)}{\partial \boldsymbol{\theta}} \underset{\mathrm{d}}{\rightarrow} N(\mathbf{0}, \boldsymbol{\Sigma}), \operatorname{Avar}(\hat{\boldsymbol{\theta}})=\boldsymbol{\Psi}^{-1} \boldsymbol{\Sigma} \boldsymbol{\Psi}^{-1}$} \\
\hline
Terms for substitution & M-estimators & GMM \\
\hline
$Q_{n}(\boldsymbol{\theta})$ & $\frac{1}{n} \sum_{t=1}^{n} m\left(\mathbf{w}_{t} ; \boldsymbol{\theta}\right)$ & $-\frac{1}{2} \mathbf{g}_{n}(\boldsymbol{\theta})^{\prime} \widehat{\mathbf{W}} \mathbf{g}_{n}(\boldsymbol{\theta})$ \\
\hline
$\sqrt{n} \frac{\partial Q_{n}\left(\boldsymbol{\theta}_{0}\right)}{\boldsymbol{\partial} \boldsymbol{\theta}}$ & $\frac{1}{\sqrt{n}} \sum_{t=1}^{n} \mathbf{s}\left(\mathbf{w}_{t} ; \boldsymbol{\theta}_{0}\right)$ & $-\left[\mathbf{G}_{n}\left(\boldsymbol{\theta}_{0}\right)\right]^{\prime} \widehat{\mathbf{W}} \frac{1}{\sqrt{n}} \sum_{t=1}^{n} \mathbf{g}\left(\mathbf{w}_{t} ; \boldsymbol{\theta}_{0}\right)$ \\
\hline
$\Psi$ & $\mathrm{E}\left[\mathbf{H}\left(\mathbf{w}_{t} ; \boldsymbol{\theta}_{0}\right)\right]$ & -G' WG \\
\hline
$\boldsymbol{\Sigma}$ & long-run variance of $\mathbf{s}\left(\mathbf{w}_{t} ; \boldsymbol{\theta}_{0}\right)$ & G' WS WG \\
\hline
\end{tabular}
\end{center}
\end{table}

Note: See (7.3.2) and (7.3.3) for definition of $\mathbf{s}\left(\mathbf{w}_{t} ; \boldsymbol{\theta}\right)$ and $\mathbf{H}\left(\mathbf{w}_{t} ; \boldsymbol{\theta}\right)$. Also,

$$
\mathbf{g}_{n}(\boldsymbol{\theta}) \equiv \frac{1}{n} \sum_{t=1}^{n} \mathbf{g}\left(\mathbf{w}_{t} ; \boldsymbol{\theta}\right) \quad \text { and } \quad \mathbf{G} \equiv \mathrm{E}\left[\frac{\partial \mathbf{g}\left(\mathbf{w}_{t} ; \boldsymbol{\theta}_{0}\right)}{\partial \boldsymbol{\theta}^{\prime}}\right] .
$$

In passing, it is useful for the next section's discussion to note that $\mathbf{\Sigma}=-\Psi$ for ML with i.i.d. observations and efficient GMM (the GMM with $\mathbf{W}=\mathbf{S}^{-1}$ ). For $\mathrm{ML}, \boldsymbol{\Sigma}=\mathrm{E}\left[\mathbf{s}\left(\mathbf{w}_{t} ; \boldsymbol{\theta}_{0}\right) \mathbf{s}\left(\mathbf{w}_{t} ; \boldsymbol{\theta}_{0}\right)^{\prime}\right]$, which equals the negative of the $\boldsymbol{\Psi}$ given in (7.3.41) thanks to the information matrix equality. For efficient GMM, setting $\mathbf{W}=\mathbf{S}^{-1}$ in (7.3.46) gives


\begin{equation*}
\boldsymbol{\Sigma}=\mathbf{G}^{\prime} \mathbf{S}^{-1} \mathbf{G}, \tag{7.3.49}
\end{equation*}


which is the negative of the $\boldsymbol{\Psi}$ in (7.3.48).

\section*{QUESTIONS FOR REVIEW}
\begin{enumerate}
  \item (Score and Hessian of objective function) For M-estimators and ML in particular, we defined the score for observation $t, \mathbf{s}\left(\mathbf{w}_{t} ; \boldsymbol{\theta}\right)$, to be the gradient (vector of first derivatives) of $m\left(\mathbf{w}_{t} ; \boldsymbol{\theta}\right)$ and the Hessian for observation $t, \mathbf{H}\left(\mathbf{w}_{t} ; \boldsymbol{\theta}\right)$, to be the Hessian of $m\left(\mathbf{w}_{t} ; \boldsymbol{\theta}\right)$. Let the score (without the qualifier "for observation $t^{\prime \prime}$ ) be the gradient of the objective function $Q_{n}(\boldsymbol{\theta})$ and denote it as $\mathbf{s}_{n}(\boldsymbol{\theta})$. Let the Hessian (without the qualifier "for observation $t$ ") be the Hessian of $Q_{n}(\boldsymbol{\theta})$ and denote it as $\mathbf{H}_{n}(\boldsymbol{\theta})$. Under the conditions of Proposition 7.9, show that
\end{enumerate}

$$
\mathrm{E}\left[\mathbf{s}_{n}\left(\boldsymbol{\theta}_{0}\right)\right]=\mathbf{0},-\frac{1}{n} \mathrm{E}\left[\mathbf{H}_{n}\left(\boldsymbol{\theta}_{0}\right)\right]=\mathrm{E}\left[\mathbf{s}_{n}\left(\boldsymbol{\theta}_{0}\right) \mathbf{s}_{n}\left(\boldsymbol{\theta}_{0}\right)^{\prime}\right] .
$$

Hint: $\left\{\mathbf{s}\left(\mathbf{w}_{t} ; \boldsymbol{\theta}_{0}\right)\right\}$ is i.i.d.\\
2. (Conditional information matrix equality for the linear regression model) For the linear regression model of Example 7.10, verify that

$$
\mathrm{E}\left[\mathbf{s}\left(\mathbf{w}_{t} ; \boldsymbol{\theta}_{0}\right) \mid \mathbf{x}_{t}\right]=\mathbf{0} \text { and }-\mathrm{E}\left[\mathbf{H}\left(\mathbf{w}_{t} ; \boldsymbol{\theta}_{0}\right) \mid \mathbf{x}_{t}\right]=\mathrm{E}\left[\mathbf{s}\left(\mathbf{w}_{t} ; \boldsymbol{\theta}_{0}\right) \mathbf{s}\left(\mathbf{w}_{t} ; \boldsymbol{\theta}_{0}\right)^{\prime} \mid \mathbf{x}_{t}\right] .
$$

\begin{enumerate}
  \setcounter{enumi}{2}
  \item (Avar for the linear regression model) For the linear regression model of Example 7.10, let $\left(\widehat{\boldsymbol{\beta}}, \hat{\sigma}^{2}\right)$ be the ML estimate of $\left(\boldsymbol{\beta}, \sigma^{2}\right)$. Write down the first estimator of $\operatorname{Avar}(\hat{\boldsymbol{\theta}})$ defined in (7.3.14). Does $\widehat{\operatorname{Avar}(\widehat{\boldsymbol{\beta}})}$ equal $\hat{\sigma}^{2}\left(\frac{1}{n} \sum \mathbf{x}_{t} \mathbf{x}_{t}^{\prime}\right)^{-1}$ ? [Answer: No.]
  \item (GMM with optimal orthogonality conditions) Assume that the relevant technical conditions are satisfied for the hypothetical density $f\left(\mathbf{w}_{i} ; \boldsymbol{\theta}\right)$ so that the information matrix equality holds. Using Proposition 7.10, show that when the $\mathbf{g}$ function in the orthogonality conditions is given as in (7.3.38), the asymptotic variance matrix of the GMM estimator equals the inverse of the information matrix for observation $t$ (which is the lower bound in (7.3.37)). Hint: Since $K$ equals $p, \mathbf{G}$ is a square matrix and the expression for $\operatorname{Avar}(\hat{\boldsymbol{\theta}})$ in
\end{enumerate}

Proposition 7.10 reduces to $\operatorname{Avar}(\hat{\boldsymbol{\theta}})=\mathbf{G}^{-1} \mathbf{S G}^{-1}$. The choice of $\mathbf{g}$ implies $\mathbf{G}=\mathrm{E}\left[\mathbf{H}\left(\mathbf{w}_{t} ; \boldsymbol{\theta}_{0}\right)\right]$ and $\mathbf{S}=\mathrm{E}\left[\mathbf{s}\left(\mathbf{w}_{t} ; \boldsymbol{\theta}_{0}\right) \mathbf{s}\left(\mathbf{w}_{t} ; \boldsymbol{\theta}_{0}\right)^{\prime}\right]$.\\
5. (Taylor expansion of the sampling error for GMM) Show that (7.3.47) can be written as the Taylor expansion (7.3.43) under the conditions of Proposition 7.10 , by taking the following steps.\\
(a) Show that $\mathbf{B}^{-1}=\Psi^{-1}+\mathbf{Y}_{n}$ where $\mathbf{Y}_{n} \rightarrow_{\mathrm{p}} \mathbf{0}$.\\
(b) Show that $\mathbf{c}=\sqrt{n} \frac{\partial Q_{n}\left(\theta_{0}\right)}{\partial \theta}+\mathbf{y}_{n}$ where $\mathbf{y}_{n} \rightarrow_{\mathrm{p}}$ 0. Hint: $\mathbf{G}_{n}(\hat{\boldsymbol{\theta}})-\mathbf{G}_{n}\left(\boldsymbol{\theta}_{0}\right)= \left(\mathbf{G}_{n}(\hat{\boldsymbol{\theta}})-\mathbf{G}\right)-\left(\mathbf{G}_{n}\left(\boldsymbol{\theta}_{0}\right)-\mathbf{G}\right) \rightarrow_{\mathrm{p}} \mathbf{0}$ and $\widehat{\mathbf{W}} \frac{1}{\sqrt{n}} \sum_{t=1}^{n} \mathbf{g}\left(\mathbf{w}_{t} ; \boldsymbol{\theta}_{0}\right)$ converges in distribution to a random variable.\\
(c) Show that $-\mathbf{B}^{-1} \mathbf{c}=-\Psi^{-1} \sqrt{n} \frac{\partial Q_{n}\left(\theta_{0}\right)}{\partial \theta}+\mathbf{x}_{n}$, where $\mathbf{x}_{n} \equiv-\left[\mathbf{Y}_{n} \sqrt{n} \frac{\partial Q_{n}\left(\theta_{0}\right)}{\partial \theta}+\right. \left.\boldsymbol{\Psi}^{-1} \mathbf{y}_{n}+\mathbf{Y}_{n} \mathbf{y}_{n}\right]$.\\
(d) Show that $\mathbf{x}_{n} \rightarrow_{\mathrm{p}} \mathbf{0}$.\\
6. (Why the GMM distance is deflated by 2 ) Verify that, if the negative of the distance $\mathbf{g}_{n}(\boldsymbol{\theta})^{\prime} \widehat{\mathbf{S}}^{-1} \mathbf{g}_{n}(\boldsymbol{\theta})$ were not divided by 2 in the definition of the efficient GMM objective function, then we would not have $\boldsymbol{\Sigma}=-\boldsymbol{\Psi}$.

\subsection*{7.4 Hypothesis Testing}
As is well known for ML, there is a trio of statistics - the Wald, Lagrange multiplier (LM), and likelihood ratio (LR) statistics - that can be used for testing the null hypothesis. The three statistics are asymptotically equivalent in that they share the same asymptotic distribution (of $\chi^{2}$ ). In this section we show that the asymptotic equivalence of the trinity can be extended to GMM by developing an argument applicable to both ML and GMM. The key to the argument is the common format for the sampling error derived at the end of the previous section. On the first reading, you may wish to read the next subsection and then jump to the last subsection.

\section*{The Null Hypothesis}
We have already considered, in the context of the linear regression model, the problem of testing a set of $r$ possibly nonlinear restrictions (see Section 2.4). To recapitulate, let $\theta_{0}$ be the $p$-dimensional model parameter. The null hypothesis can\\
be expressed as


\begin{equation*}
\mathrm{H}_{0}: \underset{(r \times 1)}{\mathbf{a}\left(\boldsymbol{\theta}_{0}\right)}=\mathbf{0} . \tag{7.4.1}
\end{equation*}


We assume that $\mathbf{a}(\cdot)$ is continuously differentiable. Also, let


\begin{equation*}
\underset{(r \times p)}{\mathbf{A}(\boldsymbol{\theta})} \equiv \frac{\partial \mathbf{a}(\boldsymbol{\theta})}{\partial \boldsymbol{\theta}^{\prime}} \tag{7.4.2}
\end{equation*}


be the Jacobian of $\mathbf{a}(\boldsymbol{\theta})$. We assume that


\begin{equation*}
\underset{(r \times p)}{\mathbf{A}_{0}} \equiv \mathbf{A}\left(\boldsymbol{\theta}_{0}\right) \text { is of full row rank. } \tag{7.4.3}
\end{equation*}


This ensures that the $r$ restrictions are not redundant. This rank condition implies that $r \leq p$.

Let $\hat{\boldsymbol{\theta}}$ be the extremum estimator in question. It is either ML or GMM. It solves the unconstrained optimization problem (7.1.1). The Wald statistic for testing the null hypothesis utilizes the unconstrained estimator. The LM statistic, in contrast, utilizes the constrained estimator, denoted $\tilde{\theta}$, which solves


\begin{equation*}
\max _{\boldsymbol{\theta} \in \Theta} Q_{n}(\boldsymbol{\theta}) \text { s.t. } \mathbf{a}(\boldsymbol{\theta})=\mathbf{0} \text {. } \tag{7.4.4}
\end{equation*}


As was shown in Section 7.2, the true parameter value $\boldsymbol{\theta}_{0}$ solves the "limit unconstrained problem" where $Q_{n}$ in the unconstrained optimization problem (7.1.1) is replaced by the limit function $Q_{0}$. It also solves the "limit constrained problem" where $Q_{n}$ in the above constrained optimization problem is replaced by $Q_{0}$, because $\boldsymbol{\theta}_{0}$ satisfies the constraint. It turns out that the uniform convergence in probability of $Q_{n}(\cdot)$ to $Q_{0}(\cdot)$ assures that the limit of the constrained estimator $\tilde{\boldsymbol{\theta}}$ is the solution to the limit constrained problem, which is $\boldsymbol{\theta}_{0}$. For a proof, see Newey and McFadden (1994, Theorem 9.1) for GMM and Gallant and White (1988, Lemma 7.3(b)) for general extremum estimators. It is also possible to show that, if all the conditions ensuring the asymptotic normality of the unconstrained estimator are satisfied, then $\tilde{\boldsymbol{\theta}}$, too, is asymptotically normal under the null. For a proof for GMM, see Newey and McFadden (1994, Theorem 9.2); for ML, see, e.g., Amemiya (1985, Section 4.5). ${ }^{18}$

\footnotetext{${ }^{18}$ If you are interested in the proof, it is just the mean value version of the discussion in the subsection on the LM statistic below, which is based on Taylor expansions.
}\section*{The Working Assumptions}
The argument for showing that the Wald, LM, and LR statistics are all asymptotically $\chi^{2}(r)$ is based on the truth of the following conditions.\\
(A) $\sqrt{n}\left(\hat{\boldsymbol{\theta}}-\boldsymbol{\theta}_{0}\right)$ admits the Taylor expansion (7.3.43), which is reproduced here


\begin{equation*}
\sqrt{n}\left(\hat{\boldsymbol{\theta}}-\boldsymbol{\theta}_{0}\right)=-\boldsymbol{\Psi}^{-1} \sqrt{n} \frac{\partial Q_{n}\left(\boldsymbol{\theta}_{0}\right)}{\partial \boldsymbol{\theta}}+o_{p} . \tag{7.4.5}
\end{equation*}


(B) $\sqrt{n} \frac{\partial Q_{n}\left(\boldsymbol{\theta}_{0}\right)}{\partial \boldsymbol{\theta}} \rightarrow_{\mathrm{d}} N(\mathbf{0}, \boldsymbol{\Sigma}), \boldsymbol{\Sigma}$ positive definite.\\
(C) $\sqrt{n}\left(\tilde{\boldsymbol{\theta}}-\boldsymbol{\theta}_{0}\right)$ converges in distribution to a random variable.\\
(D) $\boldsymbol{\Sigma}=-\boldsymbol{\Psi}$.

We have shown in the previous section that the first two conditions are satisfied for conditional ML under the hypothesis of Proposition 7.9, for unconditional ML under the hypothesis suitably modified, and for GMM under the hypothesis of Proposition 7.10. As just noted, under the same hypothesis, the constrained estimator $\tilde{\boldsymbol{\theta}}$ is consistent and asymptotically normal, ensuring condition (C). So conditions (A)-(C) are implied by the hypothesis assuring the asymptotic normality of the extremum estimator.

The Wald and LM statistics, if defined appropriately so as not to depend on condition (D), are asymptotically $\chi^{2}(r)$. But their expressions can be simplified if condition (D) is invoked. Furthermore, without the condition, the LR statistic would not be asymptotically $\chi^{2}$. As was noted at the end of the previous section, condition (D) is satisfied for ML and also for efficient GMM with $\mathbf{W}=\mathbf{S}^{-1}$. Therefore, the part of the discussion that relies on condition (D) is not applicable to nonefficient GMM.

\section*{The Wald Statistic}
Throughout the book we have had several occasions to derive the Wald statistic using the "delta method" of Lemma 2.5. Here, we provide a slightly different derivation based on the Taylor expansion. Applying the Mean Value Theorem to $\mathbf{a}(\hat{\boldsymbol{\theta}})$ around $\boldsymbol{\theta}_{0}$ and noting that $\mathbf{a}\left(\boldsymbol{\theta}_{0}\right)=\mathbf{0}$ under the null, the mean value expansion of $\sqrt{n} \mathbf{a}(\hat{\boldsymbol{\theta}})$ is


\begin{equation*}
\underset{(r \times 1)}{\sqrt{n}} \mathbf{a}(\hat{\boldsymbol{\theta}})=\underset{(r \times p)}{\mathbf{A}(\overline{\boldsymbol{\theta}})} \underset{(p \times 1)}{\sqrt{n}\left(\hat{\boldsymbol{\theta}}-\boldsymbol{\theta}_{0}\right)} . \tag{7.4.6}
\end{equation*}


Derivation of the Taylor expansion from the mean value expansion is more transparent than in the previous section and proceeds as follows. Since $\overline{\boldsymbol{\theta}}$, lying between\\
$\boldsymbol{\theta}_{0}$ and $\hat{\boldsymbol{\theta}}$, is consistent and since $\mathbf{A}(\cdot)$ is continuous, $\mathbf{A}(\overline{\boldsymbol{\theta}})$ converges in probability to $\mathbf{A}_{0}\left(\equiv \mathbf{A}\left(\boldsymbol{\theta}_{0}\right)\right)$. Furthermore, the term multiplying $\mathbf{A}(\overline{\boldsymbol{\theta}}), \sqrt{n}\left(\hat{\boldsymbol{\theta}}-\boldsymbol{\theta}_{0}\right)$, converges to a random variable. So $\left(\mathbf{A}(\overline{\boldsymbol{\theta}})-\mathbf{A}_{0}\right) \sqrt{n}\left(\hat{\boldsymbol{\theta}}-\boldsymbol{\theta}_{0}\right)$ vanishes (converges to zero in probability) by Lemma 2.4(b). Denoting this term by $o_{p}$, the mean value expansion can be rewritten as a Taylor expansion:


\begin{equation*}
\underset{(r \times 1)}{\sqrt{n}} \mathbf{a}(\hat{\boldsymbol{\theta}})=\underset{(r \times p)}{\mathbf{A}_{0}} \sqrt{n} \underset{(p \times 1)}{\left(\hat{\boldsymbol{\theta}}-\boldsymbol{\theta}_{0}\right)}+\underset{(r \times 1)}{\boldsymbol{o}_{p}} \tag{7.4.7}
\end{equation*}


(Note that this argument would not be valid if $\sqrt{n}\left(\hat{\boldsymbol{\theta}}-\boldsymbol{\theta}_{0}\right)$ did not converge to a random variable.)

Substituting the Taylor expansion (7.4.5) into this, we obtain an equation linking $\sqrt{n} \mathbf{a}(\hat{\boldsymbol{\theta}})$ to $\sqrt{n} \frac{\partial Q_{n}\left(\boldsymbol{\theta}_{0}\right)}{\partial \boldsymbol{\theta}}$ :


\begin{equation*}
\underset{(r \times 1)}{\sqrt{n}} \mathbf{a}(\hat{\boldsymbol{\theta}})=-\underset{(r \times p)}{\mathbf{A}_{0}} \underset{(p \times p)}{\mathbf{\Psi}^{-1}} \sqrt{n} \frac{\partial Q_{n}\left(\boldsymbol{\theta}_{0}\right)}{\partial \boldsymbol{\theta}}+o_{p} \tag{7.4.8}
\end{equation*}


(The $o_{p}$ here is $\mathbf{A}_{0}$ times the $o_{p}$ in (7.4.5) plus the $o_{p}$ in (7.4.7), but it still converges to zero in probability.) It immediately follows from this expression and the asymptotic normality of $\sqrt{n} \frac{\partial Q_{n}\left(\boldsymbol{\theta}_{0}\right)}{\partial \boldsymbol{\theta}}$ (condition (B)) that $\sqrt{n} \mathbf{a}(\hat{\boldsymbol{\theta}})$ converges to a normal distribution with mean zero and


\begin{align*}
\underset{(r \times r)}{\operatorname{Avar}(\mathbf{a}(\hat{\boldsymbol{\theta}}))} & =\mathbf{A}_{0} \boldsymbol{\Psi}^{-1} \boldsymbol{\Sigma} \boldsymbol{\Psi}^{-1} \mathbf{A}_{0}^{\prime} \\
& =\underset{(r \times p)}{\mathbf{A}_{0}} \underset{(p \times p)}{\boldsymbol{\Sigma}^{-1}} \underset{(p \times r)}{\mathbf{A}_{0}^{\prime}} \quad \text { (if condition (D) is invoked) } \tag{7.4.9}
\end{align*}


Since the $r \times p$ matrix $\mathbf{A}_{0}$ is of full row rank and $\boldsymbol{\Sigma}$ is positive definite (by condition (B)), this $r \times r$ matrix is positive definite. Therefore, the associated quadratic form


\begin{equation*}
n \mathbf{a}(\hat{\boldsymbol{\theta}})^{\prime}\left[\mathbf{A}_{0} \boldsymbol{\Sigma}^{-1} \mathbf{A}_{0}^{\prime}\right]^{-1} \mathbf{a}(\hat{\boldsymbol{\theta}}) \tag{7.4.10}
\end{equation*}


is asymptotically $\chi^{2}(r)$ under the null. As noted above, $\mathbf{A}(\hat{\boldsymbol{\theta}}) \rightarrow_{\mathrm{p}} \mathbf{A}_{0}$. So if there is available a consistent estimator $\widehat{\boldsymbol{\Sigma}}$ of $\boldsymbol{\Sigma}$, then $\mathbf{A}(\hat{\boldsymbol{\theta}}) \widehat{\boldsymbol{\Sigma}}^{-1} \mathbf{A}(\hat{\boldsymbol{\theta}})^{\prime}$ is consistent for $\mathbf{A}_{0} \boldsymbol{\Sigma}^{-1} \mathbf{A}_{0}^{\prime}$, which implies by Lemma 2.4(d) that the Wald statistic defined by


\begin{equation*}
W \equiv n \mathbf{a}(\hat{\boldsymbol{\theta}})^{\prime}\left[\mathbf{A}(\hat{\boldsymbol{\theta}}) \widehat{\mathbf{\Sigma}}^{-1} \mathbf{A}(\hat{\boldsymbol{\theta}})^{\prime}\right]^{-1} \mathbf{a}(\hat{\boldsymbol{\theta}}) \tag{7.4.11}
\end{equation*}


too is asymptotically $\chi^{2}(r)$ under the null.\\
How the consistent estimator $\widehat{\boldsymbol{\Sigma}}$ of $\boldsymbol{\Sigma}(=-\boldsymbol{\Psi})$ is obtained depends on whether the extremum estimator is ML or efficient GMM. For ML, it is given by a consistent\\
estimate of $-\boldsymbol{\Psi}=-\mathrm{E}\left[\mathbf{H}\left(\mathbf{w}_{i} ; \boldsymbol{\theta}_{0}\right)\right]$ :


\begin{equation*}
-\frac{\partial^{2} Q_{n}(\hat{\boldsymbol{\theta}})}{\partial \boldsymbol{\theta} \partial \boldsymbol{\theta}^{\prime}}=-\frac{1}{n} \sum_{t=1}^{n} \mathbf{H}\left(\mathbf{w}_{t} ; \hat{\boldsymbol{\theta}}\right), \tag{7.4.12}
\end{equation*}


or by a consistent estimate of $\boldsymbol{\Sigma}=\mathrm{E}\left[\mathbf{s}\left(\mathbf{w}_{t} ; \boldsymbol{\theta}_{0}\right) \mathbf{s}\left(\mathbf{w}_{t} ; \boldsymbol{\theta}_{0}\right)^{\prime}\right]$ :


\begin{equation*}
\frac{1}{n} \sum \mathbf{s}\left(\mathbf{w}_{t} ; \hat{\boldsymbol{\theta}}\right) \mathbf{s}\left(\mathbf{w}_{t} ; \hat{\boldsymbol{\theta}}\right)^{\prime} \tag{7.4.13}
\end{equation*}


For efficient GMM, $\boldsymbol{\Sigma}=\mathbf{G}^{\prime} \mathbf{S}^{-1} \mathbf{G}$ (see (7.3.49)), which is consistently estimated by


\begin{equation*}
\widehat{\mathbf{\Sigma}}=\widehat{\mathbf{G}}^{\prime} \widehat{\mathbf{S}}^{-1} \widehat{\mathbf{G}} \quad \text { with } \underset{(K \times p)}{\widehat{\mathbf{G}}} \equiv \mathbf{G}_{n}(\hat{\boldsymbol{\theta}})=\frac{1}{n} \sum_{t=1}^{n} \frac{\partial \mathbf{g}\left(\mathbf{w}_{t} ; \hat{\boldsymbol{\theta}}\right)}{\partial \boldsymbol{\theta}^{\prime}} \tag{7.4.14}
\end{equation*}


As already mentioned, $\widehat{\mathbf{S}}$ is an estimate of the long-run variance matrix constructed from the estimated series $\left\{\mathbf{g}\left(\mathbf{w}_{t} ; \hat{\boldsymbol{\theta}}\right)\right\}$.

\section*{The Lagrange Multiplier (LM) Statistic}
Let $\gamma_{n}(r \times 1)$ be the Lagrange multiplier for the constrained problem (7.4.4). The constrained estimator $\tilde{\theta}$ satisfies the first-order conditions


\begin{align*}
\sqrt{n} \frac{\partial Q_{n}(\tilde{\boldsymbol{\theta}})}{\partial \boldsymbol{\theta}}+\mathbf{A}(\tilde{\boldsymbol{\theta}})^{\prime} \sqrt{n} \boldsymbol{\gamma}_{n} & =\underset{(p \times 1)}{\mathbf{0}},  \tag{7.4.15}\\
\sqrt{n} \mathbf{a}(\tilde{\boldsymbol{\theta}}) & =\underset{(r \times 1)}{\mathbf{0}} . \tag{7.4.16}
\end{align*}


We seek the Taylor expansion of these first-order conditions. By condition (C), $\sqrt{n}\left(\tilde{\boldsymbol{\theta}}-\boldsymbol{\theta}_{0}\right)$ converges to a random variable. So $\sqrt{n} \mathbf{a}(\tilde{\boldsymbol{\theta}})$ admits the Taylor expansion


\begin{equation*}
\sqrt{n} \mathbf{a}(\tilde{\boldsymbol{\theta}})=\mathbf{A}_{0} \sqrt{n}\left(\tilde{\boldsymbol{\theta}}-\boldsymbol{\theta}_{0}\right)+o_{p} . \tag{7.4.17}
\end{equation*}


It is left as Review Question 2 to show that the following Taylor expansion holds for both ML and GMM:


\begin{equation*}
\sqrt{n} \frac{\partial Q_{n}(\tilde{\theta})}{\partial \boldsymbol{\theta}}=\sqrt{n} \frac{\partial Q_{n}\left(\boldsymbol{\theta}_{0}\right)}{\partial \boldsymbol{\theta}}+\Psi \sqrt{n}\left(\tilde{\boldsymbol{\theta}}-\boldsymbol{\theta}_{0}\right)+o_{p} \tag{7.4.18}
\end{equation*}


An implication of this equation is that $\sqrt{n} \frac{\partial Q_{n}(\tilde{\theta})}{\partial \theta}$ converges to a random variable. So $\sqrt{n} \boldsymbol{\gamma}_{n}$, satisfying (7.4.15), converges to a random variable. Consequently, since\\
$\mathbf{A}(\tilde{\boldsymbol{\theta}}) \rightarrow_{\mathrm{p}} \mathbf{A}_{0}$, we obtain


\begin{equation*}
\mathbf{A}(\tilde{\boldsymbol{\theta}})^{\prime} \sqrt{n} \boldsymbol{\gamma}_{n}=\mathbf{A}_{0}^{\prime} \sqrt{n} \boldsymbol{\gamma}_{n}+\left(\mathbf{A}(\tilde{\boldsymbol{\theta}})-\mathbf{A}_{0}\right)^{\prime} \sqrt{n} \boldsymbol{\gamma}_{n}=\mathbf{A}_{0}^{\prime} \sqrt{n} \boldsymbol{\gamma}_{n}+o_{p} \tag{7.4.19}
\end{equation*}


Substituting these three equations into the first-order conditions, we obtain

\[
\left[\begin{array}{cc}
\underset{(p \times p)}{\boldsymbol{\Psi}} & \underset{(p \times r)}{\mathbf{A}_{0}^{\prime}}  \tag{7.4.20}\\
\underset{(r \times p)}{\mathbf{A}_{0}} & \underset{(r \times r)}{\mathbf{0}}
\end{array}\right]\left[\begin{array}{c}
\sqrt{n}\left(\tilde{\boldsymbol{\theta}}-\boldsymbol{\theta}_{0}\right) \\
(p \times 1) \\
\sqrt{n} \boldsymbol{\gamma}_{n} \\
(r \times 1)
\end{array}\right]=\left[\begin{array}{c}
-\sqrt{n} \frac{\partial Q_{n}\left(\boldsymbol{\theta}_{0}\right)}{\partial \theta} \\
(p \times 1) \\
\mathbf{0} \\
(r \times 1)
\end{array}\right]+o_{p}
\]

Using the formula for partitioned inverses, this can be solved to yield


\begin{align*}
\sqrt{n} \boldsymbol{\gamma}_{n} & =-\left(\mathbf{A}_{0} \boldsymbol{\Psi}^{-1} \mathbf{A}_{0}^{\prime}\right)^{-1} \mathbf{A}_{0} \boldsymbol{\Psi}^{-1} \sqrt{n} \frac{\partial Q_{n}\left(\boldsymbol{\theta}_{0}\right)}{\partial \boldsymbol{\theta}}+o_{p}  \tag{7.4.21}\\
\sqrt{n}\left(\tilde{\boldsymbol{\theta}}-\boldsymbol{\theta}_{0}\right) & =-\left[\boldsymbol{\Psi}^{-1}-\boldsymbol{\Psi}^{-1} \mathbf{A}_{0}^{\prime}\left(\mathbf{A}_{0} \boldsymbol{\Psi}^{-1} \mathbf{A}_{0}^{\prime}\right)^{-1} \mathbf{A}_{0} \boldsymbol{\Psi}^{-1}\right] \sqrt{n} \frac{\partial Q_{n}\left(\boldsymbol{\theta}_{0}\right)}{\partial \boldsymbol{\theta}}+o_{p} \tag{7.4.22}
\end{align*}


The rest is the same as in the derivation of the Wald statistic. From (7.4.21), $\sqrt{n} \boldsymbol{\gamma}_{n}$ converges to a normal distribution with mean zero and variance given by


\begin{align*}
\underset{(r \times r)}{\operatorname{Avar}\left(\boldsymbol{\gamma}_{n}\right)} & =\left(\mathbf{A}_{0} \boldsymbol{\Psi}^{-1} \mathbf{A}_{0}^{\prime}\right)^{-1} \mathbf{A}_{0} \boldsymbol{\Psi}^{-1} \boldsymbol{\Sigma} \boldsymbol{\Psi}^{-1} \mathbf{A}_{0}^{\prime}\left(\mathbf{A}_{0} \boldsymbol{\Psi}^{-1} \mathbf{A}_{0}^{\prime}\right)^{-1} \\
& =\left(\mathbf{A}_{0} \boldsymbol{\Sigma}^{-1} \mathbf{A}_{0}^{\prime}\right)^{-1} \quad \text { (if condition (D) is invoked) } \tag{7.4.23}
\end{align*}


Since this $r \times r$ matrix is positive definite, the associated quadratic form


\begin{equation*}
n \boldsymbol{\gamma}_{n}^{\prime}\left(\mathbf{A}_{0} \boldsymbol{\Sigma}^{-1} \mathbf{A}_{0}^{\prime}\right) \boldsymbol{\gamma}_{n} \tag{7.4.24}
\end{equation*}


is asymptotically $\chi^{2}(r)$ under the null. If there is available a consistent estimator, denoted $\tilde{\mathbf{\Sigma}}$, of $\boldsymbol{\Sigma}$, then $\mathbf{A}(\tilde{\boldsymbol{\theta}}) \tilde{\mathbf{\Sigma}}^{-1} \mathbf{A}(\tilde{\boldsymbol{\theta}})^{\prime}$ is consistent for $\mathbf{A}_{0} \boldsymbol{\Sigma}^{-1} \mathbf{A}_{0}^{\prime}$. So the LM statistic defined by the first line below is asymptotically $\chi^{2}(r)$ under the null. This statistic can be beautified by the first-order conditions (7.4.15):


\begin{align*}
L M & \equiv n \underset{(1 \times r)}{\boldsymbol{\gamma}_{n}^{\prime}} \underbrace{\left[\mathbf{A}(\tilde{\boldsymbol{\theta}}) \widetilde{\mathbf{\Sigma}}^{-1} \mathbf{A}(\tilde{\boldsymbol{\theta}})^{\prime}\right]}_{(r \times r)} \underset{(r \times 1)}{\boldsymbol{\gamma}_{n}} \\
& =n\left(\mathbf{A}(\tilde{\boldsymbol{\theta}})^{\prime} \boldsymbol{\gamma}_{n}\right)^{\prime} \widetilde{\mathbf{\Sigma}}^{-1}\left(\mathbf{A}(\tilde{\boldsymbol{\theta}})^{\prime} \boldsymbol{\gamma}_{n}\right) \\
& =n\left(\frac{\partial Q_{n}(\tilde{\boldsymbol{\theta}})}{\partial \boldsymbol{\theta}}\right)_{(1 \times p)}^{\prime}{\underset{(p \times p)}{\widetilde{\mathbf{\Sigma}}}}^{-1}\left(\frac{\partial Q_{n}(\tilde{\boldsymbol{\theta}})}{\partial \boldsymbol{\theta}}\right) \quad\left(\text { since } \mathbf{A}(\tilde{\boldsymbol{\theta}})^{\prime} \boldsymbol{\gamma}_{n}=-\frac{\partial Q_{n}(\tilde{\boldsymbol{\theta}})}{\partial \boldsymbol{\theta}} \text { by }(7.4 .15)\right) \tag{7.4.25}
\end{align*}


The last expression for the LM statistic is the reason why the statistic is also called the score test statistic. It is the distance from zero of the score evaluated at the constrained estimate $\tilde{\boldsymbol{\theta}}$. You should not be fooled by its beauty, however: despite being a quadratic form associated with a $p$-dimensional vector, the degrees of freedom of its $\chi^{2}$ asymptotic distribution are $r$, not $p$, as is clear from the derivation.

For ML, $\widetilde{\boldsymbol{\Sigma}}$ is given by (7.4.12) or (7.4.13), evaluated at $\tilde{\boldsymbol{\theta}}$. For GMM, $\boldsymbol{\Sigma}$ is consistently estimated by


\begin{equation*}
\widetilde{\boldsymbol{\Sigma}}=\widetilde{\mathbf{G}}^{\prime} \widetilde{\mathbf{S}}^{-1} \widetilde{\mathbf{G}} \text { with } \underset{(K \times p)}{\widetilde{\mathbf{G}}} \equiv \mathbf{G}_{n}(\tilde{\boldsymbol{\theta}})=\frac{1}{n} \sum_{t=1}^{n} \frac{\partial \mathbf{g}\left(\mathbf{w}_{t} ; \tilde{\boldsymbol{\theta}}\right)}{\partial \boldsymbol{\theta}^{\prime}} \text {. } \tag{7.4.26}
\end{equation*}


Here, $\widetilde{\mathbf{S}}$ is an estimate of the long-run variance matrix constructed from the estimated series $\left\{\mathbf{g}\left(\mathbf{w}_{t} ; \tilde{\boldsymbol{\theta}}\right)\right\}$.

\section*{The Likelihood Ratio (LR) Statistic}
The LR statistic is defined as $2 n$ times $Q_{n}(\hat{\boldsymbol{\theta}})-Q_{n}(\tilde{\boldsymbol{\theta}})$. (For efficient GMM, calling this the likelihood ratio statistic is a slight abuse of the language because $Q_{n}$ for GMM is not the likelihood function, but we will continue to use the term for both ML and GMM.) It is not hard to show for ML and GMM (see Review Questions 6 and 7) that


\begin{equation*}
L R \equiv 2 n \cdot\left[Q_{n}(\hat{\boldsymbol{\theta}})-Q_{n}(\tilde{\boldsymbol{\theta}})\right]=-n(\hat{\boldsymbol{\theta}}-\tilde{\boldsymbol{\theta}})^{\prime} \Psi(\hat{\boldsymbol{\theta}}-\tilde{\boldsymbol{\theta}})+o_{p} . \tag{7.4.27}
\end{equation*}


The validity of this expression does not depend on condition (D).\\
From this expression, it is easy to derive the asymptotic distribution of the LR statistic. Subtracting (7.4.22) from (7.4.5), we obtain


\begin{equation*}
\sqrt{n}(\hat{\boldsymbol{\theta}}-\tilde{\boldsymbol{\theta}})=-\boldsymbol{\Psi}^{-1} \mathbf{A}_{0}^{\prime}\left(\mathbf{A}_{0} \boldsymbol{\Psi}^{-1} \mathbf{A}_{0}^{\prime}\right)^{-1} \mathbf{A}_{0} \boldsymbol{\Psi}^{-1} \sqrt{n} \frac{\partial Q_{n}\left(\boldsymbol{\theta}_{0}\right)}{\partial \boldsymbol{\theta}}+o_{p} . \tag{7.4.28}
\end{equation*}


Substituting this into (7.4.27), the LR statistic can be written as


\begin{equation*}
L R=-\left(\sqrt{n} \frac{\partial Q_{n}\left(\boldsymbol{\theta}_{0}\right)}{\partial \boldsymbol{\theta}}\right)^{\prime} \boldsymbol{\Psi}^{-1} \mathbf{A}_{0}^{\prime}\left(\mathbf{A}_{0} \boldsymbol{\Psi}^{-1} \mathbf{A}_{0}^{\prime}\right)^{-1} \mathbf{A}_{0} \boldsymbol{\Psi}^{-1}\left(\sqrt{n} \frac{\partial Q_{n}\left(\boldsymbol{\theta}_{0}\right)}{\partial \boldsymbol{\theta}}\right)+o_{p} . \tag{7.4.29}
\end{equation*}


Now invoke assumption (D). The LR statistic becomes


\begin{equation*}
L R=\left(\sqrt{n} \frac{\partial Q_{n}\left(\boldsymbol{\theta}_{0}\right)}{\partial \boldsymbol{\theta}}\right)^{\prime} \boldsymbol{\Sigma}^{-1} \mathbf{A}_{0}^{\prime}\left[\mathbf{A}_{0} \boldsymbol{\Sigma}^{-1} \mathbf{A}_{0}^{\prime}\right]^{-1} \mathbf{A}_{0} \boldsymbol{\Sigma}^{-1}\left(\sqrt{n} \frac{\partial Q_{n}\left(\boldsymbol{\theta}_{0}\right)}{\partial \boldsymbol{\theta}}\right)+o_{p} . \tag{7.4.30}
\end{equation*}


This is asymptotically $\chi^{2}(r)$ because the variance of the limiting distribution of $\mathbf{A}_{0} \boldsymbol{\Sigma}^{-1}\left(\sqrt{n} \frac{\partial Q_{n}\left(\boldsymbol{\theta}_{0}\right)}{\partial \theta}\right)$ is the positive definite matrix in brackets.

This completes the proof of the asymptotic equivalence of the trio of statistics. But it should be clear from the proof that something stronger than asymptotic equivalence (that the three statistics share the same asymptotic distribution) holds. We have shown in the proof that the difference between the Wald statistic and (7.4.10) is $o_{p}$. But by (7.4.8) the difference between the latter and the quadratic form in the expression for the LR statistic (7.4.30) is $o_{p}$ when $\boldsymbol{\Sigma}=\boldsymbol{-} \boldsymbol{\Psi}$. So the numerical difference between the Wald statistic and the LR statistic is $o_{p}$ when $\boldsymbol{\Sigma}=-\boldsymbol{\Psi}$. The same applies to the LM statistic. The difference between it and (7.4.24) is $o_{p}$. But the difference between the latter and the quadratic form in (7.4.30) is again $o_{p}$ when $\boldsymbol{\Sigma}=-\boldsymbol{\Psi}$.

\section*{Summary of the Trinity}
To summarize the rather lengthy discussion of this section,\\
Proposition 7.11 (The trinity): Let $\hat{\theta}$ the extremum estimator defined in (7.1.1). Consider the null hypothesis (7.4.1) where $\mathbf{a}(\cdot)$ is continuously differentiable (so the Jacobian $\mathbf{A}(\boldsymbol{\theta}) \equiv \frac{\partial \mathbf{a}(\boldsymbol{\theta})}{\partial \boldsymbol{\theta}^{\prime}}$ is continuous) and the rank condition (7.4.3) is satisfied. Let $\tilde{\theta}$ be the constrained extremum estimator defined in (7.4.4). Assume:

\begin{itemize}
  \item the hypothesis of Proposition 7.9, if the extremum estimator is conditional ML,
  \item the hypothesis appropriately modified as indicated right below Proposition 7.9, in the case of unconditional ML,
  \item the hypothesis of Proposition 7.10 with $\mathbf{W}=\mathbf{S}^{-1}$ and that there is available a consistent estimator $\widehat{\mathbf{S}}$ of $\mathbf{S}$ constructed from $\hat{\boldsymbol{\theta}}$ and $\widetilde{\mathbf{S}}$ constructed from $\tilde{\boldsymbol{\theta}}$, in the case of efficient GMM.
\end{itemize}

Define the Wald, $L M$, and $L R$ statistics as in Table 7.2. Then\\
(a) the constrained extremum estimator $\tilde{\theta}$ is consistent and asymptotically normal,\\
(b) the three statistics all converge in distribution to $\chi^{2}(r)$ under the null where $r$ is the number of restrictions in the null,\\
(c) moreover, the numerical difference between the three statistics converges in probability to zero under the null.

Table 7.2: Trinity

\begin{center}
\begin{tabular}{|l|l|l|}
\hline
 & Wald: $\underset{(1 \times r)}{n \mathbf{a}(\hat{\boldsymbol{\theta}})^{\prime}}\left[\underset{(r \times p)}{\mathbf{A}(\hat{\boldsymbol{\theta}})} \underset{(p \times p)}{\widehat{\boldsymbol{\Sigma}}^{-1}} \mathbf{A}(\hat{\boldsymbol{\theta}})^{\prime}\right]^{-1} \underset{(p \times r)}{\mathbf{a}(\hat{\boldsymbol{\theta}})} \operatorname{LM}: n\left(\frac{\partial Q_{n}(\tilde{\boldsymbol{\theta}})}{\underset{(1 \times p)}{\partial \boldsymbol{\theta}}}\right)^{\prime} \underset{(p \times p)}{\widetilde{\boldsymbol{\Sigma}}^{-1}}\left(\frac{\partial Q_{n}(\tilde{\boldsymbol{\theta}})}{\underset{(p \times 1)}{\partial \boldsymbol{\theta}}}\right)$ LR: $2 n \cdot\left[Q_{n}(\hat{\boldsymbol{\theta}})-Q_{n}(\tilde{\boldsymbol{\theta}})\right]$ &  \\
\hline
Terms for substitution & Conditional ML & Efficient GMM \\
\hline
$Q_{n}(\boldsymbol{\theta})$ & $\frac{1}{n} \sum_{t=1}^{n} \log f\left(y_{t} \mid \mathbf{x}_{t} ; \boldsymbol{\theta}\right)$ & $-\frac{1}{2} \mathbf{g}_{n}(\boldsymbol{\theta})^{\prime} \widehat{\mathbf{S}}^{-1} \mathbf{g}_{n}(\boldsymbol{\theta})$ \\
\hline
$\widehat{\mathbf{\Sigma}}$ & $-\frac{\partial^{2} Q_{n}(\hat{\boldsymbol{\theta}})}{\partial \boldsymbol{\theta} \partial \boldsymbol{\theta}^{\prime}}$ or $\frac{1}{n} \sum \mathbf{s}\left(\mathbf{w}_{t} ; \hat{\boldsymbol{\theta}}\right) \mathbf{s}\left(\mathbf{w}_{t} ; \hat{\boldsymbol{\theta}}\right)^{\prime}$ & $\widehat{\mathbf{G}}^{\prime} \widehat{\mathbf{S}}^{-1} \widehat{\mathbf{G}}, \underset{(K \times p)}{\widehat{\mathbf{G}}} \equiv \mathbf{G}_{n}(\hat{\boldsymbol{\theta}})$ \\
\hline
$\widetilde{\mathbf{\Sigma}}$ & replace $\hat{\boldsymbol{\theta}}$ by $\tilde{\boldsymbol{\theta}}$ in above & $\widetilde{\mathbf{G}}^{\prime} \widetilde{\mathbf{S}}^{-1} \widetilde{\mathbf{G}}, \underset{(K \times p)}{\widetilde{\mathbf{G}}} \equiv \mathbf{G}_{n}(\tilde{\boldsymbol{\theta}})$ \\
\hline
\end{tabular}
\end{center}

Note: See (7.3.2) and (7.3.3) for definitions of $\mathbf{s}\left(\mathbf{w}_{t} ; \boldsymbol{\theta}\right)$ and $\mathbf{H}\left(\mathbf{w}_{t} ; \boldsymbol{\theta}\right)$. Also,

$$
\mathbf{g}_{n}(\boldsymbol{\theta}) \equiv \frac{1}{n} \sum_{t=1}^{n} \mathbf{g}\left(\mathbf{w}_{t} ; \boldsymbol{\theta}\right), \quad \mathbf{G}_{n}(\boldsymbol{\theta}) \equiv \frac{\partial \mathbf{g}_{n}(\boldsymbol{\theta})}{\partial \boldsymbol{\theta}^{\prime}}, \quad \text { and } \quad \mathbf{G} \equiv \mathrm{E}\left[\frac{\partial \mathbf{g}\left(\mathbf{w}_{t} ; \boldsymbol{\theta}_{0}\right)}{\partial \boldsymbol{\theta}^{\prime}}\right] .
$$

The same middle matrix is used in $Q_{n}(\hat{\boldsymbol{\theta}})$ and $Q_{n}(\tilde{\boldsymbol{\theta}})$ for GMM.

\section*{QUESTIONS FOR REVIEW}
\begin{enumerate}
  \item Verify that the LM statistic for ML can be written as
\end{enumerate}

$$
n\left(\frac{1}{n} \sum_{t=1}^{n} \mathbf{s}\left(\mathbf{w}_{t} ; \tilde{\boldsymbol{\theta}}\right)\right)^{\prime}\left\{\frac{1}{n} \sum \mathbf{s}\left(\mathbf{w}_{t} ; \tilde{\boldsymbol{\theta}}\right) \mathbf{s}\left(\mathbf{w}_{t} ; \tilde{\boldsymbol{\theta}}\right)^{\prime}\right\}^{-1}\left(\frac{1}{n} \sum_{t=1}^{n} \mathbf{s}\left(\mathbf{w}_{t} ; \tilde{\boldsymbol{\theta}}\right)\right) .
$$

\begin{enumerate}
  \setcounter{enumi}{1}
  \item Derive (7.4.18). Hint: For ML, replace $\hat{\boldsymbol{\theta}}$ by $\tilde{\boldsymbol{\theta}}$ in (7.3.5). For GMM, derive
\end{enumerate}

$$
\sqrt{n} \frac{\partial Q_{n}(\tilde{\boldsymbol{\theta}})}{\partial \boldsymbol{\theta}}=-\mathbf{G}_{n}(\tilde{\boldsymbol{\theta}})^{\prime} \widehat{\mathbf{W}} \frac{1}{\sqrt{n}} \sum_{t=1}^{n} \mathbf{g}\left(\mathbf{w}_{t} ; \boldsymbol{\theta}_{0}\right)-\mathbf{G}_{n}(\tilde{\boldsymbol{\theta}})^{\prime} \widehat{\mathbf{W}} \mathbf{G}_{n}(\overline{\boldsymbol{\theta}}) \sqrt{n}\left(\tilde{\boldsymbol{\theta}}-\boldsymbol{\theta}_{0}\right),
$$

where $\overline{\boldsymbol{\theta}}$ lies between $\boldsymbol{\theta}_{0}$ and $\tilde{\boldsymbol{\theta}}$. The derivation of this should be as easy as the derivation of (7.3.31).\\
3. Explain why the LR statistic would not be asymptotically $\chi^{2}$ if $\mathbf{\Sigma}=-\mathbf{\Psi}$ did not hold.\\
4. For ML, suppose the trio of statistics are calculated as indicated in Table 7.2, but suppose $\mathbf{w}_{t}$ is actually serially correlated. Which one is still asymptotically $\chi^{2}$ ? [Answer: None.]\\
5. Suppose $\mathbf{\Sigma}=-\mathbf{\Psi}$ does not hold but suppose that a consistent estimate $\widetilde{\Psi}$ of $\mathbf{\Psi}$ is available along with the consistent estimate $\boldsymbol{\Sigma}$ of $\boldsymbol{\Sigma}$. Propose an LM statistic that is asymptotically $\chi^{2}(r)$. Hint: Use the first equality in (7.4.23), which is valid without $\mathbf{\Sigma}=-\mathbf{\Psi}$. The answer is


\begin{equation*}
n\left(\frac{\partial Q_{n}(\tilde{\boldsymbol{\theta}})}{\partial \boldsymbol{\theta}}\right)^{\prime} \widetilde{\Psi}^{-1} \mathbf{A}(\tilde{\boldsymbol{\theta}})^{\prime}\left[\mathbf{A}(\tilde{\boldsymbol{\theta}}) \widetilde{\Psi}^{-1} \widetilde{\boldsymbol{\Sigma}} \widetilde{\Psi}^{-1} \mathbf{A}(\tilde{\boldsymbol{\theta}})^{\prime}\right]^{-1} \mathbf{A}(\tilde{\boldsymbol{\theta}}) \widetilde{\Psi}^{-1}\left(\frac{\partial Q_{n}(\tilde{\boldsymbol{\theta}})}{\partial \boldsymbol{\theta}}\right) . \tag{7.4.31}
\end{equation*}


\begin{enumerate}
  \setcounter{enumi}{5}
  \item (Optional, proof of (7.4.27)) Prove (7.4.27) for M-estimators. Hint: Take for granted the "second-order" mean value expansion around $\hat{\boldsymbol{\theta}}$,
\end{enumerate}


\begin{equation*}
Q_{n}(\tilde{\boldsymbol{\theta}})=Q_{n}(\hat{\boldsymbol{\theta}})+\frac{\partial Q_{n}(\hat{\boldsymbol{\theta}})}{\partial \boldsymbol{\theta}^{\prime}}(\tilde{\boldsymbol{\theta}}-\hat{\boldsymbol{\theta}})+\frac{1}{2}(\tilde{\boldsymbol{\theta}}-\hat{\boldsymbol{\theta}})^{\prime} \frac{\partial^{2} Q_{n}(\overline{\boldsymbol{\theta}})}{\partial \boldsymbol{\theta} \partial \boldsymbol{\theta}^{\prime}}(\tilde{\boldsymbol{\theta}}-\hat{\boldsymbol{\theta}}), \tag{7.4.32}
\end{equation*}


where $\overline{\boldsymbol{\theta}}$ lies between $\hat{\boldsymbol{\theta}}$ and $\tilde{\boldsymbol{\theta}}$. By the first-order condition, $\frac{\partial Q_{n}(\hat{\boldsymbol{\theta}})}{\partial \boldsymbol{\theta}^{\prime}}=\mathbf{0}$. Also, $\frac{\partial^{2} Q_{n}(\overline{\boldsymbol{\theta}})}{\partial \boldsymbol{\theta} \partial \boldsymbol{\theta}^{\prime}} \rightarrow \mathrm{p}\left[\mathbf{H}\left(\mathbf{w}_{t} ; \boldsymbol{\theta}_{0}\right)\right]$.\\
7. (Optional, proof of (7.4.27) for GMM) Prove (7.4.27) for GMM. Hint: From the mean value expansion of $\mathbf{g}_{n}(\boldsymbol{\theta})$ around $\hat{\boldsymbol{\theta}}$, derive the Taylor expansion


\begin{equation*}
\mathbf{g}_{n}(\tilde{\boldsymbol{\theta}})=\mathbf{g}_{n}(\hat{\boldsymbol{\theta}})+\mathbf{G}_{n}(\hat{\boldsymbol{\theta}})(\tilde{\boldsymbol{\theta}}-\hat{\boldsymbol{\theta}})+o_{p} . \tag{7.4.33}
\end{equation*}


Substitute this into $Q_{n}(\tilde{\boldsymbol{\theta}})=-\frac{1}{2} \mathbf{g}_{n}(\tilde{\boldsymbol{\theta}})^{\prime} \widehat{\mathbf{W}} \mathbf{g}_{n}(\tilde{\boldsymbol{\theta}})$. Then use the first-order condition that $\mathbf{G}_{n}(\hat{\boldsymbol{\theta}})^{\prime} \widehat{\mathbf{W}} \mathbf{g}_{n}(\hat{\boldsymbol{\theta}})=\mathbf{0}$.

\subsection*{7.5 Numerical Optimization}
So far, we have not been concerned about the computational aspect of extremum estimators. In the ML estimation of the linear regression model, for example, the objective function $Q_{n}(\boldsymbol{\theta})$ is quadratic in $\boldsymbol{\theta}$, so there is a closed-form formula for the maximizer, which is the OLS estimator. The objective function is quadratic for linear GMM also, with the linear GMM estimator being the closed-form formula. In most other cases, however, no such closed-form formulas are available, and some numerical algorithm must be employed to locate the maximum. A number of algorithms are available, but this section covers only the two most important ones. More details can be found in, e.g., Judge et al. (1985, Appendix B). The first is used for calculating M-estimators, while the second is used for GMM.

\section*{Newton-Raphson}
We first consider M-estimators, whose objective function $Q_{n}(\boldsymbol{\theta})$ is (7.1.2). Since $Q_{n}(\boldsymbol{\theta})$ is twice continuously differentiable for M-estimators, there is a secondorder Taylor expansion


\begin{equation*}
Q_{n}(\boldsymbol{\theta}) \cong Q_{n}\left(\hat{\boldsymbol{\theta}}_{j}\right)+\mathbf{s}_{n}\left(\hat{\boldsymbol{\theta}}_{j}\right)^{\prime}\left(\boldsymbol{\theta}-\hat{\boldsymbol{\theta}}_{j}\right)+\frac{1}{2}\left(\boldsymbol{\theta}-\hat{\boldsymbol{\theta}}_{j}\right)^{\prime} \mathbf{H}_{n}\left(\hat{\boldsymbol{\theta}}_{j}\right)\left(\boldsymbol{\theta}-\hat{\boldsymbol{\theta}}_{j}\right), \tag{7.5.1}
\end{equation*}


where $\hat{\boldsymbol{\theta}}_{j}$ is the estimate in the $j$-th round of the iterative procedure to be described in a moment, and $\mathbf{s}_{n}$ and $\mathbf{H}_{n}$ are the gradient and the Hessian of the objective function:


\begin{equation*}
\mathbf{s}_{n}(\boldsymbol{\theta}) \equiv \frac{\partial Q_{n}(\boldsymbol{\theta})}{\partial \boldsymbol{\theta}^{\prime}}, \quad \mathbf{H}_{n}(\boldsymbol{\theta}) \equiv \frac{\partial^{2} Q_{n}(\boldsymbol{\theta})}{\partial \boldsymbol{\theta} \partial \boldsymbol{\theta}^{\prime}} . \tag{7.5.2}
\end{equation*}


The $(j+1)$-th round estimator $\hat{\boldsymbol{\theta}}_{j+1}$ is the maximizer of the quadratic function on the right-hand side of (7.5.1). It is given by


\begin{equation*}
\hat{\boldsymbol{\theta}}_{j+1}=\hat{\boldsymbol{\theta}}_{j}-\left[\mathbf{H}_{n}\left(\hat{\boldsymbol{\theta}}_{j}\right)\right]^{-1} \mathbf{s}_{n}\left(\hat{\boldsymbol{\theta}}_{j}\right) . \tag{7.5.3}
\end{equation*}


This iterative procedure is called the Newton-Raphson algorithm. If the objective function is concave, the algorithm often converges quickly to the global maximum.

For ML estimators, when the global maximum $\hat{\boldsymbol{\theta}}$ is thus obtained, the estimate of $\operatorname{Avar}(\hat{\boldsymbol{\theta}})$ is obtained as $-\mathbf{H}_{n}(\hat{\boldsymbol{\theta}})^{-1}$ (recall that $\mathbf{H}_{n}(\boldsymbol{\theta}) \equiv \frac{1}{n} \sum \mathbf{H}\left(\mathbf{w}_{t} ; \boldsymbol{\theta}\right)$ ).

\section*{Gauss-Newton}
Now turn to GMM, whose objective function is $Q_{n}(\boldsymbol{\theta})=-\frac{1}{2} \mathbf{g}_{n}(\boldsymbol{\theta})^{\prime} \widehat{\mathbf{W}} \mathbf{g}_{n}(\boldsymbol{\theta})$ (see (7.1.3)). As in the derivation of the asymptotic distribution, we work with a linearization of the $\mathbf{g}_{n}(\boldsymbol{\theta})$ function. The first-order Taylor expansion of $\mathbf{g}_{n}(\boldsymbol{\theta})$ around $\hat{\boldsymbol{\theta}}_{j}$ is


\begin{align*}
\underset{(K \times 1)}{\mathbf{g}_{n}(\boldsymbol{\theta})} & \cong \mathbf{g}_{n}\left(\hat{\boldsymbol{\theta}}_{j}\right)+\underset{(K \times p)}{\mathbf{G}_{n}\left(\hat{\boldsymbol{\theta}}_{j}\right)\left(\boldsymbol{\theta}-\hat{\boldsymbol{\theta}}_{j}\right)} \\
& =\left[\mathbf{g}_{n}\left(\hat{\boldsymbol{\theta}}_{j}\right)-\mathbf{G}_{n}\left(\hat{\boldsymbol{\theta}}_{j}\right) \hat{\boldsymbol{\theta}}_{j}\right]-\left[-\mathbf{G}_{n}\left(\hat{\boldsymbol{\theta}}_{j}\right)\right] \boldsymbol{\theta} \\
& =\mathbf{v}_{j}-\mathbf{G}_{j} \boldsymbol{\theta}, \tag{7.5.4}
\end{align*}


where


\begin{equation*}
\mathbf{v}_{j} \equiv \mathbf{g}_{n}\left(\hat{\boldsymbol{\theta}}_{j}\right)-\mathbf{G}_{n}\left(\hat{\boldsymbol{\theta}}_{j}\right) \hat{\boldsymbol{\theta}}_{j}, \quad \mathbf{G}_{j} \equiv-\mathbf{G}_{n}\left(\hat{\boldsymbol{\theta}}_{j}\right) . \tag{7.5.5}
\end{equation*}


(Recall that $\mathbf{G}_{n}(\boldsymbol{\theta}) \equiv \frac{\partial \mathbf{g}_{n}(\boldsymbol{\theta})}{\partial \boldsymbol{\theta}^{\prime}}$.) If the $\mathbf{g}_{n}(\boldsymbol{\theta})$ function in the expression for the GMM objective function were this linear function $\mathbf{v}_{j}-\mathbf{G}_{j} \boldsymbol{\theta}$, then the objective function would be quadratic in $\theta$ and the maximizer (or the minimizer of the GMM distance) would be the linear GMM estimator:


\begin{equation*}
\hat{\boldsymbol{\theta}}_{j+1}=\left(\mathbf{G}_{j}^{\prime} \widehat{\mathbf{W}} \mathbf{G}_{j}\right)^{-1} \mathbf{G}_{j}^{\prime} \widehat{\mathbf{W}} \mathbf{v}_{j} . \tag{7.5.6}
\end{equation*}


This is the $(j+1)$-th round estimate in the Gauss-Newton algorithm. ${ }^{19}$ Unlike in Newton-Raphson, there is no need to calculate second-order derivatives.

\section*{Writing Newton-Raphson and Gauss-Newton in a Common Format}
There are also similarities between Gauss-Newton and Newton-Raphson. Substituting (7.5.5) into (7.5.6) and rearranging, we obtain

\footnotetext{${ }^{19}$ The Gauss-Newton algorithm was originally developed for nonlinear least squares. See Review Question 2 for how Gauss-Newton is applied to NLS. The algorithm presented here is therefore its GMM adaptation.
}

\begin{equation*}
\hat{\boldsymbol{\theta}}_{j+1}=\hat{\boldsymbol{\theta}}_{j}-\left[-\mathbf{G}_{n}\left(\hat{\boldsymbol{\theta}}_{j}\right)^{\prime} \widehat{\mathbf{W}} \mathbf{G}_{n}\left(\hat{\boldsymbol{\theta}}_{j}\right)\right]^{-1}\left[-\mathbf{G}_{n}\left(\hat{\boldsymbol{\theta}}_{j}\right)^{\prime} \widehat{\mathbf{W}} \mathbf{g}_{n}\left(\hat{\boldsymbol{\theta}}_{j}\right)\right] . \tag{7.5.7}
\end{equation*}


The second bracketed term is none other than the gradient evaluated at $\hat{\boldsymbol{\theta}}_{j}$ for the GMM objective function $Q_{n}(\boldsymbol{\theta})=-\frac{1}{2} \mathbf{g}_{n}(\boldsymbol{\theta})^{\prime} \widehat{\mathbf{W}} \mathbf{g}_{n}(\boldsymbol{\theta})$. The role played by the Hessian of the objective function in (7.5.3) is played here by the first bracketed term, which is the estimate based on $\hat{\boldsymbol{\theta}}_{j}$ of the $-\mathbf{G}^{\prime}$ WG in Table 7.1. Therefore, the analogy between M-estimators and GMM noted in Table 7.1 also holds here. Furthermore, if $\mathbf{g}_{n}$ is linear, then the Hessian equivalent (the first bracketed term in (7.5.7)) is the Hessian of the GMM objective function. So the analogy is exact: the Gauss-Newton algorithm coincides with the Newton-Raphson algorithm.

\section*{Equations Nonlinear in Parameters Only}
In the Gauss-Newton algorithm (7.5.7), when $\mathbf{g}_{n}(\boldsymbol{\theta})\left(\equiv \frac{1}{n} \sum \mathbf{g}\left(\mathbf{w}_{t} ; \boldsymbol{\theta}\right)\right)$ is not linear in $\theta$, it is generally necessary to take averages over observations to evaluate the gradient and the Hessian equivalent in each iteration. (This is also generally true in Newton-Raphson if the objective function is not quadratic in $\boldsymbol{\theta}$.) For large data sets it is a very CPU time-intensive operation. There is, however, one case where the averaging over observations in each iteration is not needed. That occurs in a class of generalized nonlinear instrumental variable estimation (a special case of nonlinear GMM where $\mathbf{g}\left(y_{t}, \mathbf{z}_{t}, \boldsymbol{\theta}\right)$ can be written as $\left.a\left(y_{t}, \mathbf{z}_{t} ; \boldsymbol{\theta}\right) \mathbf{x}_{t}\right)$. Suppose that the equation $a\left(y_{t}, \mathbf{z}_{t} ; \boldsymbol{\theta}\right)$ takes the form


\begin{equation*}
a\left(y_{t}, \mathbf{z}_{t} ; \boldsymbol{\theta}\right)=a_{0}\left(y_{t}, \mathbf{z}_{t}\right)+\underset{(1 \times q)}{\mathbf{a}_{1}\left(y_{t}, \mathbf{z}_{t}\right)^{\prime}} \boldsymbol{\alpha}(\boldsymbol{\theta}(\boldsymbol{\theta})) \tag{7.5.8}
\end{equation*}


where $a_{0}(\cdot)$ and $\mathbf{a}_{1}(\cdot)$ are known functions of $\left(y_{t}, \mathbf{z}_{t}\right)$ (but not functions of $\boldsymbol{\theta}$ ). A function of this form is said to be nonlinear in parameters only or linear in variables. It can be easily seen that


\begin{align*}
\mathbf{g}_{n}(\boldsymbol{\theta}) & \equiv \frac{1}{n} \sum_{t=1}^{n} \mathbf{g}\left(y_{t}, \mathbf{z}_{t} ; \boldsymbol{\theta}\right)=\frac{1}{n} \sum_{t=1}^{n} a\left(y_{t}, \mathbf{z}_{t} ; \boldsymbol{\theta}\right) \cdot \mathbf{x}_{t} \\
& =\left(\frac{1}{n} \sum_{t=1}^{n} a_{0}\left(y_{t}, \mathbf{z}_{t}\right) \cdot \mathbf{x}_{t}\right)+\left(\frac{1}{n} \sum_{t=1}^{n} \mathbf{x}_{t} \mathbf{a}_{1}\left(y_{t}, \mathbf{z}_{t}\right)^{\prime}\right) \underset{(K \times 1)}{\alpha}(\boldsymbol{\theta}),  \tag{7.5.9}\\
\mathbf{G}_{n}(\boldsymbol{\theta}) & =\frac{\partial \mathbf{g}_{n}(\boldsymbol{\theta})}{\partial \boldsymbol{\theta}^{\prime}}=\left(\frac{1}{n} \sum_{t=1}^{n} \mathbf{x}_{t} \mathbf{a}_{1}\left(y_{t}, \mathbf{z}_{t}\right)^{\prime}\right) \frac{\partial \boldsymbol{\alpha}(\boldsymbol{\theta})}{\partial \boldsymbol{\theta}^{\prime}} . \tag{7.5.10}
\end{align*}


These equations make it clear that the sample averages need to be calculated just once, before the start of iterations.

\section*{QUESTIONS FOR REVIEW}
\begin{enumerate}
  \item (Does $Q_{n}$ increase during iterations?) Set $\boldsymbol{\theta}=\hat{\boldsymbol{\theta}}_{j+1}$ in (7.5.1) and rewrite it as
\end{enumerate}

$$
Q_{n}\left(\hat{\boldsymbol{\theta}}_{j+1}\right)-Q_{n}\left(\hat{\boldsymbol{\theta}}_{j}\right) \cong \mathbf{s}_{n}\left(\hat{\boldsymbol{\theta}}_{j}\right)^{\prime}\left(\hat{\boldsymbol{\theta}}_{j+1}-\hat{\boldsymbol{\theta}}_{j}\right)+\frac{1}{2}\left(\hat{\boldsymbol{\theta}}_{j+1}-\hat{\boldsymbol{\theta}}_{j}\right)^{\prime} \mathbf{H}_{n}\left(\hat{\boldsymbol{\theta}}_{j}\right)\left(\hat{\boldsymbol{\theta}}_{j+1}-\hat{\boldsymbol{\theta}}_{j}\right) .
$$

Consider the Newton-Raphson algorithm (7.5.3). Suppose that $\hat{\boldsymbol{\theta}}_{j+1}$ is sufficiently close to $\hat{\boldsymbol{\theta}}_{j}$ so that the sign of $Q_{n}\left(\hat{\boldsymbol{\theta}}_{j+1}\right)-Q_{n}\left(\hat{\boldsymbol{\theta}}_{j}\right)$ is the same as that of the right hand side of the equation. Show that $Q_{n}\left(\hat{\boldsymbol{\theta}}_{j+1}\right) \geq Q_{n}\left(\hat{\boldsymbol{\theta}}_{j}\right)$ if $Q_{n}(\boldsymbol{\theta})$ is concave. Hint: Using (7.5.3), the right-hand side can be written as

$$
-\frac{1}{2}\left(\hat{\boldsymbol{\theta}}_{j+1}-\hat{\boldsymbol{\theta}}_{j}\right)^{\prime} \mathbf{H}_{n}\left(\hat{\boldsymbol{\theta}}_{j}\right)\left(\hat{\boldsymbol{\theta}}_{j+1}-\hat{\boldsymbol{\theta}}_{j}\right) .
$$

\begin{enumerate}
  \setcounter{enumi}{1}
  \item (Gauss-Newton for NLS) In NLS (nonlinear least squares), the objective function is given by (7.1.19). Let $\hat{\boldsymbol{\theta}}_{j}$ be the estimate in the $j$-th round and consider the linear approximation
\end{enumerate}


\begin{align*}
\varphi\left(\mathbf{x}_{t} ; \boldsymbol{\theta}\right) & \cong \varphi\left(\mathbf{x}_{t} ; \hat{\boldsymbol{\theta}}_{j}\right)+\left(\frac{\partial \varphi\left(\mathbf{x}_{t} ; \hat{\boldsymbol{\theta}}_{j}\right)}{\partial \boldsymbol{\theta}}\right)^{\prime}\left(\boldsymbol{\theta}-\hat{\boldsymbol{\theta}}_{j}\right) \\
(\boldsymbol{p} \times \mathbf{l}) &  \tag{7.5.11}\\
& =\left[\varphi\left(\mathbf{x}_{t} ; \hat{\boldsymbol{\theta}}_{j}\right)-\left(\frac{\partial \varphi\left(\mathbf{x}_{t} ; \hat{\boldsymbol{\theta}}_{j}\right)}{\partial \boldsymbol{\theta}}\right)^{\prime} \hat{\boldsymbol{\theta}}_{j}\right]+\left(\frac{\partial \varphi\left(\mathbf{x}_{t} ; \hat{\boldsymbol{\theta}}_{j}\right)}{\partial \boldsymbol{\theta}}\right)^{\prime} \boldsymbol{\theta}
\end{align*}


Replace $\varphi\left(\mathbf{x}_{t} ; \boldsymbol{\theta}\right)$ in the $m$ function by this linear function in $\boldsymbol{\theta}$. Show that the maximizer of this linearized problem is


\begin{align*}
\hat{\boldsymbol{\theta}}_{j+1}=\hat{\boldsymbol{\theta}}_{j}+\left[\frac{1}{n} \sum_{t=1}^{n}\left(\frac{\partial \varphi\left(\mathbf{x}_{t} ; \hat{\boldsymbol{\theta}}_{j}\right)}{\partial \boldsymbol{\theta}}\right)\left(\frac{\partial \varphi\left(\mathbf{x}_{t} ; \hat{\boldsymbol{\theta}}_{j}\right)}{\partial \boldsymbol{\theta}}\right)^{\prime}\right]^{-1} & \\
& \left\{\frac{1}{n} \sum_{t=1}^{n}\left(\frac{\partial \varphi\left(\mathbf{x}_{t} ; \hat{\boldsymbol{\theta}}_{j}\right)}{\partial \boldsymbol{\theta}}\right) \cdot\left[y_{t}-\varphi\left(\mathbf{x}_{t} ; \hat{\boldsymbol{\theta}}_{j}\right)\right]\right\} . \tag{7.5.12}
\end{align*}


\section*{ANALYTICAL EXERCISES}
\begin{enumerate}
  \item (Kullback-Leibler information inequality) Let $f(y \mid \mathbf{x} ; \boldsymbol{\theta})$ be a parametric family of hypothetical conditional density functions, with the true density function given by $f\left(y \mid \mathbf{x} ; \boldsymbol{\theta}_{0}\right)$. Suppose $\mathrm{E}[\log f(y \mid \mathbf{x} ; \boldsymbol{\theta})]$ exists and is finite for all $\boldsymbol{\theta}$. The Kullback-Leibler information inequality states that
\end{enumerate}

$$
\begin{aligned}
& \operatorname{Prob}\left[f(y \mid \mathbf{x} ; \boldsymbol{\theta}) \neq f\left(y \mid \mathbf{x} ; \boldsymbol{\theta}_{0}\right)\right]>0 \\
& \quad \Rightarrow \mathrm{E}[\log f(y \mid \mathbf{x} ; \boldsymbol{\theta})]<\mathrm{E}\left[\log f\left(y \mid \mathbf{x} ; \boldsymbol{\theta}_{0}\right)\right] .
\end{aligned}
$$

Prove this by taking the following steps. In the proof, for simplicity only, assume that $f(y \mid \mathbf{x} ; \boldsymbol{\theta})>0$ for all $(y, \mathbf{x})$ and $\boldsymbol{\theta}$, so that it is legitimate to take logs of the density $f(y \mid \mathbf{x} ; \boldsymbol{\theta})$.\\
(a) Suppose $\operatorname{Prob}\left[f(y \mid \mathbf{x} ; \boldsymbol{\theta}) \neq f\left(y \mid \mathbf{x} ; \boldsymbol{\theta}_{0}\right)\right]>0$. Let $\mathbf{w}=\left(y, \mathbf{x}^{\prime}\right)^{\prime}$ and define

$$
a(\mathbf{w}) \equiv f(y \mid \mathbf{x} ; \boldsymbol{\theta}) / f\left(y \mid \mathbf{x} ; \boldsymbol{\theta}_{0}\right) .
$$

Verify that $a(\mathbf{w}) \neq 1$ with positive probability, so that $a(\mathbf{w})$ is a nonconstant random variable.\\
(b) The strict version of Jensen's inequality states that\\
if $c(x)$ is a strictly concave function and $x$ is a nonconstant random variable, then $\mathrm{E}[c(x)]<c(\mathrm{E}(x))$.

Using this fact, show that

$$
\mathrm{E}[\log a(\mathbf{w})]<\log \{\mathrm{E}[a(\mathbf{w})]\} .
$$

Hint: $\log (x)$ is strictly concave.\\
(c) Show that $\mathrm{E}[a(\mathbf{w})]=1$. Hint: The conditional mean of $a(\mathbf{w})$ equals 1 because

$$
\begin{aligned}
\mathrm{E}[a(\mathbf{w}) \mid \mathbf{x}] & =\int a(\mathbf{w}) f\left(y \mid \mathbf{x} ; \boldsymbol{\theta}_{0}\right) \mathrm{d} y \\
& =\int \frac{f(y \mid \mathbf{x} ; \boldsymbol{\theta})}{f\left(y \mid \mathbf{x} ; \boldsymbol{\theta}_{0}\right)} f\left(y \mid \mathbf{x} ; \boldsymbol{\theta}_{0}\right) \mathrm{d} y \\
& =\int f(y \mid \mathbf{x} ; \boldsymbol{\theta}) \mathrm{d} y \\
& =1
\end{aligned}
$$

The last equality holds because $f(y \mid \mathbf{x} ; \boldsymbol{\theta})$ is a density function for any $\mathbf{x}$ and $\boldsymbol{\theta}$. Note well that the conditional expectation is taken with respect to the true conditional density $f\left(y \mid \mathbf{x} ; \boldsymbol{\theta}_{0}\right)$.\\
(d) Finally, show the desired result.\\
2. (Information matrix equality) For ML, the score vector and the Hessian for observation $(y, \mathbf{x})$ can be rewritten (with $\left.\mathbf{w} \equiv\left(y, \mathbf{x}^{\prime}\right)^{\prime}\right)$ as:

$$
\begin{aligned}
\underset{(p \times 1)}{\mathbf{s}(\mathbf{w} ; \boldsymbol{\theta})} & =\frac{\partial \log f(y \mid \mathbf{x} ; \boldsymbol{\theta})}{\partial \boldsymbol{\theta}} \\
\underset{(p \times p)}{\mathbf{H}(\mathbf{w} ; \boldsymbol{\theta})} & =\frac{\partial s(\mathbf{w} ; \boldsymbol{\theta})}{\partial \boldsymbol{\theta}^{\prime}}=\frac{\partial^{2} \log f(y \mid \mathbf{x} ; \boldsymbol{\theta})}{\partial \boldsymbol{\theta} \partial \boldsymbol{\theta}^{\prime}}
\end{aligned}
$$

As usual, for simplicity, assume $f(y \mid \mathbf{x}, \boldsymbol{\theta})>0$ for all $(y, \mathbf{x})$ and $\boldsymbol{\theta}$, so it is legitimate to take logs of the density function.\\
(a) Assuming that integration (i.e., taking expectations) and differentiation can be interchanged, show that $\mathrm{E}\left[\mathbf{s}\left(\mathbf{w} ; \boldsymbol{\theta}_{0}\right)\right]=\mathbf{0}$. Hint: Since $f(y \mid \mathbf{x} ; \boldsymbol{\theta})$ is a hypothetical density, its integral is unity:

$$
\int f(y \mid \mathbf{x} ; \boldsymbol{\theta}) \mathrm{d} y=1
$$

Differentiate both sides of this equation with respect to $\theta$ and then change the order of differentiation and integration to obtain the identity

$$
\int \mathbf{s}(\mathbf{w} ; \boldsymbol{\theta}) f(y \mid \mathbf{x} ; \boldsymbol{\theta}) \mathrm{d} y=\underset{(p \times 1)}{\mathbf{0}}
$$

Set $\boldsymbol{\theta}=\boldsymbol{\theta}_{0}$ and obtain $\mathrm{E}\left[\mathbf{s}\left(\mathbf{w} ; \boldsymbol{\theta}_{0}\right) \mid \mathbf{x}\right]=\mathbf{0}$.\\
(b) Show the information matrix equality:

$$
-\mathrm{E}\left[\mathbf{H}\left(\mathbf{w} ; \boldsymbol{\theta}_{0}\right)\right]=\mathrm{E}\left[\mathbf{s}\left(\mathbf{w} ; \boldsymbol{\theta}_{0}\right) \mathbf{s}\left(\mathbf{w} ; \boldsymbol{\theta}_{0}\right)^{\prime}\right]
$$

Hint: Differentiating both sides of the above identity in the previous hint and assuming that the order of differentiation and integration can be interchanged, we obtain

$$
\int \frac{\partial}{\partial \boldsymbol{\theta}^{\prime}} \underbrace{[\mathbf{s}(\mathbf{w} ; \boldsymbol{\theta}) f(y \mid \mathbf{x} ; \boldsymbol{\theta})]}_{(p \times 1)} \mathrm{d} y=\underset{(p \times p)}{\mathbf{0}} .
$$

Show that the integrand can be written as

$$
\begin{aligned}
& \frac{\partial}{\partial \theta^{\prime}}[\mathbf{s}(\mathbf{w} ; \boldsymbol{\theta}) f(y \mid \mathbf{x} ; \boldsymbol{\theta})] \\
& \quad=\mathbf{H}(\mathbf{w} ; \boldsymbol{\theta}) f(y \mid \mathbf{x} ; \boldsymbol{\theta})+\mathbf{s}(\mathbf{w} ; \boldsymbol{\theta}) \mathbf{s}(\mathbf{w} ; \boldsymbol{\theta})^{\prime} f(y \mid \mathbf{x} ; \boldsymbol{\theta}) .
\end{aligned}
$$

\begin{enumerate}
  \setcounter{enumi}{2}
  \item (Trinity for linear regression model) Consider the linear regression model with normal errors, whose conditional density for observation $t$ is
\end{enumerate}

$$
\log f\left(y_{t} \mid \mathbf{x}_{t} ; \boldsymbol{\beta}, \sigma^{2}\right)=-\frac{1}{2} \log (2 \pi)-\frac{1}{2} \log \left(\sigma^{2}\right)-\frac{\left(y_{t}-\mathbf{x}_{t}^{\prime} \boldsymbol{\beta}\right)^{2}}{2 \sigma^{2}} .
$$

Let $\left(\widehat{\boldsymbol{\beta}}, \hat{\sigma}^{2}\right)$ be the unrestricted ML estimate of $\boldsymbol{\theta}=\left(\boldsymbol{\beta}^{\prime} \sigma^{2}\right)^{\prime}$ and let $\left(\tilde{\boldsymbol{\beta}}, \tilde{\sigma}^{2}\right)$ be the restricted ML estimate subject to the constraint $\mathbf{R} \boldsymbol{\beta}=\mathbf{c}$ where $\mathbf{R}$ is an $r \times K$ matrix of known constants. Assume that $\Theta=\mathbb{R}^{K} \times \mathbb{R}_{++}$and that $\mathrm{E}\left(\mathbf{x}_{t} \mathbf{x}_{t}^{\prime}\right)$ is nonsingular. Also, let

$$
\widehat{\mathbf{\Sigma}}=\left[\begin{array}{cc}
\frac{1}{\hat{\sigma}^{2}} \frac{1}{n} \sum_{t=1}^{n} \mathbf{x}_{t} \mathbf{x}_{t}^{\prime} & \mathbf{0} \\
\mathbf{0}^{\prime} & \frac{1}{2\left(\hat{\sigma}^{2}\right)^{2}}
\end{array}\right], \quad \widetilde{\mathbf{\Sigma}}=\left[\begin{array}{cc}
\frac{1}{\hat{\sigma}^{2}} \frac{1}{n} \sum_{t=1}^{n} \mathbf{x}_{t} \mathbf{x}_{t}^{\prime} & \mathbf{0} \\
\mathbf{0}^{\prime} & \frac{1}{2\left(\hat{\sigma}^{2}\right)^{2}}
\end{array}\right] .
$$

(a) Verify that $\widehat{\boldsymbol{\beta}}$ minimizes the sum of squared residuals. So it is the OLS estimator. Verify that $\widetilde{\beta}$ minimizes the sum of squared residuals subject to the constraint $\mathbf{R} \boldsymbol{\beta}=\mathbf{c}$. So it is the restricted least squares estimator.\\
(b) Let $Q_{n}(\boldsymbol{\theta})=\frac{1}{n} \sum_{t=1}^{n} \log f\left(y_{t} \mid \mathbf{x}_{t} ; \boldsymbol{\beta}, \sigma^{2}\right)$. Show that

$$
\begin{aligned}
Q_{n}(\hat{\theta}) & =-\frac{1}{2} \log (2 \pi)-\frac{1}{2}-\frac{1}{2} \log \left(\frac{S S R_{U}}{n}\right), \\
Q_{n}(\tilde{\theta}) & =-\frac{1}{2} \log (2 \pi)-\frac{1}{2}-\frac{1}{2} \log \left(\frac{S S R_{R}}{n}\right),
\end{aligned}
$$

where $S S R_{U}\left(\equiv \sum\left(y_{t}-\mathbf{x}_{t}^{\prime} \widehat{\boldsymbol{\beta}}\right)^{2}\right)$ is the unrestricted sum of squared residuals and $S S R_{R}\left(\equiv \sum\left(y_{t}-\mathbf{x}_{t}^{\prime} \tilde{\boldsymbol{\beta}}\right)^{2}\right)$ is the restricted sum of squared residuals.\\
Hint: Show that $\hat{\sigma}^{2}=S S R_{U} / n$ and $\tilde{\sigma}^{2}=S S R_{R} / n$.\\
(c) Verify that the $\widehat{\mathbf{\Sigma}}$ given here, although not the same as $-\frac{1}{n} \sum_{t=1}^{n} \mathbf{H}\left(\mathbf{w}_{t} ; \hat{\boldsymbol{\theta}}\right)$, is consistent for $-\mathrm{E}\left[\mathbf{H}\left(\mathbf{w}_{t} ; \boldsymbol{\theta}_{0}\right)\right]$. Verify that the $\tilde{\boldsymbol{\Sigma}}$ given here, although not the same as $-\frac{1}{n} \sum_{t=1}^{n} \mathbf{H}\left(\mathbf{w}_{t} ; \tilde{\boldsymbol{\theta}}\right)$, is consistent for $-\mathrm{E}\left[\mathbf{H}\left(\mathbf{w}_{t} ; \boldsymbol{\theta}_{0}\right)\right]$. Hint: From the discussion in Example 7.8, $\hat{\boldsymbol{\theta}}$ is consistent. As mentioned in Section 7.4, $\tilde{\theta}$ too is consistent under the null. Given the consistency of $\hat{\theta}$ and $\tilde{\theta}$, it should be easy to show that $\hat{\sigma}^{2}$ and $\tilde{\sigma}^{2}$ are consistent for $\sigma^{2}$.\\
(d) Show that the Wald, LM, and LR statistics using $\widehat{\boldsymbol{\Sigma}}$ and $\widetilde{\boldsymbol{\Sigma}}$ given here in the formulas in Table 7.2 can be written as

$$
\begin{aligned}
W & =n \cdot \frac{(\mathbf{R} \widehat{\boldsymbol{\beta}}-\mathbf{c})^{\prime}\left[\mathbf{R}\left(\mathbf{X}^{\prime} \mathbf{X}\right)^{-1} \mathbf{R}^{\prime}\right]^{-1}(\mathbf{R} \widehat{\boldsymbol{\beta}}-\mathbf{c})}{S S R_{U}} \\
L M & =n \cdot \frac{(\mathbf{y}-\mathbf{X} \tilde{\boldsymbol{\beta}})^{\prime} \mathbf{P}(\mathbf{y}-\mathbf{X} \tilde{\boldsymbol{\beta}})}{S S R_{R}} \\
L R & =n \cdot\left\{\log \left(\frac{S S R_{R}}{n}\right)-\log \left(\frac{S S R_{U}}{n}\right)\right\}
\end{aligned}
$$

where $\mathbf{y}(n \times 1)$ and $\mathbf{X}(n \times K)$ are the data vector and matrix associated with $y_{t}$ and $\mathbf{x}_{t}$, and $\mathbf{P}=\mathbf{X}\left(\mathbf{X}^{\prime} \mathbf{X}\right)^{-1} \mathbf{X}^{\prime}$. Hint: The $\mathbf{A}(\boldsymbol{\theta})(r \times(K+1))$ in Table 7.2 is

$$
(\underset{(r \times K)}{\mathbf{R}} \vdots \underset{(r \times 1)}{\mathbf{0}}) .
$$

(e) Show that the three statistics can also be written as

$$
\begin{aligned}
W & =n \cdot \frac{S S R_{R}-S S R_{U}}{S S R_{U}}, \\
L M & =n \cdot \frac{S S R_{R}-S S R_{U}}{S S R_{R}}, \\
L R & =n \cdot \log \left(\frac{S S R_{R}}{S S R_{U}}\right)
\end{aligned}
$$

Hint: As we have shown in an analytical exercise of Chapter 1,

$$
\begin{aligned}
S S R_{R}-S S R_{U} & =(\widehat{\boldsymbol{\beta}}-\widetilde{\boldsymbol{\beta}})^{\prime}\left(\mathbf{X}^{\prime} \mathbf{X}\right)(\widehat{\boldsymbol{\beta}}-\widetilde{\boldsymbol{\beta}}) \\
& =(\mathbf{R} \widehat{\boldsymbol{\beta}}-\mathbf{c})^{\prime}\left[\mathbf{R}\left(\mathbf{X}^{\prime} \mathbf{X}\right)^{-1} \mathbf{R}^{\prime}\right]^{-1}(\mathbf{R} \widehat{\boldsymbol{\beta}}-\mathbf{c}) \\
& =(\mathbf{y}-\mathbf{X} \widetilde{\boldsymbol{\beta}})^{\prime} \mathbf{P}(\mathbf{y}-\mathbf{X} \widetilde{\boldsymbol{\beta}})
\end{aligned}
$$

(f) Show that $W \geq L R \geq L M$. (These inequalities do not always hold in nonlinear regression models.)

\section*{ANSWERS TO SELECTED QUESTIONS}
2a. Since $f(y \mid \mathbf{x} ; \boldsymbol{\theta})$ is a hypothetical density, its integral is unity:


\begin{equation*}
\int f(y \mid \mathbf{x} ; \boldsymbol{\theta}) \mathrm{d} y=1 \tag{1}
\end{equation*}


This is an identity, valid for any $\theta \in \Theta$. Differentiating both sides of this identity with respect to $\theta$, we obtain


\begin{equation*}
\frac{\partial}{\partial \boldsymbol{\theta}} \int f(y \mid \mathbf{x} ; \boldsymbol{\theta}) \mathrm{d} y=\underset{(p \times 1)}{\mathbf{0}} \tag{2}
\end{equation*}


If the order of differentiation and integration can be interchanged, then


\begin{equation*}
\frac{\partial}{\partial \boldsymbol{\theta}} \int f(y \mid \mathbf{x} ; \boldsymbol{\theta}) \mathrm{d} y=\int \frac{\partial}{\partial \boldsymbol{\theta}} f(y \mid \mathbf{x} ; \boldsymbol{\theta}) \mathrm{d} y \tag{3}
\end{equation*}


But by the definition of the score, $\mathbf{s}(\mathbf{w} ; \boldsymbol{\theta}) f(y \mid \mathbf{x} ; \boldsymbol{\theta})=\frac{\partial}{\partial \boldsymbol{\theta}} f(y \mid \mathbf{x} ; \boldsymbol{\theta})$. Substituting this into (3), we obtain


\begin{equation*}
\int \mathbf{s}(\mathbf{w} ; \boldsymbol{\theta}) f(y \mid \mathbf{x} ; \boldsymbol{\theta}) \mathrm{d} y=\underset{(p \times 1)}{\mathbf{0}} \tag{4}
\end{equation*}


This holds for any $\boldsymbol{\theta} \in \Theta$, in particular, for $\boldsymbol{\theta}_{0}$. Setting $\boldsymbol{\theta}=\boldsymbol{\theta}_{0}$, we obtain


\begin{equation*}
\int \mathbf{s}\left(\mathbf{w} ; \boldsymbol{\theta}_{0}\right) f\left(y \mid \mathbf{x} ; \boldsymbol{\theta}_{0}\right) \mathrm{d} y=\mathrm{E}\left[\mathbf{s}\left(\mathbf{w} ; \boldsymbol{\theta}_{0}\right) \mid \mathbf{x}\right]=\underset{(p \times 1)}{\mathbf{0}} \tag{5}
\end{equation*}


Then, by the Law of Total Expectations, we obtain the desired result.

\section*{References}
Amemiya, T., 1985, Advanced Econometrics, Cambridge: Harvard University Press.\\
Davidson, R., and J. MacKinnon, 1993, Estimation and Inference in Econometrics, New York: Oxford University Press.\\
Gallant, R., and H. White, 1988, A Unified Theory of Estimation and Inference for Nonlinear Dynamic Models, New York: Basil Blackwell.

Gourieroux, C., and A. Monfort, 1995, Statistics and Econometric Models, New York: Cambridge University Press.\\
Hansen, L. P., and K. Singleton, 1982, "Generalized Instrumental Variables Estimation of Nonlinear Rational Expectations Models," Econometrica, 50, 1269-1286.\\
Judge, G., W. Griffiths, R. Hill, H. Lütkepohl, and T, Lee, 1985, The Theory and Practice of Econometrics (2d ed.), New York: Wiley.\\
Newey, W., and D. McFadden, 1994, "Large Sample Estimation and Hypothesis Testing," Chapter 36 in R. Engle and D. McFadden (eds.), Handbook of Econometrics, Volume IV, New York: North-Holland.\\
White, H., 1994, Estimation, Inference, and Specification Analysis, New York: Cambridge University Press.


\end{document}
